% FIXME Ao final, deixe descomentada a linha correspondente ao numero de paginas
% que a sua defesa possui.
\documentclass[12pt, letterpaper, twoside]{book}

%%%%%%%%%%%%%%%%%%%%%%%%%%%% Configuração: pacotes %%%%%%%%%%%%%%%%%%%%%%%%%%%%
%%%%%%%%%%%%%%%%%%%%%%%%%%%%%% Pacotes: básicos %%%%%%%%%%%%%%%%%%%%%%%%%%%%%% 
\usepackage[utf8]{inputenc}
\usepackage{cmap}
\usepackage[T1]{fontenc}
\usepackage[english,brazilian]{babel}
\usepackage{indentfirst}
\usepackage[top=3cm,bottom=3cm,right=2cm,left=2cm]{geometry}


%\usepackage[hyperref=true, url=true, doi=true, natbib=true, backref=true, backrefstyle=three, backend=biber, refsegment=chapter, citestyle=numeric-comp, style=chem-rsc, maxcitenames=5, date=year, block=none, maxbibnames=100]{biblatex}
\usepackage[style=alphabetic,citestyle=authoryear,backend=biber]{biblatex}
\addbibresource{biblatex_style_samples/sample.bib}
\addbibresource{biblatex_style_samples/sample-abel.bib}


%%%%%%%%%%%%%%%%%%%%%%%%%%%%%%% Pacotes: layoyt %%%%%%%%%%%%%%%%%%%%%%%%%%%%%%%
\usepackage{etoolbox}  % É preciso para mudar o layout do frontmatter


%%%%%%%%%%%%%%%%%%%%%%%%%%%%%%% Pacotes: links %%%%%%%%%%%%%%%%%%%%%%%%%%%%%%%
\usepackage{url}
\usepackage{breakurl}
\usepackage{hyperref}
\hypersetup{
	colorlinks=true,
	linkcolor=blue,
	filecolor=magenta,      
	urlcolor=cyan,
	citecolor= red,
}
%\usepackage{cite}
%\renewcommand{\citeleft}{\textcolor{red}{[}}
%\renewcommand{\citeright}{\textcolor{red}{]}}

% FIXME Comente o pacote abaixo quando for concluir sua defesa e for entregar a
% versão final.
%\usepackage{showkeys}


%%%%%%%%%%%%%%%%%%%%%%%%%%%%%%%% Pacotes: ams %%%%%%%%%%%%%%%%%%%%%%%%%%%%%%%% 
\usepackage{amsmath}
\usepackage{amsfonts}
\usepackage{amssymb}
\usepackage{amsthm}
\usepackage{breqn}


%%%%%%%%%%%%%%%%%%%%%%%%%%%%%% Pacotes: tabelas %%%%%%%%%%%%%%%%%%%%%%%%%%%%%%
\usepackage{multicol}
\usepackage{multirow}
\usepackage{array}
\usepackage{booktabs}
\usepackage{longtable}


%%%%%%%%%%%%%%%%%%%%%%%%%%%%%% Pacotes: cores %%%%%%%%%%%%%%%%%%%%%%%%%%%%%%%% 
\usepackage[usenames,dvipsnames,svgnames,table]{xcolor}


%%%%%%%%%%%%%%%%%%%%%%%%%%%%%% Pacotes: figuras %%%%%%%%%%%%%%%%%%%%%%%%%%%%%% 
\usepackage{pdfpages}
\usepackage{graphicx}
\usepackage{tikz}
\usetikzlibrary{fit}
\usepackage{wrapfig}


%%%%%%%%%%%%%%%%%%%%%%%%%%%%% Pacotes: algoritmos %%%%%%%%%%%%%%%%%%%%%%%%%%%%% 
\usepackage{algorithmic}
\usepackage[chapter]{algorithm}
\floatname{algorithm}{Algoritmo}
\renewcommand{\listalgorithmname}{Lista de Algoritmos}


%%%%%%%%%%%%%%%%%%%%%%%%%%%%%% Pacotes: códigos %%%%%%%%%%%%%%%%%%%%%%%%%%%%%% 
\usepackage{textcomp}
\usepackage{listings}
\renewcommand\lstlistingname{Código}
\renewcommand\lstlistlistingname{Lista de Códigos}


%%%%%%%%%%%%%%%%%%%%%%%%%%%%%%% Pacotes: index %%%%%%%%%%%%%%%%%%%%%%%%%%%%%%% 
\usepackage{makeidx}
\makeindex


%%%%%%%%%%%%%%%%%%%%%%%%%%%%%%% Pacotes: fontes %%%%%%%%%%%%%%%%%%%%%%%%%%%%%% 
\usepackage{lmodern} \normalfont
\DeclareFontShape{T1}{lmr}{bx}{sc} { <-> ssub * cmr/bx/sc }{}
\usepackage{mathrsfs}


% TODO Inserir pacotes adicionais aqui.
\usepackage{color,soul}
\usepackage{bm}
\usepackage{subcaption}
\usepackage{relsize}
\usepackage{tikz}  % Arquivo com os pacotes.

%%%%%%%%%%%%%%%%%%%%%%%%% Configuração: dados pessoais %%%%%%%%%%%%%%%%%%%%%%%%%
% FIXME Substituir 'Nome completo do aluno' pelo seu nome.
\newcommand{\autor}{Pedro Leal Pazzini da Silva}
% FIXME Se for do sexo feminino, descomente a linha a seguir.
% \def\femaleAuthor{}

% FIXME Substituir 'Título da defesa' pelo título da defesa.
\newcommand{\titulo}{Agrupamento de séries temporais de consumo de carga}
% FIXME Se estiver no programa de mestrado, descomente a linha a seguir.
% \def\mestrado{}
% FIXME Deixe descomente apenas a linha referente ao departamento.
% \def\matematica{}
\def\aplicada{}
% \def\estatistica{}

% FIXME Substituir 'Nome completo do orientador' pelo nome completo do seu
% orientador.
\newcommand{\orientador}{Cristiano Leite de Castro}
% FIXME Se for orientado por uma mulher, descomente a linha a seguir.
% \def\femaleOrientador{}

% FIXME Substituir 'Nome completo do coorientador' pelo nome completo do seu
% coorientador. Caso não tenha coorientador, comente a linha a seguir.
\newcommand{\coorientador}{Frederico Gadelha Guimarães}
% FIXME Se for coorientado por uma mulher, descomente a linha a seguir.
% \def\femaleCoorientador{}

% FIXME Substituir 'Ano' pelo ano em que ocorreu sua defesa.
\newcommand{\ano}{2016}

% FIXME: ppazzini: Local da defesa
\newcommand{\local}{Belo Horizonte}
\newcommand{\universidade}{Universidade Federal de Minas Gerais }
\newcommand{\ppgee}{Programa de Pós-Graduação em Engenharia Elétrica }
\newcommand{\mestrado}{True}

%%%%%%%%%%%%%%%%%%%%%%%%%%% Configuração: definições %%%%%%%%%%%%%%%%%%%%%%%%%%%
\input{configuracao}  % Arquivo com algumas configurações.

%%%%%%%%%%%%%%%%%%%%%%%%%% Início do texto da defesa %%%%%%%%%%%%%%%%%%%%%%%%%%%
\begin{document}
% WARNING Todas as paginas deverão ser numeradas.
%
% As paginas iniciais são numeradas com algoritmos romanos em sua forma
% minuscula.
\frontmatter
%
\thispagestyle{plain}
\includegraphics[width=.94in, height=1in,
keepaspectratio=true]{figuras/aplicacoes_ufmg/negativo_completa_ufmg.jpg}
\vspace*{1cm}
\begin{center}
  % O tamanho da fonte deve ser 16pt.
  % Deve-se utilizar caixa alta.
  {\Large\textsc{\autor}}
\end{center}
\vspace{4cm}
\begin{center}
  % O tamanho da fonte deve ser 16pt em negrito.
  % Deve-se utilizar caixa alta.
  {\Large\textbf{\textsc{\titulo}}}
\end{center}
\vfill
\begin{center}
  % O tamanho da fonte deve ser 12pt em negrito.
  % Deve-se utilizar caixa alta.
  \textbf{\local \\ \ano}
\end{center}
  % Não edite esse arquivo.
\newpage\mbox{}\thispagestyle{plain}\newpage  % Pagina em branco.
%
% WARNING A folha de rosto precisa ser assinada pelos orientadores.
% WARNING Você deve escanear como preto-e-branco.
% FIXME Substitua arquivo folha-de-rosto.pdf por uma copia escaneada, comente
% esta linha e descomente a próxima.
\thispagestyle{plain}
% WARNING Não modifique este arquivo.
\includegraphics[width=.94in, height=1in,
keepaspectratio=true]{figuras/aplicacoes_ufmg/negativo_completa_ufmg.jpg}
\begin{center}
  {\large\textbf{\textsc{\universidade}}
  \vspace{.5cm}

  \ppgee}
\end{center}
\vfill
\begin{center}
  {\large\textbf{\textsc{\autor}}}
\end{center}
\vfill
\begin{center}
  {\Large\textbf{\textsc{\titulo}}}
\end{center}
\vfill

\begin{flushright}
  \begin{minipage}[c]{.5\textwidth}
    \ifx\mestrado\undefined
    Tese
    \else
    Dissertação
    \fi
    apresentada ao \ppgee da \universidade como parte dos requisitos exigidos
    para a obtenção do título de
    \ifx\mestrado\undefined
    \ifx\femaleAuthor\undefined
    Doutor
    \else
    Doutora
    \fi
    \else
    \ifx\femaleAuthor\undefined
    Mestre
    \else
    Mestra
    \fi
    \fi
    em Engenharia Elétrica.
  \end{minipage}
\end{flushright}
\vspace{.5cm}

\noindent
\textbf{Orientador\ifx\femaleOrientador\undefined
\else
a\fi: \orientador
}
\vspace{.25cm}

\ifx\coorientador\undefined
\else
\noindent
\textbf{Coorientador\ifx\femaleCoorientador\undefined
\else
a\fi: \coorientador
}
\vspace{.5cm}
\fi

\noindent
\begin{minipage}[c]{.5\textwidth}
  {\footnotesize\textsc{Este exemplar corresponde à versão final da
  \ifx\mestrado\undefined
  tese
  \else
  dissertação
  \fi
  defendida
  \ifx\femaleAuthor\undefined
  pelo aluno
  \else
  pela aluna
  \fi
  \autor,
  e orientada pel\ifx\femaleOrientador\undefined
  o\else
  a\fi{} Prof\ifx\femaleOrientador\undefined
  \else
  a\fi. Dr\ifx\femaleOrientador\undefined
  \else
  a\fi. \orientador.
  }}
\end{minipage}
\vspace{1cm}

\noindent
{\small\textbf{Assinatura
\ifx\femaleOrientador\undefined
do Orientador
\else
da Orientadora
\fi
}

\vspace{.5cm}
\noindent
\rule[1pt]{7cm}{.5pt}  % Linha para assinatura do orientador
}
\vspace{.5cm}

\ifx\coorientador\undefined
\else
\noindent
{\small\textbf{Assinatura
\ifx\femaleCoorientador\undefined
do Coorientador
\else
da Coorientadora
\fi
}

\vspace{.5cm}
\noindent
\rule[1pt]{7cm}{.5pt}  % Linha para assinatura do coorientador
}
\fi
\vfill
\begin{center}
  {\small\textbf{\textsc{ \local \\ \ano}}}
\end{center}

% \includepdf{folha-de-rosto.pdf}
%
% WARNING A ficha catalográfica deve estar no verso da folha de rosto.
% WARNING Você deve escanear como preto-e-branco.
% FIXME O arquivo ficha-catalografica.pdf deve ser sobrescrito com uma cópia
% do arquivo pdf que a biblioteca lhe enviar.
%\includepdf{ficha-catalografica}
%
% WARNING A folha de aprovação deve ser assinada pelos membros da banca apos a
% defesa.
% WARNING Você deve escanear como preto-e-branco.
% FIXME Substitua o arquivo folha-de-aprovacao.pdf por uma copia escaneada.
%\includepdf{folha-de-aprovacao}
\chapter*{}  % Capítulo vazio.
\begin{center}
  \large{\textbf{Abstract}}
\end{center}

\selectlanguage{english}

The union of computational, measurement and communication systems applied to electrical power systems will form the Smart Grids. One of its main components are the Smart Meters, which are measurement systems that periodically measure the consumed power in electrical loads, and that send the measurements to the electricity companies. The research, and mainly the installation of smart meters is more developed in the north hemisphere countries, although it already started in Brazil. When this measures become ubiquitous in the Brazilian electrical power system, a considerable amount of data will be produced, and the extraction of useful information from it will be necessary. One of the possible data mining tasks that could be applied to this data, and which is the object of this dissertation, is the clustering of load curves, from which is expected to get information and knowlodge that can enhance the distribution service, the load forecast , the planning of expansion, and others. As the data of the smart meters installed in the Brazilian electrical system is not publicly available, the load curves of Australian residences were clustered. Since that, in this country, the climatic variations that could influence in the consuming profile are very similar to the climatical brazilian variations. The problem of clustering load profiles is a time series clustering problem, since that for each measurement there is the temporal component associated. Thus, in this master dissertation was presented, initially, a study of clustering of time series, and this study is focused in the preprocessing, dissimilarity measure and algorithm applied, as well as the method of evaluating the results. Initially, after an empirical study in labeled datasets, the general clustering of time series strategy was proposed, and it is composed by: Z normalization in the preprocessing step, CID-DTW as the dissimilarity the hierarchical clustering algorithm using the average linkage and the silhouette And Calinski-Harabasz index for the evaluation of the obtained partitions. In tha case study, with the australian dataset, such strategy generated highly unbalanced partitions, which motivated the proposition of a multiobjective approach, in which partitions with a good compromise between the validation index values, wich is a measure of the desired low variation intragroup and high variance intergroup, and balanced partitions, in which the variance of the number of instances in each cluster is low. The multi-objective approach indicates the tha max normalization, the k-means algorithm and the euclidean distance is the most appropiate configuration for the clustering of load profiles generated by the smart meters.

\vspace{.5cm}
\textbf{Keywords}:
% FIXME Remover a linha abaixo.
Clustering, time series, load curves, smart meters.
% TODO Inserir as palavras-chave em inglês aqui.
\selectlanguage{brazilian}

\vspace{.1\textheight}
\begin{center}
  \large{\textbf{Resumo}}
\end{center}

% TODO Inserir o resumo em português aqui.
A união de sistemas computacionais, de medição e de comunicação aplicada aos sistemas elétricos de potência dará forma às chamadas Redes Inteligentes. Um de seus componentes principais são os chamados medidores inteligentes, que são instrumentos que medem periodicamente a potência consumida em cargas elétricas, e enviam estes dados de medição para as distribuidoras de energia. A pesquisa, e principalmente a instalação de medidores inteligentes é mais intensa nos países do hemisfério norte, no entanto ela já se iniciou também no Brasil. A partir do momento em que esses medidores se tornarem onipresentes no sistema elétrico de potência brasileiro, uma massa de dados considerável será gerada, e se fará necessária a extração de informação útil dela. Uma das possíveis tarefas de mineração de dados a ser realizada nesta massa de dados, e que será objeto desta dissertação de mestrado, será o agrupamento das cargas, a partir do qual se espera obter informações e conhecimento que possa aprimorar o serviço de distribuição, a previsão de demanda, o planejamento de expansão da rede, dentre outras possibilidades. Uma vez que ainda não estão disponíveis os dados dos medidores inteligentes instalados no Brasil, foi realizado o agrupamento de curvas de carga medidas em residências australianas, já que neste país, as variações climáticas que podem influenciar no consumo de energia elétrica são bem similares às variações climáticas brasileiras. O problema de agrupamento de curvas de carga, é um problema de agrupamento de séries temporais, já que a cada medição realizada pelos medidores inteligentes existe uma componente temporal associada. Assim, nesta dissertação de mestrado foi apresentado, inicialmente, um estudo das estratégias de agrupamento de séries temporais no que diz respeito ao pré-processamento empregado, métrica de dissimilaridade e algoritmo de agrupamento utilizados, bem como o método de avaliação dos resultados obtidos. Inicialmente, em um estudo empírico em base de dados rotuladas, foi definida uma estratégia genérica para o agrupamento de séries temporais, composta por: normalização Z na etapa de pré-processamento, escolha da métrica de dissimilaridade CID-DTW e o algoritmo hierárquico com \emph{linkagem average} para a realização do agrupamento em si, e os índices de validação silhouette e Calinski-Harabasz para a avaliação das partições obtidas. No estudo de caso, com os dados australianos, tal estratégia gerou grupos altamente desbalanceados, o que motivou a proposição de uma abordagem multi-objetiva, na qual se procura obter partições com um bom índice de validação, que por sua vez indica a tão desejada baixa variância intragrupo e alta variância intergrupo, e partições balanceadas, nas quais o número de instâncias em cada grupo possui baixa variância. A abordagem multi-objetiva indica que a normalização \emph{max}, aliada ao algoritmo \emph{k-means}, fazendo uso da distância euclidiana como métrica de dissimilaridade é a configuração mais adequada para o agrupamento de curvas de carga geradas por medidores inteligentes.

\vspace{.2cm}
\textbf{Palavras-chave}:

% TODO Inserir as palavras-chave aqui.
Agrupamento, séries temporais, curvas de carga, medidores inteligentes.
%
\tableofcontents
%
% FIXME Se não for incluir a dedicatória, comentar a linha abaixo.
%\input{dedicatoria}
%
% FIXME Se não for incluir os agradecimentos, comentar a linha abaixo.
\input{agradecimentos}
%
% FIXME Comentar a linha abaixo se não desejar listar as figuras
% apresentadas.
\renewcommand{\listfigurename}{Lista de Ilustrações}  % CCPG 228/2013
\listoffigures
%
% FIXME Comentar a linha abaixo se não desejar listar as tabelas
% apresentadas.
\listoftables
%
% FIXME Comentar a linha abaixo de não for apresentar as
% abreviações e siglas utilizadas.
\chapter*{Lista de Abreviaturas e Siglas\markboth{Lista de Abreviaturas e Siglas}{}}  % \markboth{}{} é utilizado para corrigir o cabeçalho.

\begin{description}
  % FIXME Remover as duas abreviações/siglas abaixo e incluir as que serão
  % utilizadas.
 % \item[FIXME] Indica que algo deve ser consertado.
  %\item[TODO] Indica que algo deve ser feito.
  \item [REI] Redes Elétricas Inteligentes
  \item [MI] Medidores Inteligentes
  \item[SVD] \emph{Singular Value Decomposition} (Decomposição em valores singulares)
  \item[DFT] \emph{Discrete Fourier Transform} (Transformada discreta de Fourier)
  \item[DBI] \emph{Davies-Bouldin Index}
  \item[DTW] \emph{Dynamic Time Warping}
  \item[EDR] \emph{Edit distance on real sequence}
  \item[RAM] \emph{Random Access Memory}, ou memória de acesso aleatória, que é a memória principal da maioria dos sistemas computacionais.
  \item[DBI] \emph{Davies-Bouldin index}, ou índice Davies-Bouldin, que é um índice de validação interna de grupos.
  \item[CID] \emph{Complex Invariant Distance}
  \item[VI] \emph{Variation of Information}, ou índice de variação da informação, que é um índice de validação externa de grupos.
  \item[ANEEL] Agência Nacional de Energia Elétrica

\end{description}

%
% FIXME Comentar a linha abaixo se não for apresentar os
% símbolos utilizados.
\chapter*{Lista de Símbolos\markboth{Lista de Símbolos}{}}  % \markboth{}{} é utilizado para corrigir o cabeçalho.

\begin{description}
    % FIXME Remover os dois símbolos abaixo e incluir as que serão
    % utilizadas.
    \item[$L_1$] Norma $L_1$, distância \emph{Manhattan}, ou distância \emph{cityblock}.
    \item[$L_2$] Norma $L_2$ ou distância Euclidiana.
    \item[$L_\infty$] Norma $L_\infty$ ou distância Chebyshev.
    \item[$dist(\bm{a},\bm{b})$] Métrica de dissimilaridade entre as séries temporais $\bm{a}$ e $\bm{b}$.
\end{description}

%
% FIXME Comentar a linha abaixo se não desejar listar os algoritmos
% apresentados.
\listofalgorithms
%
% FIXME Comentar a linha abaixo se não desejar listar os trechos de código
% apresentados.
%\lstlistoflistings
%
% FIXME Ao final, se seu trabalho tiver menos de 100 páginas descomente as
% três linhas a seguir.
% \makeatletter
% \@openrightfalse
% \makeatother
%
% As paginas com o conteúdo da tese são numeradas com algoritmos arábicos.
\mainmatter
%
% FIXME Remover as 3 linhas abaixo.
%\input{ex_cap1}
%\input{ex_cap2}
%\input{ex_cap3}
%\input{ex_cap4}
%
% TODO Inserir os arquivos referentes ao corpo da tese.
%
\chapter{Introdução}
Ao conjunto de tecnologias que visam modernizar o sistema elétrico de potência, dá-se o nome de Redes Elétricas Inteligentes (\emph{Smart Grids}), doravante REI. Tal conjunto de tecnologias faz uso de sistemas de medição, computação e comunicações digitais de forma a gerar um fluxo de informações bidirecional entre o fornecedor do sistema de energia e o cliente final. Os maiores impactos da adoção de tais tecnologias ocorrerão na distribuição de energia elétrica e nos pequenos consumidores. Pelo número crescente de investimentos em pesquisa de REI em todo mundo, percebe-se que esta se apresenta como uma forte tendência de modernização dos sistemas elétricos de potência em diversos países ~\parencite{REI}.

Uma das bases dessa nova tecnologia é a instalação de sensores, conhecidos como \emph{smart meters} ou medidores inteligentes (MI),  em diversos pontos da rede, de forma a se obter dados relevantes da mesma. Com a instalação desses sensores em cada consumidor, seja ele residencial, comercial ou industrial, será possível medir periodicamente a potência consumida por ele.  De posse dos dados gerados pelos medidores inteligentes, os consumidores podem ser classificados por padrões de consumo e a partir dessa classificação espera-se obter, por exemplo, as seguintes vantagens:
\begin{itemize}
	\item Realização de uma nova política de tarifas diferenciadas por tipo de consumidor residencial.
	\item Aprimoramento das técnicas de previsão de demanda por parte dos fornecedores
	de energia.
	\item Auxílio dos consumidores na análise do seu consumo de energia para a elaboração
	de estratégias de redução deste.
	\item Fomento de um mercado competitivo de fornecimento de energia.
	\item Identificação automática de falhas e fraudes.
\end{itemize}

De forma geral, a análise inteligente de dados extraídos dos \emph{smart meters} permitirá
às companhias o desenvolvimento de novas capacidades em termos de gerenciamento e
planejamento da rede, bem como novos serviços direcionados ao consumidor visando a
eficiência energética ~\parencite{ReviewElectric}.

O grande volume de dados gerados pelos \emph{smart meters} demandarão algoritmos e estratégias capazes de extraírem informação útil dessa grande massa de dados. A tarefa de se encontrar subgrupos em determinado conjunto de dados baseando-se somente nas características dos elementos do conjunto é denominada agrupamento, que, do ponto de vista da Inteligência Computacional, faz parte do paradigma de aprendizado não-supervisionado. Ou seja, o agrupamento se caracteriza pela necessidade de se formar grupos dos dados, de forma que os dados são mais semelhantes entre si quando contidos no mesmo grupo, ao passo que estes são mais díspares possível dos dados pertencentes a outros grupos. 

Por não-supervisionado, entende-se que a tarefa de agrupamento se dará apenas a partir das características dos dados, e que não existe rotulação dos mesmos, ou seja, não se sabe a priori à qual classe, ou grupo, os dados pertencem e, em certos problemas, não se sabe também o número de grupos nos quais os dados se encontram naturalmente divididos.

O problema de se dividir as curvas de carga geradas por MI em grupos é um problema de agrupamento de séries temporais, pois na maioria das aplicações, os dados de consumo são medidos em intervalos regulares, caracterizando, assim, uma série temporal. Ademais, não sabemos de antemão a qual grupo cada carga pertence e nem o número de grupos que contêm os dados, logo devido à ausência de dados rotulados, constitui-se como uma tarefa de agrupamento.

\section{Objetivos e contribuições}

O objetivo geral desta dissertação de mestrado é de investigar as técnicas consideradas estado da arte para o agrupamento de séries temporais, visando à proposta de uma metodologia para agrupamento de curvas de carga de consumidores residenciais. Tal investigação envolveu o estudo teórico-empírico de diferentes formas de normalização de dados, métodos de agrupamento, medidas de dissimilaridade e índices de validação.

A Agência Nacional de Energia Elétrica, ANEEL, em uma resolução de 2010 ~\parencite{ANEEL}, definiu os critérios de divisão tarifária dos seus consumidores. A primeira divisão que é feita, é baseada na tensão de atendimentos dos consumidores que são divididos em:
\begin{itemize}
	\item grupo A: tensão de atendimento superior à 2,3 kV, que pode ainda ser subdivido em:
	\begin{itemize}
		\item subgrupo A1: atendimento em tensão igual ou superior à 230 kV,
		\item subgrupo A2: atendimento em tensão de 88 kV a 130 kV,
		\item subgrupo A3: atendimento em tensão de 69 kV,
		\item subgrupo A4: atendimento em tensão de 2,3 kV a 44 kV,
		\item subgrupo AS: tendimento em tensão inferior a 2,3 kV mas por sistema subterrâneo,
	\end{itemize}
	\item e grupo B: tensão de atendimento igual ou inferior à 2,3kV, que também é subdividido em:
	\begin{itemize}
		\item subgrupo B1: atendimento residencial,
		\item subgrupo B2: atendimento rural,
		\item subgrupo B3: atendimento das demais classes,
		\item subgrupo B4: atendimento da iluminação pública.
	\end{itemize}
\end{itemize}

Existem ainda, mais subdivisões, para alguns dos subgrupos do grupo B, e especificamente o subgrupo B1, pode ser dividido, ainda, em:
\begin{itemize}
	\item atendimento residencial,
	\item e atendimento residencial de Baixa Renda -TSEE (Tarifa Social de Energia Elétrica).
\end{itemize}
Esta divisão implica em um cálculo tarifário diferenciado para as famílias de baixa renda, e seus critérios baseiam-se não nos perfis de curva de carga, mas sim em fatores de renda e de consumo, como pode ser visto em ~\parencite{TSEE}.

O cálculo tarifário das cargas do grupo A, que são aquelas que demandam uma maior potência, sendo em geral, cargas industriais e comerciais, possuem um cálculo tarifário bem mais complexo do que o cálculo tarifário residencial. As empresas podem contratar uma potência mínima das concessionárias de energia, pagando multas, inclusive, se o seu consumo ultrapassar o valor contratado. Além deste fator, para as cargas do grupo A, o cálculo tarifário varia de acordo com o mês e a hora do dia de consumo, o que motiva as empresas a fazerem um uso mais racional da energia.

Dessa maneira, o objetivo deste trabalho é a realização da proposição de uma técnica que gere uma segmentação maior dos consumidores residencias, baseada em perfis de consumo e não em critérios de renda, de forma a ser possível a proposição de uma política tarifária diferenciada para os consumidores residencias, assim como já é feito atualmente para os consumidores do grupo A.

Seria extremamente interessante que os dados a serem utilizados durante o trabalho fossem o mais representativos possível à realidade brasileira, ou seja, sem as influências da grande variação de temperatura que ocorre nos países do hemisfério norte, que é onde se encontra mais avançada a pesquisa e implantação de REI e, consequentemente, de instalação de MI e a disponibilidade de dados para pesquisa. Em ~\parencite{REI} é feita uma menção à um projeto da CEMIG em Sete Lagoas denominado "Cidade do Futuro". Tal projeto é reconhecido como pioneiro no Brasil no que se refere à pesquisa e instalação de equipamentos de REI no atual sistema elétrico de potência brasileiro. No entanto, estes dados ainda não são de domínio público. Dessa maneira, neste trabalho, optou-se por utilizar dados de cargas residenciais da Austrália, obtidos de ~\parencite{dadosAustralia}.

Os dados consistem de cerca de três anos e meio de medições do consumo de aproximadamente 13000 residências australianas, com uma taxa de amostragem de trinta minutos. Por questões de privacidade, somente as curvas de demanda foram disponibilizadas, e outros dados que poderiam ser úteis na tarefa de agrupamento como: localização das residências, número e situação econômica dos moradores, entre outros, não foram fornecidos. Dessa maneira, a única informação disponível para a realização do agrupamento das cargas são as curvas de demanda. Existem dados inconsistentes e incompletos o que demandou uma fase de pré-processamento dos mesmos, no qual o número de curvas a serem agrupadas foi reduzido para 427.

Para a realização do agrupamento de curvas de carga fazem-se necessárias as escolhas de um tipo de normalização, de uma métrica de dissimilaridade para se comparar as cargas, de um algoritmo de agrupamento para se obter a partição dos dados, e de um índice de validação que indique a qualidade da partição obtida. 

A etapa de normalização tem como objetivos principais reduzir a dimensionalidade das séries temporais e  facilitar a comparação entre séries temporais que possuam amplitudes ou \emph{offsets} muito discrepantes. As abordagens mais comuns para a redução de dimensionalidade consistem na \emph{Piecewise aggregate approximation} ~\parencite{SAX} ou na transformada discreta de Fourier ~\parencite{FFT}. Já a eliminação de distorções, tais como amplitude e \emph{offset}, pode ser realizada por meio da normalização Z ~\parencite{trillions}, ou por meio da normalização \emph{max} como utilizada em ~\parencite{Chicco}. Em geral, em problemas de agrupamento de cargas é utilizada a normalização \emph{max} sem nenhum tipo de redução de dimensionalidade, como pode ser visto em ~\parencite{Chicco, Flath2012,Tsekouras20081494}, no entanto, outras abordagens foram testadas.

Uma vez definido a normalização a ser aplicada nos dados, surge a necessidade de se escolher a medida de dissimilaridade a ser utilizada na comparação entre as cargas. A distância euclidiana se apresenta onipresente em estudos de agrupamento de cargas, no entanto, em outras aplicações como, por exemplo, na classificação de séries temporais, medidas de dissimilaridades como a DTW \parencite{DTW}, corrigidas por métodos de invariância à complexidade ~\parencite{CID}, têm obtido desempenho superior à distância euclidiana ~\parencite{BatistaComparativo}. Isso motivou o presente estudo a verificar a aplicabilidade de tais métricas ao problema de agrupamento de séries temporais em geral e de agrupamento de curvas de carga em específico.

Após o estabelecimento de uma métrica de dissimilaridade para a comparação das cargas, faz-se necessária a escolha de um algoritmo de agrupamento. O algoritmo mais utilizado em tarefas de agrupamento de cargas é o \emph{k-means}, devido à sua simplicidade e praticidade. No entanto, ele possui a deficiência de fazer uso de somente a distância euclidiana ~\parencite[][254]{Ullman}, o que descarta potenciais ganhos a serem obtidos com a escolha de outras medidas de dissimilaridades. O \emph{k-means} assume que o formato dos grupos é do tipo hiper-esfera (esférico) , o que pode não ser verdade para dados multidimensionais como dados de séries temporais. Ademais, existem estudos que mostram que outros algoritmos de agrupamento, como os algoritmos hierárquicos, possuem desempenho superior aos algoritmos particionais, como o \emph{k-means}, no contexto de séries temporais ~\parencite{k_shape}.

A partição dos dados a partir da aplicação do algoritmo escolhido, deve ser avaliada por um índice que indique a qualidade da mesma. Tal índice é útil, também, para comparar os resultados obtidos por diferentes escolhas de normalização, métrica de dissimilaridade e algoritmo de agrupamento. Não existe um índice que seja sistematicamente empregado no agrupamento de curvas de carga, e dessa maneira, os principais índices de validação de agrupamento são discutidos e comparados neste trabalho.

Sendo assim, devido a multiplicidade de caminhos possíveis para a realização da tarefa de agrupamento de séries temporais, a primeira fase do estudo focou em analisar, por meio de conjuntos de dados \emph{benchmark} (rotulados) de séries temporais, uma gama de metodologias de agrupamento propostas na literatura com o intuito de se obter um direcionamento para a aplicação nos dados reais australianos.

Além da tarefa de agrupamento em si, os tempos e custos, no que diz respeito à utilização de processamento, também são objetos de estudo desta pesquisa. Assim, espera-se que este trabalho auxilie no dimensionamento de um sistema computacional a ser utilizado pelas companhias de distribuição de energia elétrica para o agrupamento de curvas de caragas obtidas por MI.

%De uma forma geral, espera-se com este trabalho que, a partir de dados estrangeiros, seja possível a criação de uma metodologia de extração de informações úteis a partir dos dados obtidos de MI, quando estes forem instalados em larga escala no sistema elétrico de potência brasileiro.

Dentre as principais contribuições desta dissertação, pode-se citar:
\begin{itemize}
	\item Extensa revisão bibliográfica acerca do problema e das várias opções para a realização de agrupamento de séries temporais.
	\item Proposição de uma metodologia genérica para a realização de agrupamento de séries temporais embasada em dados empíricos de experimentos realizados com bases de dados rotuladas.
	\item Proposição de uma metodologia específica para a realização de agrupamento de curvas de cargas residenciais.
\end{itemize}

\section{Visão Geral do Texto}

O restante do texto desta dissertação de mestrado se apresenta da seguinte maneira:

O Capítulo ~\ref{cap:agrupamento_series_temporais} apresenta uma extensa revisão bibliográfica acerca dos aspectos relevantes durante a tarefa de agrupamento de séries temporais. Nele, são destacadas as principais métricas de dissimilaridade, bem como os principais algoritmos de agrupamento para a realização do agrupamento. Ao final deste capítulo, também são discutidos os métodos de avaliação dos resultados do agrupamento.

Em seguida, no Capítulo ~\ref{cap:testes_teoricos}, são realizados testes em conjuntos de dados rotulados (\emph{benchmark}), com a finalidade de se avaliar os algoritmos de agrupamento, métricas de dissimilaridades e métodos de avaliação discutidos no Capítulo ~\ref{cap:agrupamento_series_temporais}. Neste terceiro capítulo, foram realizados três experimentos com a finalidade de se definir a estratégia de agrupamento a ser empregada no estudo de caso com dados reais. O primeiro destes experimentos teve como objetivo definir o algoritmo de agrupamento a ser empregado no referido estudo de caso. Já o segundo procurou definir a métrica de dissimilaridade, enquanto o terceiro visou apontar a melhor maneira de se avaliar o resultado do agrupamento. Em todos os experimentos foram realizados testes estatísticos de significância para se definir a melhor estratégia a ser empregada no agrupamento da base de dados australiana.

Após as definições obtidas a partir dos testes realizados no Capítulo ~\ref{cap:testes_teoricos}, o estudo de caso em si, que consiste no agrupamento das cuvas de carga das residências australianas, é descrito no Capítulo ~\ref{cap:estudo_de_caso}. Finalmente, no Capítulo ~\ref{cap:conclusao}, são apresentadas as conclusões gerais do trabalho bem como algumas das propostas para sua continuidade.

\chapter{Agrupamento de séries temporais} \label{cap:agrupamento_series_temporais}

Por agrupamento, entende-se a tarefa de dividir um conjunto de dados em grupos homogêneos e bem definidos. A partir de determinada métrica de dissimilaridade, o agrupamento busca formar grupos que contêm dados que possuem baixa dissimilaridade entre si, ao passo que estes possuem alta dissimilaridade com os dados contidos nos demais grupos. O agrupamento se faz necessário quando não se possui, \emph{a priori}, a rotulação dos dados, ou seja, não se sabe a qual grupo cada instância pertence, e em alguns casos, não se sabe até mesmo o número de grupos que as instâncias estão naturalmente contidas. A partir da realização de tal tarefa, espera-se adquirir conhecimento acerca da estrutura dos dados bem como extrair informação útil dos mesmos.

O agrupamento pode ser feito em dados textuais, categóricos, binários, numéricos, ordinais, relacionais, temporais, espaço-temporais, ou qualquer combinação destes. Quando o agrupamento é realizado em dados nos quais a observação das grandezas de interesse não variam com o tempo, os dados são considerados estáticos, e caso contrário os dados são ditos dinâmicos. As séries temporais, que são objeto de estudo deste trabalho, são dados dinâmicos. Grande parte da pesquisa de agrupamento se dá na análise de dados estáticos, no entanto, existe uma tendência de crescimento no estudo em séries temporais devido às diversas aplicações nas engenharias, mercado financeiro, marketing, medicina, biologia, meteorologia, astronomia entre outros ~\parencite{Liao}.

Formalmente, uma partição $\mathcal{C}$ é a partição de um conjunto de dados $X$ composto por $n$ instâncias, e cada instância possuindo $d$ características, em $k$ sub-conjuntos mutualmente disjuntos, $\mathcal{C}_1,\mathcal{C}_2,...,\mathcal{C}_k$. A esses sub-conjuntos dá-se o nome de grupos. Dessa maneira,

\begin{equation}
\mathcal{C} = {\mathcal{C}_1,\mathcal{C}_2,...,\mathcal{C}_k} \text{,  tal que } \mathcal{C}_k \cap \mathcal{C}_l = \emptyset
\end{equation}

Uma série temporal é definida como um conjunto de observações sequenciais de uma grandeza de interesse, assim, ela é um tipo de dados com valores reais e de alta dimensionalidade. Denomina-se de tamanho da série temporal o número de observações ou amostras que ela contém.

As seções a seguir deste capítulo descrevem as principais estratégias de agrupamento de séries temporais. Por estratégia entende-se a escolha de:
\begin{itemize}
	\item um pré-processamento dos dados (normalização),
	\item uma métrica de dissimilaridade para se comparar as instâncias,
	\item um algoritmo de agrupamento para a realização do agrupamento em si,
	\item e, finalmente, um índice de validação para avaliar a qualidade da partição resultante.
\end{itemize}
Este capítulo se encontra estruturado de forma a realizar um estudo de cada um destes aspectos. Inicialmente, na primeira seção, são analisadas as opções de pré-processamento dos dados, que consistem, basicamente em normalizar, e reduzir o tamanho das séries temporais.

 Em seguida, são analisadas as métricas de dissimilaridade de séries temporais que podem ser divididas em: métricas baseadas em distâncias ou dissimilaridades, aplicadas no espectro das séries temporais e em modelos estatísticos ~\parencite{Liao}. Pelo fato da escolha de uma métrica de dissimilaridade ser um dos aspectos centrais na tarefa de agrupamento  de séries temporais as principais métricas apresentadas na literatura são discutidas na maior parte deste capítulo.
 
 A escolha da métrica de dissimilaridade se dá pela definição de um critério, ou métrica , que nos diga o quão similares duas instâncias são. Tal métrica é dependente do objetivo principal da tarefa de agrupamento. Existem determinadas aplicações que realizam o agrupamento baseado no valor da predição da série temporal, ou seja, as curvas são agrupadas em função do próximo valor esperado de cada série, enquanto em outras aplicações, quando se deseja comparar o formato das séries, o uso de métricas ponto a ponto se adequa melhor. Existam ainda outras abordagens que buscam agrupar curvas que possuam o mesmo comportamento ou tendência, ou seja, curvas que se mantêm constante, aumentam ou diminuem o seu valor em períodos de tempo próximos.

Após a discussão das métricas, serão expostos os principais algoritmos de agrupamento de séries temporais disponíveis na literatura. Estes algoritmos podem ser estocásticos como o \emph{k-means}, \emph{k-medoids}, ou de agrupamento hierárquico com diferentes tipos de \emph{linkagem}, tais como: a \emph{average}, \emph{complete} ou \emph{ward}. Existem também os algoritmos de agrupamento baseados em densidade, como o \emph{DBSCAN}, que possuem a vantagem de, não ser necessária a definição, \emph{a priori}, do número de grupos.

Finalmente, na última seção do capítulo, são apresentados os principais índices de validação de agrupamentos disponíveis na literatura. Tais índices são essenciais para se mensurar a qualidade das partições dos dados resultantes dos algoritmos de agrupamento.

Tal revisão bibliográfica é importante para a definição da melhor estratégia, ou seja, da escolha de uma normalização, métrica de dissimilaridade, algoritmo de agrupamento e índice de validação das partições resultantes a ser empregada no agrupamento das curvas de carga.

\section{Pré-processamento de séries temporais}

Antes da tarefa de agrupamento em si, recomenda-se a realização de uma fase de pré-processamento dos dados, nos quais eles são normalizados. A normalização é utilizada para corrigir possíveis distorções que o dados brutos podem apresentar devido ao seu processo de aquisição. Pode-se corrigir distorções relacionadas a escala, translação, rotação, diferença de fase, entre outras.

Além da normalização, métricas de dissimilaridade podem ser aplicadas no espectro das séries temporais, como uma estratégia de redução de dimensionalidade, para a posterior realização da tarefa de agrupamento. Tal redução é desejada, pela ocorrência do "mal da dimensionalidade", que basicamente nos diz que quando o número de dimensões de cada instância de um conjunto de dados é muito grande, as instâncias tendem a se encontrar espaçadas e equidistantes, fazendo com que seja necessário um número muito grande de instâncias para a construção de uma fronteira de decisão ~\parencite{DudaHart}. Dessa maneira, faz-se necessária um etapa de pré-processamento, e as principais abordagens se encontram detalhadas a seguir.

\subsection{Normalização} \label{sec:norm_Z}

Estudos demonstram que qualquer tarefa de mineração de dados de séries temporais sem algum tipo de normalização é sem sentido ~\parencite{Keogh_need_norm}. O tipo de normalização mais comum no contexto de séries temporais é realizado por meio da normalização Z ~\parencite{trillions}, na qual, e média $\mu$ e o desvio padrão $\sigma$ de cada série temporal de tamanho $n$ é calculada e é feita a seguinte transformação em cada medição da série $\bf{x}$

\begin{equation}
\hat{x_i} = \frac{x_i - \mu}{\sigma}, i = 1,...,n.
\end{equation}

No contexto de séries temporais, a normalização Z é extremamente importante para remover certas imperfeições e diferenças de escala que surgem a partir da geração das séries temporais, sejam estas obtidas por meio de medições ou de algum tipo de transformação de dados. Na figura ~\ref{fig:z_norm} podem ser vistas duas senoides, $s1$ e $s2$, ligeiramente defasadas mas com diferença significativa em amplitude e valor médio, ou \emph{offset}. Na mesma figura, as suas respectivas versões normalizadas podem ser vistas pelas linhas tracejadas. É notável que as curvas tracejadas são bem mais semelhantes entre si, e caso as diferenças de amplitude e valor médio não sejam relevantes para a tarefa de agrupamento ou classificação, o que é geralmente o caso, então a normalização Z é uma etapa de pré-processamento obrigatória.
\begin{figure}[h!]
	\includegraphics[width=\linewidth]{figuras/z_norm.png}
	\caption{Exemplo de normalização Z. As duas senóides s1 e s2 que possuem comportamento semelhante mas, devido à diferença de \emph{offsets} e  amplitudes, se apresentam significativamente distintas. Após a aplicação da normalização Z em cada uma delas (z-s1 e z-s2), elas são praticamente idênticas, }
	\label{fig:z_norm}
\end{figure}
 Em ~\parencite{trillions} é mencionado que em tarefas de classificação com séries temporais não normalizadas e Z-normalizadas, obteve-se uma taxa de erro 50\% maior no primeiro caso em relação ao segundo quando é inserido um \emph{offset} de apenas $\pm$5\%.
 
Além da normalização Z, podem ser utilizadas também as normalizações max e min-max, nas quais devem ser aplicadas, respectivamente, as seguintes transformações em cada uma das $n$ medições da série $\bf{x}$:

\begin{equation}
\hat{x_i} = \frac{x_i}{\max(\bf{x})}, i = 1,...,n \text{          e}
\end{equation}

\begin{equation}
\hat{x_i} = \frac{x_i-\min(\bf{x})}{\max(\bf{x}) - \min(\bf{x})},  = 1,...,n.
\end{equation}

A primeira é utilizada, em geral, quando os valores observados das séries temporais que compõem o conjunto de dados são sempre positivas, e confina a série no intervalo $[0,1]$ dividindo cada valor pelo maior valor observado em toda série. Já a segunda, também confina a série no intervalo $[0,1]$, mas necessariamente seu menor valor se torna $0$, o que nem sempre ocorre na primeira. Os efeitos de ambas normalizações podem ser vistos na figura ~\ref{fig:norms}. Como os valores de medição dos medidores inteligentes são sempre positivos, a normalização max é muito utilizada na mineração de dados de Redes Elétricas Inteligentes ~\parencite{Chicco}.

\begin{figure}[h!]
	\includegraphics[width=\linewidth]{figuras/normalizations.png}
	\caption{Exemplos de normalização max e min-max.}
	\label{fig:norms}
\end{figure}

\subsection{\emph{Piecewise aggregate approximation} - 
	PAA}

A \emph{Piecewise aggregate approximation} é uma técnica de redução de dimensionalidade na qual, basicamente, a série temporal é dividida em $w$ janelas de tamanhos iguais e o valor da série neste intervalo é definido como o valor médio da série no intervalo compreendido pela janela ~\parencite{SAX}. Assim, dada uma série temporal $\bm{x} = x_1,x_2,...,x_n$ de tamanho $n$, esta pode ser representada em  um espaço $w$ dimensional, sendo $w<n$, por um vetor $\bm{\bar{x}} =\bar{x_1},\bar{x_2},...,\bar{x_w} $, tal que seu $i$-ésimo elemento é dado por

\begin{equation}
	\bar{x_i} = \frac{w}{n} \sum_{j=\frac{n}{w}(i-1)+1}^{\frac{n}{w}i} x_j.
\end{equation}
%\begin{equation}
%	\bar{x_i} = \frac{w}{n} \mathop{\mathlarger{\mathlarger{\mathlarger{\mathlarger{\sum}}}}}_{j=\frac{n}{w}(i-1)+1}^{\frac{n}{w}i} x_j.
%\end{equation}
Na figura \ref{fig:PAA} pode ser vista a redução de dimensionalidade da senoide em azul formada por 315 pontos para a curva em vermelho, formada por 19 pontos. Em ~\parencite{fuzzy_transform} é sugerida uma técnica de redução de dimensionalidade de séries temporais que pode ser vista como uma versão fuzzy da PAA. Basicamente, a ideia é definir uma partição fuzzy para o espaço de entrada e definir um coeficiente para cada valor da partição.

\begin{figure}[h!]
	\includegraphics[width=\linewidth]{figuras/PAA.png}
	\caption{Exemplo de redução de dimensionalidade pela PAA.}
	\label{fig:PAA}
\end{figure}




%\subsubsection{Symbolic Aggregate Approximation - SAX}


%\hl{TODO...}

\subsection{Transformada discreta de Fourier - DFT}

Ao se calcular a série de Fourier para os vetores $\bm{x}$ e $\bm{y}$, a distância entre os dois vetores pode ser obtida pelo cálculo de qualquer umas das métricas de distância a serem discutidas, como a distância euclidiana, entre os coeficientes da série de Fourier discreta truncada de cada vetor 

\begin{equation}
dist(\bm{x},\bm{y}) = \bigg(\sum_{i=1}^{\theta} (\hat{x_i} -\hat{y_i})^2 \bigg)^{\frac{1}{2}},
\end{equation}
onde $\theta$ é o número de coeficientes complexos a serem utilizados, e $\hat{x_i}$ e $\hat{y_i}$ são os i-ésimos termos da transformada discreta de Fourier, ou \emph{Discrete Fourier Transform} (DFT), dos vetores $\bm{x}$ e $\bm{y}$ respectivamente.

Segundo ~\parencite{FFT}, são necessários apenas 5 coeficientes ou menos para uma boa aproximação das curvas no domínio da frequência, e a distância euclidiana no domínio da frequência é análoga à distância euclidiana no domínio do tempo pelo teorema de \emph{Parseval}, sendo, portanto, a métrica de distância mais indicada.

O método proposto tem a vantagem de realizar uma redução significativa de dimensionalidade, já que uma série temporal com 300 pontos, por exemplo, pode ser reduzida para uma análise em 5 pontos. Além disso, quando existe a presença de ruídos de alta frequência no sinal, e tais ruídos não são relevantes para o agrupamento, tal abordagem acarreta em resultados melhores.

%\subsubsection{Transformada discreta \emph{Wavelet} - DWT}

%\hl{TODO: Completar...}

%\subsubsection{Decomposição em valores singulares - SVD}

%A decomposição em valores singulares (\emph{Singular Value Decomposition}), ou SVD, é uma estratégia de redução de dimensionalidade dos dados baseada na obtenção dos autovalores e autovetores da matriz que contém os dados.  Segundo ~\parencite{Ullman}, a vantagem da SVD é a representação de conceitos que se encontram ocultos na matriz de dados originais.

%Assim, dada uma matriz $M_{nxd}$ contendo os dados, onde $n$ é o número de instâncias e $d$ é o número de dimensões das instâncias, ela pode ser decomposta em:

%\begin{equation}
%M = U \Sigma V^T,
%\end{equation}
%onde $U$ é a matriz dos autovetores de $MM^T$, enquanto $V$ e $\Sigma$ são respectivamente a matriz dos autovetores e a matriz diagonal contendo a raiz quadrada dos autovalores de $M^{T}M$.

%A diagonal da matriz $\Sigma$ deve ser formada pelos autovalores de forma ordenada, em função do seu valor absoluto, e a i-ésima coluna da matriz $V$ deverá ser o autovetor associado ao autovalor contido na i-ésima coluna da matriz $\Sigma$. Os autovalores são sempre positivos e a eles dá-se o nome de valores singulares. Assim temos, 

%\[
%\Sigma=
%\begin{bmatrix}
%\lambda_1 & 0 & ... & 0 \\
%0 & \lambda_2 & ... & 0 \\
%0 & 0 & ... & 0 \\
%0 & 0 & ... & \lambda_n 
%\end{bmatrix} \quad
%\text{e, }
%V=
%\begin{bmatrix}
%e_{1_1} & e_{2_1} & ... & e_{n_1} \\
%e_{1_2} & e_{2_2} & ... & e_{n_2} \\
%e_{1_3} & e_{2_3} & ... & e_{n_3} \\
%e_{1_4} & e_{2_4} & ... & e_{n_4} 
%\end{bmatrix}\text{,}
%\]
%onde o autovalor $\lambda_i$ está associado ao autovetor $e_i$, e $\lambda_1 \geq \lambda_2 \geq ...\geq\lambda_n > 0$.

%Caso os primeiros elementos da diagonal da matriz $\Sigma$ sejam muito maiores que os últimos elementos, então a matriz de dados $M$ pode ser reduzida para um número $r$ de dimensões, sendo $r<d$, de maneira a tornar mais claros os conceitos ocultos nos dados.

%Segundo ~\parencite{Ullman} Uma boa maneira de se escolher o número $r$, é reter os primeiros autovalores cujo a soma dos seus quadrados corresponde a pelo menos $90\%$ da soma total dos quadrados de todos os autovalores da matriz $\Sigma$. Assim, a matriz reduzida, e que representa os dados, $M_{nxr}^{'}$, é obtida por:

%\begin{equation}
%M_{nxr}^{'} = MV_{r}^{T},
%\end{equation}
%onde $V_{r}^{T}$ é a matriz $V$ somente com os seus $r$ primeiros autovetores  associados.

%Neste momento vale destacar que a decomposição em valores singulares é uma técnica de redução de dimensionalidade que é aplicada em todo o conjunto de dados, ao passo que técnicas como DFT e DWT são aplicadas a cada instância. \hl{TODO: Comentar acerca do custo computacional}


\section{Métricas de dissimilaridade} \label{sec:metricas}

Existem diversas métricas de dissimilaridade para comparar duas séries temporais, e elas divergem entre si principalmente no que diz respeito aos custos computacionais, sensibilidade à ruídos, invariância ao defasamento entre as curvas, capacidade de ser utilizada para séries temporais com diferentes tamanhos, entre outros. Além disso, as métricas podem ser divididas em dois grandes grupos, sendo o primeiro formado pelas métricas que podem ser consideradas distâncias, e aquelas que não podem. Para determinada métrica ser considerada uma distância ela deve possuir, necessariamente, as seguintes propriedades para vetores $\bm{x},\bm{y}$ e $\bm{z}$ ~\parencite{DudaHart}:

%DudaHart-pg 31

\begin{itemize}
	\item não-negatividade:
	\begin{equation}
	dist(\bm{x},\bm{y}) > 0, \forall \text{ }  \bm{x},\bm{y} \in \mathbb{R}^{n}\text{,}
	\end{equation}
	\item reflexividade:
	\begin{equation}
	dist(\bm{x},\bm{y}) = 0\text{, se, e somente se, }\bm{x}=\bm{y}\text{,}
	\end{equation}
	\item simetria:
	\begin{equation}
	dist(\bm{x},\bm{y}) = dist(\bm{y},\bm{x})\text{,}
	\end{equation}
	\item desigualdade do triângulo:
	\begin{equation}
	dist(\bm{x},\bm{z}) + dist(\bm{z},\bm{y}) \geq dist(\bm{x},\bm{y}).
	\end{equation}
\end{itemize}

Nas subseções seguintes desta seção, encontram-se algumas métricas que atendem à todas ou à algumas dessas propriedades e que serão detalhadas em cada uma das características já expostas e esperadas de uma métrica de similaridade entre duas séries temporais. Ao fim do capítulo se encontra a Tabela ~\ref{tbl:dissimlarity_summary_table} que sumariza todas  as informações discutidas nesta seção. 

Vale mencionar, neste momento, que existem métricas de dissimilaridade obtidas por meio de modelos estatísticos das séries temporais que sejam representativos de cada uma das séries temporais a serem comparadas. Supõe-se que elas foram geradas por modelos paramétricos, sendo o modelo ARIMA um dos modelos mais utilizados. A obtenção da medida de dissimilaridade entre as séries se dá a partir da comparação entre os parâmetros desses modelos. As métricas obtidas a partir de modelos estatísticos possuem aplicação maior em agrupamentos de séries temporais baseados na predição delas, não sendo assim, em geral, aplicáveis à problemas de agrupamento de curvas de carga, e por isso não serão abordadas nesta dissertação de mestrado. No entanto, algumas das abordagens neste sentido, assim como outras referências, podem ser encontradas em ~\parencite{TSCLUST}.

\subsection{Minkowski}

A distância de Minkowski é definida como
\begin{equation}
dist(\bm{x},\bm{y}) = \bigg (\sum_{i}^{n} |x_i -y_i|^p \bigg)^{\frac{1}{p}}.
\end{equation}

A métrica de Minkowski é também conhecida como norma $L_p$, e tem como casos notáveis a distância Euclidiana, ou $L_2$, quando $p=2$, a distância \emph{Manhattan}, ou $L_1$, quando $p=1$ e a distância \emph{Chebyshev}, ou $L_\infty$, quando $p\to\infty$.

\begin{figure}[h!]
	  \includegraphics[width=\linewidth]{figuras/dist_euclidiana.png}
	  \caption{Distância Euclidiana entre as duas curvas é dada pela soma das distâncias euclidianas entre as observações de cada curva no mesmo instante.}
	  \label{fig:dist_euclidiana}
\end{figure}

Na figura ~\ref{fig:dist_euclidiana} encontram-se duas séries temporais, na tabela ~\ref{minkowski_table} os valores de distância obtidos variando-se o valor de $p$. Pela tabela, pode-se notar que pequenas variações no valor de $p$ acarretam em grandes variações no valor da métrica, bem como que este valor diminui ao passo que $p$ aumenta.A métrica de Minkowski atende à todas as propriedades das métricas de distância e, dessa maneira, pode ser considerada uma distância, e o seu custo computacional é $\mathcal{O}(n)$, onde $n$ é o número de amostras ou observações das séries a serem comparadas.


\begin{table}[]
	\centering
	\caption{Distâncias Minkowski entre as curvas da figura ~\ref{fig:dist_euclidiana} para diversos valores de $p$.}
	\label{minkowski_table}
	\begin{tabular}{|c|c|c|c|c|c|c|}
		\hline
	%	\toprule
			%\begin{tabular}{llllllll}
						& p=1    & p=2   & p=3   & p=4 &p=5 & p$\to \infty$ \\ \hline
%		\midrule
		distância & 318.18 & 86.88 & 59.59 & 50.54 & 44.03 & 38.14   \\ \hline
	%	\bottomrule
	\end{tabular}
\end{table}


\subsection{\emph{Dynamic Time Warping} - DTW}

A \emph{Dynamic Time Warping} é, talvez, a métrica de dissimilaridade mais utilizada para dados de séries temporais. Proposta inicialmente para ser utilizada em problemas de detecção de fala, a DTW se mostrou útil para tarefas de mineração de dados em séries temporais ~\parencite{DTW}.

Dada duas séries temporais $\bm{x} = x_1,x_1,x_3,...,x_i,...,,x_n$ e $\bm{y} = y_1,y_2,y_3,...,y_j,...,,y_m$, suponha que seja feito a construção de um \emph{grid} $n_{x}m$, no qual na posição $(i,j)$ se encontra o valor resultante de uma métrica de distância $\delta(i,j)$ entre os pontos $x_i$ e $y_j$. Tal métrica pode ser, por exemplo, a já citada distância Minkowski. A partir deste \emph{grid}, é construído um caminho $W=w_1,w_2,w_3,...,w_k$, onde cada $w_k$ corresponde à alguma posição $(i,j)_k$ do grid, tal que a distância DTW entre as séries $\bm{x}$ e $\bm{y}$, definida por esse caminho, seja minimizada:

\begin{equation}
DTW(\bm{x},\bm{y}) = min_W[\sum_{k=1}^{p}\delta(w_k)]
\end{equation}

Assim, pode-se definir o cálculo da dissimilaridade DTW por meio da programação dinâmica da variável $\gamma(i,j)$ que representa a soma acumulada da distância escolhida na posição $(i,j)$ do \emph{grid}:

\begin{equation} \label{gamma}
\gamma(i,j) = \delta(i,j) + min[\gamma(i-1,j),\gamma(i-1,j-1),\gamma(i,j-1)].
\end{equation}

O cálculo se inicia na posição $(m,n)$, e a partir da fórmula recursiva apresentada na equação ~\ref{gamma}, constrói-se uma tabela de distâncias acumuladas. O caminho mínimo, que é a distância DTW entre as séries temporais, é obtido realizando o caminho reverso, ou seja, com o valor contido em $\gamma(m,n)$.

A ordem de complexidade do algoritmo é $\mathcal{O}(n*m)$, no entanto, a adição de algumas restrições fazem com que se alcance ganhos computacionais com pouca perda de performance. Uma primeira tentativa, é a criação de uma janela de tamanho $\omega$ tal que, $|i_k-j_k| \leq \omega$. Note que para $\omega=1$, a DTW se torna a distância euclidiana. Em ~\parencite{BatistaComparativo}, sugere-se a utilização de uma janela com tamanho de aproximadamente $10\%$ do tamanho da menor série temporal. Além da redução no tempo de processamento, existem estudos que, a partir de experimentos empíricos, afirmam que a utilização de janelas melhoram a performance de classificadores de séries temporais ao não permitirem mapeamentos patológicos entre uma pequena seção de uma curva à uma grande seção da outra ~\parencite{LB_Keogh}. Na figura ~\ref{fig:dtw_compare} pode ser vista uma comparação entre as métricas DTW, com e sem janela, e a Euclidiana, para duas curvas do conjunto de dados ECG200 do \emph{UCR Time Series} \parencite{UCRArchive}.

\begin{figure}
	%\begin{subfigure}{.32\textwidth}
		\centering
		\includegraphics[width=.95\linewidth]{figuras/ECG200_dtw_euclidiana.png}
		\caption{Distância Euclidiana entre duas curvas do conjunto de dados ECG200.}
		\label{fig:sfig1}
%	\end{subfigure}%
\end{figure}
\begin{figure}
%	\begin{subfigure}{.32\textwidth}
		\centering
		\includegraphics[width=.95\linewidth]{figuras/ECG200_dtw_nowindow.png}
		\caption{DTW sem janela entre duas curvas do conjunto de dados ECG200, na qual foi usada como distância $\delta(x_i,y_j) = (x_i-y_j)^2$..}
		\label{fig:sfig2}
%	\end{subfigure}
\end{figure}
\begin{figure}
	%	\begin{subfigure}{.32\textwidth}
			\centering
			\includegraphics[width=.95\linewidth]{figuras/ECG200_dtw_10_window.png}
			\caption{DTW com janela de tamanho $10\%$ do tamanho da série temporal, entre duas curvas do conjunto de dados ECG200, na qual foi usada como distância $\delta(x_i,y_j) = (x_i-y_j)^2$.}
			\label{fig:sfig3}
		%	\end{subfigure}
%	\caption{Comparação entre as diferentes estratégias para a obtenção da distância euclidiana e dissimilaridade DTW entre duas curvas do conjunto de dados ECG200, na qual foi usada como distância $\delta(x_i,y_j) = (x_i-y_j)^2$. }
	\label{fig:dtw_compare}
\end{figure}

A DTW não pode ser considerada uma métrica de distância pois não obedece à desigualdade do triângulo, no entanto, diferentemente da distância Minkowski, pode ser aplicada à series temporais de diferentes tamanhos e não é sensível ao deslocamentos entre as séries temporais.

Na literatura, a distância Euclidiana é uma das abordagens mais utilizadas para o agrupamento em geral, no entanto, em alguns casos nos quais os dados são séries temporais, ela pode apresentar resultados incoerentes ao se definir a similaridade entre duas séries temporais. Para discutir tal afirmação, os experimentos realizados em ~\parencite{FalhaEuclideana} serão reproduzidos aqui. 

\begin{table}[]
	\centering
	\caption{Distâncias euclidiana e DTW entre as curvas da figura ~\ref{fig:euclidiana_falha}.}
	\label{euclidiana_vs_DTW}
	\begin{tabular}{|c|c|c|}
%		\toprule
			\hline
										 				& Euclidiana& DTW\\ 
%		\midrule										 				
			\hline
		(ts1,ts2) 	& 26.95 	& 1.51 \\
					\hline					
		(ts1,ts3) 			& 23.18 	& 1.85 \\
					\hline
		%bottomrule
	\end{tabular}
\end{table}

\begin{figure}[h!]
	\includegraphics[width=\linewidth]{figuras/euclidiana_falha.png}
	\caption{Duas senóides e uma reta têm suas distâncias comparadas.}
	\label{fig:euclidiana_falha}
\end{figure}
 Na figura ~\ref{fig:euclidiana_falha}, encontram-se três séries temporais e, fixando-se a curva ts1, para compará-la com as outras duas, foram calculadas as distâncias euclidianas e DTW. Os resultados podem ser vistos na tabela ~\ref{euclidiana_vs_DTW}. Pelos dados da tabela, fica claro que a métrica euclidiana falha ao comparar as curvas em análise, já que a distância euclidiana entre a senóide ts1 e a reta ts3 foi menor que a distância euclidiana entre a senoide ts1 e a outra senoide ts2, o que contradiz o nosso senso de similaridade, que é baseado no formato da curva. Por sua vez, a métrica DTW nos diz que a senoide ts1 é mais "próxima"  à senoide ts2 do que à reta ts3, o que condiz com nosso senso de similaridade. A vantagem da DTW em relação à Euclidiana, neste caso, está no fato de a DTW ser invariante à distorção de fase que as curvas podem apresentar entre si.

\subsection{\emph{LB-Keogh}}

A DTW é reconhecida como uma das métricas mais robustas na obtenção da similaridade entre séries temporais, no entanto, seu elevado custo computacional faz com que esta seja preterida em muitas aplicações. Em uma tentativa de se aproximar a DTW com um custo computacional $\mathcal{O}(n)$, em ~\parencite{LB_Keogh} foi proposta a métrica LB-Keogh. No trabalho citado, tal métrica teve resultados próximos à DTW para diversas bases de dados e em tempo significativamente menor. Assim, dada duas séries temporais $\bm{x}$ e $\bm{y}$ de comprimento $n$, então, define-se as séries $\bm{u}$ e $\bm{l}$ como:

\begin{equation} \label{eq:lb_keogh_ini}
\begin{cases}
u_i = \max(y_{i-r}:y_{i+r}),\\
l_i = \min(y_{i-r}:y_{i+r}),
\end{cases}
\end{equation}
onde $r$ é definido como o alcance da aproximação e $1 \leq r \leq n$. As novas séries $\bm{u}$ e $\bm{l}$ formam um envelope que contém a série $\bm{y}$ sendo o limite superior desta formada pela curva $\bm{u}$ e o inferior pela curva $\bm{l}$. Dessa maneira, faz-se necessário que $\forall i \textbf{   } u_i \geq y_i \geq l_i$. Após a definição das duas curvas que formam o envelope, pode-se definir, finalmente, a distância \emph{LB-Keogh} entre as curvas $\bm{x}$ e $\bm{y}$:

\begin{equation} \label{eq:lb_keogh_final}
dist_{LB-Keogh}(\bm{x},\bm{y}) = \sqrt{\sum_{i}^{n} 
	\begin{cases}
	(x_i-u_i)^2,\text{ se }x_i > u_i,\\
	(x_i-l_i)^2,\text{ se }x_i < l_i,,\\
	0, \text{          caso contrário}.\\
	\end{cases}
	}
\end{equation}

A métrica LB-Keogh é sempre menor ou igual ao valor obtido pela DTW, e ao contrário desta última, pode ser aplicada somente em séries temporais de mesmo comprimento e também não pode ser considerada uma métrica de distância pois não obedece à desigualdade do triângulo. Um exemplo de cálculo da LB-Keogh pode ser visto nas figuras ~\ref{fig:keogh_1},~\ref{fig:keogh_2} e ~\ref{fig:keogh_3}.
\begin{figure} 
	\centering
	\includegraphics[height=9cm,keepaspectratio]{figuras/lb_keogh_1.png}
	\caption{Curvas $\bm{x}$ e $\bm{y}$.}
	\label{fig:keogh_1}
\end{figure}
\begin{figure}
	\centering
	\includegraphics[height=9cm,keepaspectratio]{figuras/lb_keogh_2.png}
	\caption{Curva $\bm{y}$ envelopada pelas curvas $\bm{u}$ e $\bm{l}$ obtidas pela equação ~\ref{eq:lb_keogh_ini}.}
	\label{fig:keogh_2}
\end{figure}
\begin{figure}
	\centering
	\includegraphics[height=9cm,keepaspectratio]{figuras/lb_keogh_3.png}
	\caption{A raiz do somatório das distâncias em preto fornece a distância LB-Keogh, como descrito pela equação ~\ref{eq:lb_keogh_final}.}
	\label{fig:keogh_3}
\end{figure}

\subsection{\emph{Edit distance on real sequence} - EDR}

Em ~\parencite{EDR}, é proposta uma nova técnica a ser aplicada para a obtenção de métricas de similaridade entre séries temporais multivariadas, com aplicação na análise de trajetórias de objetos móveis. Neste artigo, é defendido que a \emph{Edit distance on real sequence} é uma métrica mais robusta e menos sensível à ruídos e \emph{outliers} que outras métricas tradicionais, como a DTW e a distância Euclidiana.

A métrica é baseada na \emph{Edit distance}, ou distância de \emph{Levenshtein}, que possui larga aplicação na obtenção de métricas de dissimilaridades entre \emph{strings}. A EDR é uma adaptação desta distância a ser aplicada à uma sequência de números reais, ao invés de ser aplicada à uma sequência de caracteres. Apesar de ter sido proposta inicialmente para séries temporais multivariadas, em ~\parencite{Serra} são feitos experimentos com séries temporais univariadas nos quais se chegou a resultados superiores aos resultados de métricas tradicionais como a DTW e distância Euclidiana. Assim como a DTW, a métrica EDR não é sensível ao deslocamento entre as curvas e pode ser utilizada para curvas de comprimentos diferentes.

A EDR de duas curvas de tamanho $m$ e $n$ pode ser obtida por meio da seguinte programação dinâmica:

\begin{equation}
D_{ij} = 
\begin{cases} 
D_{i-1,j-1} 											   & \text{se } c(x_i,y_j) = 1, \\
1 + min(D_{i,j-1},D_{i-1,j},D_{i-1,j-1})     & \text{se } c(x_i,y_j) = 0, \\
\end{cases}
\end{equation}
onde $i=1,...,m$ e $j=1,....n$. A função de casamento $c(x_i,y_j) = \theta (\varepsilon - f(x_i,x_j))$, sendo $\theta(.)$ a função degrau unitário, $\varepsilon \in [0,\infty)$ é um parâmetro a ser escolhido, e sugerido por ~\parencite{EDR} como um quarto da variância da série que possui a maior variância das suas observações, e a função $f(.)$ é normalmente definida como o valor absoluto da diferença entre as variáveis, ou seja, $f(x_i,y_j) = |x_i-y_j|$. Além disso, a matriz $D$ deve ser iniciada da seguinte maneira: $D_{i,0} = i$ para $i=0,1,...,m$ e $D_{0,j} =j$ para $j=0,1,...,n$, e a distância EDR entre as curvas é dado pelo valor contido em $D_{m,n}$.

A EDR não pode ser considerada uma métrica de distância pois não possui a propriedade de desigualdade do triângulo e tem complexidade computacional $\mathcal{O}(m*n)$. Dado o seu elevado custo computacional, restrições devem ser feitas ao cálculo da dissimilaridade.

\subsection{\emph{Dice}, \emph{Jaccard} e \emph{Lorentzian}}

As distâncias \emph{Dice}, \emph{Jaccard} e \emph{Lorentzian} obtiveram performance superior à distância euclidiana em um estudo empírico de classificação de séries temporais realizado em ~\parencite{BatistaComparativo}, e por isso são utilizadas neste trabalho e definidas a seguir. Para duas séries temporais $\bm{x}$ e $\bm{y}$ de comprimento $n$, temos:

\begin{itemize}
	\item Dice \begin{equation}
	dist(\bm{x}, \bm{y}) =\frac{\sum_{i}^{n}(x_i-y_i)^2}{\sum_{i}^{n}x_i^2 + \sum_{i}^{n} y_i^2}
	\end{equation}
		\item Jaccard \begin{equation}
		dist(\bm{x}, \bm{y}) =\frac{\sum_{i}^{n}(x_i-y_i)^2}{\sum_{i}^{n}(x_i^2 +y_i^2-x_i y_i)}
		\end{equation}
			\item Lorentzian \begin{equation}
			dist(\bm{x}, \bm{y}) = \sum_{i}^{n} ln(1 + |x_i-y_i|)
			\end{equation}
\end{itemize}

Todas as três distâncias possuem ordem de complexidade $\mathcal{O}(n)$, e são aplicáveis somente à séries temporais com o mesmo comprimento $n$.

\subsection{\emph{Cosine}}

A métrica \emph{Cosine} é definida  a partir do cosseno do ângulo entre os vetores $\bm{x}$ e $\bm{y}$ das séries temporais:
\begin{equation}
dist(\bm{x},\bm{y}) = 1- \frac{\bm{x}.\bm{y}}{||\bm{x}||_2 ||\bm{y}||_2},
\end{equation}
onde $||\bm{x}||_2$ é a norma euclidiana do vetor $\bm{x}$. A distância \emph{Cosine} tem ordem de complexidade $\mathcal{O}(n)$ e só pode ser calculada em entre séries temporais de mesmo tamanho $n$.

\subsection{Distância de Mahalanobis}

A distância de Mahalanobis é uma distância estatística proposta em ~\parencite{mahalanobis1936generalized}, que leva em conta a covariância entre as amostras da base de dados. Útil na detecção de \emph{outliers}, a distância de Mahalanobis utiliza a matriz de covariância $\Sigma$ das amostras, e, para o caso específico no qual $\Sigma=I$, então a métrica é análoga à distância euclidiana. Assim, a distância de Mahalanobis entre os vetores $\bm{x}$ e $\bm{y}$, pertencentes ao conjunto de dados que gerou a matriz $\Sigma$, é:

\begin{equation}
dist(\bm{x},\bm{y}) = \sqrt{(\bm{x}-\bm{y})^T\Sigma^{-1}(\bm{x}-\bm{y})}.
\end{equation}

Os maiores custos computacionais da obtenção da distância de Mahalanobis residem no cálculo da matriz de covariância $\Sigma$ e da inversão da mesma. Além do custo computacional em si, numericamente, a inversão da matriz $\Sigma$ é problemática pois pode levar a diversos erros numéricos já que, em geral, a matriz $\Sigma$ é de baixo \emph{rank} para conjuntos de dados de séries temporais. O cálculo e inversão da matriz $\Sigma$ têm ordem de complexidade $\mathcal{O}(n^2)$, e a métrica só pode ser aplicada à séries temporais de mesmo tamanho.

Em ~\parencite{mahalanobisClassification}, foi realizado um estudo de classificação de séries temporais no qual a distância de Mahalanobis foi utilizada como métrica de dissimilaridade entre as curvas. Foram realizadas aproximações da inversa da matriz de covariância dos dados que fizeram com que se tivessem ganhos computacionais relevantes, possuindo, assim, custos computacionais inferiores aos custos da DTW, no entanto, com uma pequena perda de performance na acurácia.

\subsection{Distâncias baseadas na correlação de Pearson}

A correlação de Pearson entre duas séries temporais $\bm{x}$ e $\bm{y}$ te tamanho $n$ é dada por

\begin{equation} \label{d_COR}
COR (\bm{x},\bm{y}) = \frac{\sum_{i}^{n}(x_i-\bm{\bar{x}}) (y_i-\bm{\bar{y}})}{\sqrt{\sum_{i}^{n}(x_i-\bm{\bar{x}})^2}\sqrt{\sum_{i}^{n}(y_i-\bm{\bar{y}})^2}},
\end{equation}
sendo $\bar{\bm{x}}$ e $\bar{\bm{y}}$ as média das observações de $\bm{x}$ e $\bm{y}$ respectivamente. O valor da correlação de Pearson é compreendido no intervalo [$-1,1$]. Dessa maneira, as seguintes distância baseadas na correlação de Pearson podem ser definidas

\begin{equation}
dist_1(\bm{x},\bm{y})  = 1-COR(\bm{x},\bm{y}),
\end{equation}

\begin{equation}
dist_2(\bm{x},\bm{y})  = \sqrt{2(1-COR(\bm{x},\bm{y}))},
\end{equation}

\begin{equation}
dist_3(\bm{x},\bm{y}) = \sqrt{\Bigg(\frac{1-COR(\bm{x},\bm{y})}{1+COR(\bm{x},\bm{y})}\Bigg)^\beta},
\end{equation}
onde $\beta$ é um parâmetro para regular o decrescimento da distância $dist_3$. Todas as métricas podem ser aplicadas somente entre séries temporais de mesmo comprimento e possuem ordem de complexidade $\mathcal{O}(n)$, onde $n$ é o número de observações das séries temporais em comparação. 

\subsection{Distância baseada na correlação temporal}

Uma medida da correlação temporal entre duas séries temporais definida em ~\parencite{cort}, é equivalente à uma medida do quão similar é o comportamento dinâmico entre as séries, ou, em outras palavras, o quão próximo está o instante de tempo no qual elas crescem ou decrescem. Definida no intervalo $[-1,1]$ a correlação temporal entre as séries $\bm{x}$ e $\bm{y}$ de tamanho $n$ pode ser obtida por

\begin{equation} \label{d_CORT}
CORT (\bm{x},\bm{y}) = \frac{\sum_{i}^{n-1}(x_i-x_{i+1}) (y_i-y_{i+1})}{\sqrt{\sum_{i}^{n-1}(x_i-x_{i+1})^2}\sqrt{\sum_{i}^{n-1}(y_i-y_{i+1})^2}},
\end{equation}

Valores de correlação temporal próximos de $1$ indicam que as séries $\bm{x}$ e $\bm{y}$ crescem e decrescem simultaneamente e com a mesma taxa, enquanto valores próximos de $-1$ indicam que uma série cresce ao passo que a outra decresce simultaneamente, e isso ocorre com a mesma taxa. Valores próximos de $0$ indicam que estatisticamente não se pode afirmar nada a respeito da taxa e do crescimento ou decrescimento de uma curva ao saber o comportamento da outra.

A partir da equação ~\ref{d_CORT} criar uma métrica que leva em conta tanto aspectos temporais como a correlação temporal quanto a proximidade ponto a ponto entre as séries. Para tal, faz-se uso de uma função de sintonia e de alguma métrica convencional, como a métrica Euclidiana ou DTW, por exemplo. Assim, uma possível métrica de distância que leve em conta o aspecto temporal  entre as séries $\bm{x}$ e $\bm{y}$ seria:

\begin{equation}
dist(\bm{x},\bm{y}) = f(CORT (\bm{x},\bm{y})).dist_Q(\bm{x},\bm{y}) 
\end{equation}
onde $dist_Q$ é a distância convencional entre as séries e $f(.)$ é a função de sintonia que deve amplificar a distância convencional quando a correlação temporal se encontra no intervalo $[-1,0[$, reduzi-la quando a correlação temporal se encontra no intervalo $]0,1]$, e não alterá-la quando a correlação temporal for igual à $0$. A seguinte função, proposta em  ~\parencite{cort}, atende à estes critérios:

\begin{equation} \label{eq:tunning_func}
f(x) = \frac{2}{1+\exp(\gamma x)},\text{ 				} \gamma \geq 0,
\end{equation}
onde $\gamma$ é um parâmetro de controle que controla o quanto o aspecto temporal deve influenciar na métrica de distância. Valores grandes de $\gamma$ fazem com que o valor fornecido por $dist_Q$ seja mais penalizado, ou atenuado, quanto se obtém valores não nulos de correlação temporal entre as séries. Tal comportamento da função de sintonia bem como a influência do parâmetro $\gamma$ podem ser vistos na figura ~\ref{fig:tunning_cort}.

Em ~\parencite{cort} são feitos experimentos com séries temporais reais e sintéticas, utilizando-se métricas baseadas na correlação de Pearson e na correlação temporal, e como distâncias convencionais foram escolhidas a DTW e a distância Euclidiana. Após a realização do agrupamento das curvas por meio de algoritmos hierárquicos, utilizando-se as métricas citadas anteriormente, chegou-se à resultados superiores com a correlação temporal utilizando a similaridade DTW como distância convencional. Este resultado superior foi alcançado tanto para os conjuntos de dados sintéticos quanto para o conjunto de dados reais.

A correlação temporal não pode ser considerada uma métrica de distância pois não obedece à desigualdade do triângulo e como o cálculo da correlação temporal tem custo $\mathcal{O}(n)$, seu custo computacional se torna, em geral, o mesmo da distância convencional utilizada. 

\begin{figure}[h!]
	\includegraphics[width=\linewidth]{figuras/tunning_cort.png}
	\caption{Função de sintonia ~\ref{eq:tunning_func} para diferentes valores de $\gamma$.}
	\label{fig:tunning_cort}
\end{figure}

\subsection{\emph{Complex Invariant Distance} - CID} \label{sec:CID}

Em alguns casos, séries temporais "complexas", dependendo da métrica a ser utilizada, podem se tornar mais similares à uma outra série temporal "simples"{} do que à outra série tão "complexa"{} quanto ela. Por complexidade, entende-se uma série temporal com muitos picos e vales ~\parencite{CID}. A \emph{Complex Invariant Distance}, ou CID, leva em conta estimativas de complexidade das séries temporais para realizar uma correção na medida da dissimilaridade e, dessa maneira, para duas séries $\bm{x}$ e $\bm{y}$

\begin{equation}
dist_{CID} = dist_{Q}(\bm{x},\bm{y}) \text{ . } CF(\bm{x},\bm{y}),
\end{equation}
onde $dist_{Q}(.)$ é qualquer uma das dissimilaridades já vistas, como a DTW ou Euclidiana, e $CF(.)$ é um fator de correção que leva em conta uma estimativa da complexidade relativa entre as séries. Note que, pelo fato do fator de correção ser relativo, duas séries que tiverem complexidades semelhantes tenderão a ter um valor muito próximo de $dist_{Q}$, ao passo que séries que possuem complexidades distintas, penalizarão a dissimilaridade $dist_{Q}$.

Em ~\parencite{CID}, é proposto o fator de correção
\begin{equation} \label{eq:CF}
CF (\bm{x},\bm{y}) = \frac{max(CE(\bm{x}),CE(\bm{y}))}{min(CE(\bm{x}),CE(\bm{y}))},
\end{equation}
onde $CE(.)$ é um estimador de complexidade da série temporal.  A complexidade das séries pode ser estimada de diferentes maneiras, como: pela complexidade de Kolmogorov, métricas de entropia, número de trocas do sinal da derivada, entre outras citadas em ~\parencite{TSCLUST}. A medida de complexidade a ser utilizada neste trabalho, sugerida em ~\parencite{CID}, é bastante simples e intuitiva e será obtida pelo comprimento do tamanho da série temporal após esta ser "esticada", e é obtida por 

\begin{equation} \label{eq:CE}
CE(\bm{x}) = \sqrt{\sum_{i}^{n-1} (x_i - x_{i+1})^2}.
\end{equation}

O cálculo do fator de correção possui complexidade $\mathcal{O}(n)$, o que faz com que o custo computacional da métrica seja, em geral, igual ao custo da dissimilaridade $dist_{Q}$, e a pergunta se a dissimilaridade CID pode ou não ser aplicada à séries temporais com diferentes comprimentos, deve ser respondida pela pergunta se a métrica $dist_{Q}$ o faz.

A seguir é exibido um exemplo no qual se faz necessário levar em conta a complexidade entre as séries temporais. Na figura ~\ref{fig:cid} podem ser vistas três séries temporais, sendo uma simples, $ts1$, e duas complexas, $ts2$ e $ts3$. Foram calculadas as distâncias euclidianas e as distâncias euclidianas corrigidas pelo fator de complexidade da Equação ~\ref{eq:CF}. Além dos valores das distâncias euclidianas e CID-euclidiana, na Tabela ~\ref{cid_table}, na primeira coluna, podem ser vistas as complexidades de cada curva. 

Note que, como esperado, as curvas com maior ruído, ou com maior número de picos e vales, possuem complexidade aproximadamente quatro vezes maior do que a curva mais suave, ou mais simples. A submatriz de distância euclidiana entre as curvas nos mostra que as curvas complexas, $ts2$ e $ts3$, estão mais distantes entre si do que da curva mais simples, $ts1$. Tal resultado contradiz nossa intuição de proximidade entre as três curvas analisadas. No entanto, após o fator de correção que leva em conta a complexidade, a distância euclidiana foi corrigida, já que a distância CID-euclidiana entre as curvas complexas é menor que entre cada uma e a curva mais simples, $ts1$.

Este resultado é, em geral, o mesmo para outras métricas de dissimilaridade como a DTW. Em ~\parencite{CID} foi realizado o agrupamento hierárquico de formas geométricas, após estas serem convertidas para séries temporais (vide Figuras ~\ref{fig:cid_1} e ~\ref{fig:cid_2}), e a métrica CID obteve resultados melhores que os resultados com a distância Euclidiana pura, tanto para dados sintéticos como para dados reais. No teste de classificação empírico realizado em ~\parencite{BatistaComparativo}, a CID, utilizando a DTW como distância convencional, obteve a melhor acurácia dentre outras 45 métricas utilizada, para 20 conjuntos de dados de séries temporais diferentes.

\begin{figure}[h!]
	\includegraphics[width=\linewidth]{figuras/CID.png}
	\caption{Três curvas diferentes sendo duas complexas e uma simples. A estimativa de complexidade de cada uma bem como as distâncias CID e euclidiana entre elas podem ser vistas na tabela ~\ref{cid_table}}
	\label{fig:cid}
\end{figure}

\begin{table}[]
	\centering
	\caption{Medida de complexidade, como definida na equação ~\ref{eq:CE}, das curvas da figura ~\ref{fig:cid} e distâncias euclidiana e CID-euclidiana entre elas.}
	\label{cid_table}
	\begin{tabular}{c|c|ccc|ccc}
		%\begin{tabular}{llllllll}
		\toprule
		& Complex.  & Eucl.(ts1)   & Eucl.(ts2)   &Eucl.(ts3)  & CID.(ts1)   & CID.(ts2)   &CID.(ts3) \\
		\midrule

ts1 & 0.560 & 0.0 & 1.408 & 1.354 & 0.0 & 5.090 & 4.777\\
ts2 & 2.025 & 1.408 & 0.0 & 2.027 & 5.090 & 0.0  & 2.078\\		
ts3 & 1.975 & 1.354 & 2.027 & 0.0 & 4.777 & 2.078  & 0.0\\
		\bottomrule
	\end{tabular}
\end{table}

\subsection{Discussão geral das métricas de dissimilaridade}

Após a exposição das diversas métricas de dissimilaridade, faz-se necessária a sumarização de todas as informações nesta seção, que pode ser vista na Tabela ~\ref{tbl:dissimlarity_summary_table}, bem como de uma discussão geral acerca das métricas de dissimilaridade entre séries temporais. A escolha da métrica a ser utilizada no agrupamento depende sensivelmente do domínio do problema e de como a série temporal foi gerada. Para algumas aplicações, em que o formato da série é relevante, métricas que medem a dissimilaridade ponto a ponto são mais indicadas. Em contrapartida, para aplicações onde o interesse é verificar se as séries crescem ou decrescem de forma similar, independentemente do seu formato, métricas baseadas em correlação temporal são mais apropriadas. Alpem disso, as séries temporais a serem agrupadas, ou classificadas, podem possuir diversas distorções, devido à sua natureza ou pela maneira que elas foram obtidas, e isso deve ser levado em conta na escolha da métrica, uma vez que determinadas métricas são naturalmente invariantes a algumas distorções.. 

Em problemas de mineração de dados de séries temporais, é desejável que a métrica de dissimilaridade e o pré-processamento escolhidos alcancem uma ou mais das seguintes invariâncias, que são detalhadas em \parencite{CID} e serão brevemente discutidas nas subseções a seguir. Vale ressaltar que ao se tentar remover algumas das distorções em um problema específico, resultados piores podem ser obtidos, de forma que cada conjunto de dados de séries temporais deve ser analisado antes de se definir quais invariâncias são necessárias para a comparação entre as séries temporais.

\subsubsection{Invariância à amplitude/\emph{offset}}

Quando as séries temporais a serem comparadas são idênticas, mas medidas com diferentes escalas, estas possuem amplitudes distintas, e as métricas de dissimilaridades já discutidas podem apresentar valores muito altos, o que pode indicar erroneamente que as séries são dissimilares. Além da amplitude, séries que são idênticas, mas possuem um valor médio diferente (\emph{offset}), podem apresentar também valores muito grandes de dissimilaridade. A solução apresentada em \parencite{CID}, e discutida na seção ~\ref{sec:norm_Z}, para se alcançar a invariância à amplitude e \emph{offset} é a utilização da normalização Z nas curvas das séries temporais, em uma etapa de pré-processamento. Exemplos de curvas que devem ser comparadas com invariância à amplitude e \emph{offset} podem ser vistas na Figura ~\ref{fig:inv_offset_amplitude}.

\begin{figure}[h!]
	\includegraphics[width=\linewidth]{figuras/invariancias/amplitude_offset.png}
	\caption{À direita duas curvas, não normalizadas, que, quando comparadas, se parecem bastante distintas, quando na verdade não o são. À esquerda um eletrocardiograma, a ser comparado com um sinal de medição contínuo. As duas primeiras comparações têm um bom casamento, mas as subsequentes apresentam valores maiores em uma comparação ponto a ponto devido ao crescente \emph{offset} que a curva começa a apresentar. Extraído de ~\parencite{CID}}
	\label{fig:inv_offset_amplitude}
\end{figure}

\subsubsection{Invariância à escala local ou ao deslocamento}

Muito comum em sinais biológicos, a invariância à escala local ou ao deslocamento das séries, se manifesta quando o alinhamento ótimo entre as séries ocorre devido à um entortamento local. A dissimilaridade DTW é uma métrica invariante ao deslocamento, e tal propriedade pode ser vista na Figura ~\ref{fig:dtw_compare}

\begin{figure}[h!]
	\includegraphics[width=\linewidth]{figuras/invariancias/deslocamento.png}
	\caption{Duas curvas cujo alinhamento ótimo é obtido por meio da DTW. Extraído de ~\parencite{CID}}
	\label{fig:inv_deslocamento}
\end{figure}


\subsubsection{Invariância à escala uniforme}

Em alguns casos, duas séries temporais que são idênticas ou da mesma classe mas, por questões de medição, possuem tamanhos distintos, podem ter sua dissimilaridade afetada ao se utilizar alguma métrica de dissimilaridade como a DTW, por exemplo. A solução apresentada em \parencite{CID} para tal é a multiplicação de uma das curvas por um fator $\kappa$, que irá esticar ou reduzir uma das curvas para melhorar o casamento entre elas. A principal dificuldade em se obter a invariância à escala uniforme na comparação de séries temporais é a definição do valor de $\kappa$, que na prática, é obtido por meio de diversos testes dentro de determinado intervalo na etapa de pré-processamento. Um alinha mento ótimo obtido por meio da escala uniforme pode ser visto na Figura ~\ref{fig:inv_escala_uniforme}.

\begin{figure}[h!]
	\includegraphics[width=\linewidth]{figuras/invariancias/escala_uniforme.png}
	\caption{Duas curvas cujo alinhamento ótimo é obtido após a curva ser esticada por um fator $\kappa=1.41$, que foi obtido por força bruta. Extraído de ~\parencite{CID}.}
	\label{fig:inv_escala_uniforme}
\end{figure}

\subsubsection{Invariância ao defasamento}

A invariância ao defasamento é muito importante quando se compara duas séries temporais periódicas. Segundo \parencite{CID} algumas heurísticas foram propostas para se obter a invariância ao defasamento entre duas séries temporais, mas nenhuma delas é robusta o suficiente, de forma que tal invariância deve ser obtida por meio do teste de todos os alinhamentos possíveis entre as curvas na etapa de pré-processamento. Dessa maneira, este é um problema em aberto na literatura, e nenhuma das dissimilaridades ou técnicas de pré-processamento discutidas alcançam a invariância ao deslocamento. A Figura ~\ref{fig:inv_defasamento} demonstra a obtenção desta invariância para duas séries temporais.

\begin{figure}[h!]
	\includegraphics[width=\linewidth]{figuras/invariancias/defasamento.png}
	\caption{Duas curvas que após sucessivas tentativas de defasamento de uma delas, possuem casamento ótimo. Extraído de ~\parencite{CID}.}
	\label{fig:inv_defasamento}
\end{figure}

\subsubsection{Invariância à oclusão}

Algumas séries temporais podem apresentar subsequências que não possuem valores observados, por algum erro na medição, ou pela própria característica do problema. A invariância à oclusão pode ser obtida ao se ignorar determinadas partes das séries temporais que não possuem casamento. Algumas métricas de dissimilaridade como a DTW e EDR são mais robustas à oclusão do que a distância Euclidiana. A comparação entre as séries temporais geradas a partir dos imagens de dois crânios, na qual se faz necessária a invariância à oclusão, pode ser vistas na Figura ~\ref{fig:inv_oclusao}.

\begin{figure}[h!]
	\includegraphics[width=\linewidth]{figuras/invariancias/oclusao.png}
	\caption{Duas curvas geradas a partir do formato de dois crânios a serem comparados. A falta da região do nariz de um dos crânios pode influenciar significativamente na comparação entre as métricas. Extraído de ~\parencite{CID}.}
	\label{fig:inv_oclusao}
\end{figure}

\subsubsection{Invariância à complexidade}

A complexidade de uma série temporal, como definida na seção \ref{sec:CID}, diz respeito ao número de picos e vales que a série apresenta. Observa-se que em diferentes domínios, instâncias de diferentes classes possuem complexidades distintas, o que torna tal atributo da série temporal relevante para as tarefas de classificação e agrupamento destas. A melhor maneira de se levar em conta as complexidades das séries temporais ao compará-las, é a utilização do fator de correção CID exposto na seção ~\ref{sec:CID}. As figuras ~\ref{fig:cid_1} e ~\ref{fig:cid_2} demonstram um problema no qual a invariância à complexidade influencia os resultados em uma tarefa de agrupamento.

\begin{figure}[h!]
	\includegraphics[width=\linewidth]{figuras/invariancias/cid_1.png}
	\caption{Diversas formas geométricas e as séries temporais geradas para cada uma delas. As curvas foram geradas pela distância do ponto central da figura até cada ponto do seu contorno. Extraído de ~\parencite{CID}.}
	\label{fig:cid_1}
\end{figure}
\begin{figure}[h!]
	\includegraphics[width=\linewidth]{figuras/invariancias/cid_2.png}
	\caption{Agrupamentos hierárquicos com \emph{linkagem} average realizados nas curvas; à esquerda com a distância euclidiana e à direita com a distância euclidiana corrigida pela CID. Extraído de ~\parencite{CID}.}
	\label{fig:cid_2}
\end{figure}

\begin{table}[]
	\centering
	\caption{Sumário das métricas de dissimilaridades discutidas em ~\ref{sec:metricas}.}
	\label{tbl:dissimlarity_summary_table}
	\resizebox{\columnwidth}{!}{%
	\begin{tabular}{ccccccc}
		%\begin{tabular}{llllllll}
		\toprule
		Métrica 		& Distância? 		& Ordem de complexidade & Inv. Oclusão & Inv. Deslocamento & Inv. Complexidade & Ref. \\
		\midrule
		Minkowski 	  & 	\checkmark & $\mathcal{O}(n)$			 & {} 				 & {}						& {} 						 & \parencite{DudaHart} \\
		DTW			    & 	{}      		   & $\mathcal{O}(n_1*n_2)$ 		   & {} 			   &  \checkmark     & {} 						 & \parencite{DTW} \\
		LB-Keogh	& 	{}      		   & $\mathcal{O}(n)$ 		   & {} 			   &  \checkmark     & {} 				& \parencite{LB_Keogh} \\
		EDR			    & 	{}      		   & $\mathcal{O}(n_1*n_2)$ 		   & {} 			   &  \checkmark     & {} 						 & \parencite{EDR} \\
		Dice	& 	{}      		   & $\mathcal{O}(n)$ 		   & {} 			   & {}    & {} 				& \parencite{BatistaComparativo} \\
		Jaccard	& 	{}      		   & $\mathcal{O}(n)$ 		   & {} 			   &  {}     & {} 				& \parencite{BatistaComparativo} \\
		Lorentzian	& 	{}      		   & $\mathcal{O}(n)$ 		   & {} 			   &  {}     & {} 				& \parencite{BatistaComparativo} \\
		Cosine	& 	{}      		   & $\mathcal{O}(n)$ 		   & {} 			   &  {}     & {} 				& \parencite{BatistaComparativo} \\
		Mahalanobis	& 	{}      		   & $\mathcal{O}(n^2)$ 		   & {} 			   &  {}     & {} 				& \parencite{mahalanobis1936generalized} \\
		Pearson	& 	{}      		   & $\mathcal{O}(n)$ 		   & {} 			   &  {}     & {} 				& \parencite{TSCLUST} \\
		CORT			      & 	{}      	& $\mathcal{O}(n)$ 		   & {} 			   &  {}                      & {}				         	 & \parencite{cort} \\
		CID			      & 	{}      		  & $\mathcal{O}(n)$ 		   & {} 			   &  {}                      & \checkmark         	 & \parencite{CID} \\
		\midrule
	\end{tabular}
} % --> resizebox
\end{table}

\section{Algoritmos de agrupamento}

Dada uma métrica de dissimilaridade e um conjunto de dados, existem diferentes algoritmos que realizam a tarefa de agrupamento dos dados, e segundo ~\parencite[][243]{Ullman} estes podem ser divididos em:

\begin{enumerate}
	\item  algoritmos hierárquicos ou aglomerativos, nos quais, inicialmente, cada instância representa um grupo formado por apenas ela mesma, e a partir daí os grupos são combinados sucessivamente a partir de um critério que leva em conta a dissimilaridade entre os grupos,
	\item algoritmos particionais, ou de atribuição de pontos, nos quais cada ponto, ou instância, é atribuído inicialmente à um grupo de uma solução inicial. Em seguida, os pontos são iterativamente realocados entre os grupos, até determinado critério de parada para a obtenção da partição final. Quando os pontos são atribuídos a somente um grupo, dá-se o nome de \emph{hard-clustering}, enquanto quando atribui-se graus de pertinência do ponto à cada um dos grupos, têm-se o \emph{soft-clustering}, ou \emph{fuzzy clustering}.
\end{enumerate}
No que diz respeito ao uso de memória secundária (memória RAM), os algoritmos de agrupamento podem ser, também, divididos em:

\begin{enumerate}
	\item algoritmos que assumem que as instâncias se encontram no espaço euclidiano, ou seja, que a métrica de dissimilaridade utilizada é a distância Euclidiana. Essa distinção existe pois no espaço euclidiano existe o conceito de centróide, que é o ponto médio das instâncias que compõem o grupo, e este ponto é tido como o ponto mais representativo deste. Note que este ponto mais representativo não é, necessariamente, um ponto do conjunto de dados. Assim, em espaços não euclidianos, tal conceito de um ponto médio não existe e, dessa maneira, deve-se utilizar outras estratégias para se definir um ponto representativo dos grupos.
	\item algoritmos que assumem que os dados são pequenos, o suficiente, no que diz respeito ao consumo de memória, para serem armazenados na memória RAM. Caso isso não seja verdade, outras estratégias já partem do pressuposto que é inviável mensurar a dissimilaridade entre todos os pontos, de maneira que diversos algoritmos já são descartados. 
\end{enumerate}

Além dos critérios supracitados, os algoritmos de agrupamento podem, ainda, serem classificados em: aqueles baseados em características, nos quais exigem como entrada uma matriz $NxD$ contento os dados, que são formados por $N$ instâncias, sendo cada uma composta por $D$ características, e os algoritmos de agrupamento baseados em dissimilaridades, os quais esperam uma matriz simétrica $NxN$, onde o elemento $ij$ desta matriz contém a dissimilaridade entre o i-ésimo e o j-ésimo elemento da matriz $NxD$ que contém as instâncias. A seguir seguem breves comentários de alguns dos principais algoritmos de agrupamento utilizados no contexto de agrupamento de séries temporais.

\subsection{Algoritmos particionais}

\subsubsection{\emph{k-means}}

O \emph{k-means}, ou k-médias, ou algoritmo de Lloyd, é uma heurística que utiliza o conceito de centróide para a obtenção da partição final. Basicamente, dado um número $k$ de grupos, o objetivo do algoritmo é encontrar uma partição dos dados que minimize o soma do quadrado da distância entre as instâncias de cada grupo e seu centróide. A resolução deste problema é NP-difícil, no entanto em ~\parencite{Lloyd} é proposto um algoritmo de busca local que fez com que o \emph{k-means} se tornasse um dos algoritmos de agrupamento mais populares.

O algoritmo \emph{k-means} se encontra detalhado no pseudocódigo ~\ref{alg:k_means}, e pelo fato de ser um algoritmo de busca local, deve ser executado diversas vezes, para tentar se chegar ao mínimo global. O \emph{k-means} foi o precursor de diversas outras técnicas de agrupamento, que são variações dele ou outras técnicas realmente inspiradas no \emph{k-means}. Dentre as principais podemos destacar: \emph{k-means++} ~\parencite{k-means++} , \emph{k-medoids} ~\parencite{k-medoids} e \emph{Fuzzy-k-means } ~\parencite{fuzzy_k-means}.

Uma das desvantagens do \emph{k-means} é que se deve definir \emph{a priori} o número de grupos da partição. Essa informação nem sempre se encontra disponível de antemão, o que faz com que sejam realizadas partições com valores iterativos de $k$, e a partir dos valores dos índices de validação de agrupamento, pode-se ter uma ideia do melhor valor de $k$, dentro de determinado intervalo.

A utilização do \emph{k-means} para o agrupamento de séries temporais fica limitada para os casos nos quais a distância euclidiana foi utilizada como métrica de dissimilaridade entre as instâncias, pois somente no espaço euclidiano faz sentido a utilização de pontos médios ou centroides. Em outras palavras, isso significa que várias das métricas promissoras detalhadas na seção ~\ref{sec:metricas}, como a DTW, CID, entre outras, não podem ser utilizadas no agrupamento por meio do algoritmo \emph{k-means}.

\begin{algorithm}
	\caption{Algoritmo \emph{k-means}.}
	\label{alg:k_means}
	\begin{algorithmic}[1]
	\STATE{\textbf{Entradas:  número k de grupos, matriz X de dados com n instâncias e critério de parada $\epsilon$}}
		\STATE{\textbf{Saída:  partição $\mathcal{G}$ dos dados }}
	
	\FOR{$i \leftarrow 1$ \TO $k$}
		\STATE{$j \leftarrow rand(0,n)$} // escolhe aleatoriamente um índice da matriz de dados
		\STATE{$\mathcal{C}_i \leftarrow {x_j}$} //define o centróide como a instância escolhida;
		\STATE{$\mathcal{G}_i \leftarrow {\emptyset}$} // Inicia o i-ésimo grupo vazio
	\ENDFOR
	
	\STATE{$\phi \leftarrow \epsilon $}	// soma das distâncias de cada ponto ao seu centroide

	\WHILE{$\phi \geq \epsilon$}	
		\FOR{$p \leftarrow 1$ \TO $n$}
				\FOR{$i \leftarrow 1$ \TO $k$}
						\STATE{$d_i \leftarrow {dist(x_p,C_i)}$} // calcula a distância do ponto a cada centroide
				\ENDFOR
			
				\STATE{$ s \leftarrow {\arg \min_{i} (d_i)}$}	// Encontra o centroide mais próximo ao ponto p
				\STATE{$\mathcal{G}_s \leftarrow {x_p}$} // Associa o ponto p ao grupo do seu centroide mais próximo			
		\ENDFOR

		\STATE{$\phi \leftarrow 0$}	// soma das distâncias de cada ponto ao seu centroide
		\FOR{$i \leftarrow 1$ \TO $k$}
			\STATE{$\mathcal{C}_i \leftarrow {centroid(\mathcal{G}_i)}$} // Atualiza o centroide do i-ésimo grupo
			\FOR{$j \leftarrow 1$ \TO $n$}			
				\IF{$x_j \in \mathcal{G}_i$}
					\STATE{$\phi \leftarrow \phi + dist(x_j,C_i)$}	// soma das distâncias de cada ponto ao seu centroide
				\ENDIF
			\ENDFOR
		\ENDFOR
		
	\ENDWHILE
		
	\end{algorithmic}
\end{algorithm}

\subsubsection{\emph{K-medoids}}

Essencialmente, o \emph{k-medoids} é uma variação do \emph{k-means}, na qual a única diferença é que não se faz o uso de um ponto médio, ou centróide, das instâncias de um mesmo grupo, mas sim de um ponto mediano, ou medóide, que melhor representa as instâncias de cada grupo. O algoritmo se encontra detalhado no pseudocódigo ~\ref{alg:k_medoids}, onde pode ser visto que a única diferença em relação ao pseudocódigo ~\ref{alg:k_means} se encontra na linha $18$, ao realizar a atualização dos medóides de cada grupo, ao invés da atualização dos centróides.  Dessa maneira, diferentemente do \emph{k-means}, o \emph{k-medoids} consegue trabalhar com qualquer métrica de dissimilaridade, além da distância euclidiana.

O medóide pode ser definido de diversas maneiras, sendo a mais comum como a instância do grupo cuja a soma da distância, ou dissimilaridade, entre ela e as demais instâncias do  mesmo grupo é mínima para aquele grupo ~\parencite[][253]{Ullman}. Então dado um conjunto de dados $X$ composto por $n$ instâncias $x_i, i=0,...,n$, o medóide $x_{med}$ é definido como

\begin{equation}
x_{med} = \arg \min_{x_i} \sum_{j=0}^{n} diss(x_i,x_j), \text{			}i=0,...,n
\end{equation}

\begin{algorithm}
	\caption{Algoritmo \emph{k-medoids}.}
	\label{alg:k_medoids}
	\begin{algorithmic}[1]
		\STATE{\textbf{Entradas:  número k de grupos, matriz X de dados com n instâncias e critério de parada $\epsilon$}}
		\STATE{\textbf{Saída:  partição $\mathcal{G}$ dos dados }}
		
		\FOR{$i \leftarrow 1$ \TO $k$}
		\STATE{$j \leftarrow rand(0,n)$} // escolhe aleatoriamente um índice da matriz de dados
		\STATE{$\mathcal{C}_i \leftarrow {x_j}$} //define o medóide como a instância escolhida;
		\STATE{$\mathcal{G}_i \leftarrow {\emptyset}$} // Inicia o i-ésimo grupo vazio
		\ENDFOR
		
		\STATE{$\phi \leftarrow \epsilon $}	// soma das distâncias de cada ponto ao seu medóide
		
		\WHILE{$\phi \geq \epsilon$}	
		\FOR{$p \leftarrow 1$ \TO $n$}
		\FOR{$i \leftarrow 1$ \TO $k$}
		\STATE{$d_i \leftarrow {dist(x_p,C_i)}$} // calcula a distância do ponto a cada medóide
		\ENDFOR
		
		\STATE{$ s \leftarrow {\arg \min_{i} (d_i)}$}	// Encontra o medóide mais próximo ao ponto p
		\STATE{$\mathcal{G}_s \leftarrow {x_p}$} // Associa o ponto p ao grupo do seu medóide mais próximo			
		\ENDFOR
		
		\STATE{$\phi \leftarrow 0$}	// soma das distâncias de cada ponto ao seu medóide
		\FOR{$i \leftarrow 1$ \TO $k$}
		\STATE {\textbf{$\mathcal{C}_i \leftarrow {medoid(\mathcal{G}_i)}$}} // Atualiza o medóide do i-ésimo grupo
		\FOR{$j \leftarrow 1$ \TO $n$}			
		\IF{$x_j \in \mathcal{G}_i$}
		\STATE{$\phi \leftarrow \phi + dist(x_j,C_i)$}	// soma das distâncias de cada ponto ao seu medóide
		\ENDIF
		\ENDFOR
		\ENDFOR
		
		\ENDWHILE
		
	\end{algorithmic}
\end{algorithm}

\subsection{Algoritmos hierárquicos}

Os algoritmos hierárquicos, possuem a vantagem de não ser necessário informar, \emph{a priori}, o número de grupos que os dados se encontram divididos, no entanto, em geral, possuem ordem de complexidade de  $\mathcal{O}(n^2log(n))$, onde $n$ é o número de instâncias, e esta complexidade é superior em relação aos algoritmos baseados em atribuição de pontos. Outra característica relevante dos algoritmos de agrupamento hierárquicos é que estes são determinísticos, ou seja, para uma mesma entrada, e mesmos parâmetros, sempre retornam a mesma partição. Eles conseguem trabalhar tanto com matrizes de características como com matrizes de dissimilaridade ~\parencite[][245]{Ullman}.

Os algoritmos hierárquicos podem ser divididos em algoritmos divisivos ou aglomerativos. Nos primeiros, inicialmente é considerado um grande grupo que contém todos os dados e que é iterativamente particionado até que cada instância esteja contida em grupo formado única e exclusivamente por ela mesma. Já nos segundos ocorre o processo inverso, no qual, inicialmente, cada instância está contida única e exclusivamente em seu próprio grupo, e iterativamente os grupos são fundidos para formar um novo grupo, até que se tenha somente um grande grupo contendo todos os dados. O pseudocódigo de um algoritmo hierárquico aglomerativo se encontra detalhado em ~\ref{alg:aglomerative_clustering}.

O resultado de ambas as abordagens, é a estrutura de uma árvore binária chamada dendrograma, na qual as instâncias se encontram nas folhas da árvore e os ramos representam a união sucessiva das instâncias em um mesmo grupo. O tamanho do ramo que une dois grupos representa a distância, ou dissimilaridade, entre eles, de forma que ramos maiores representam uma maior distância entre os grupos. Dessa maneira, visualmente, pela altura dos ramos em um dendrograma, é possível se ter uma ideia do número de grupos que melhor representa a partição dos dados. Um exemplo de dendrograma pode ser visto na figura ~\ref{fig:dendrograma_exemplo}, o qual é resultado do agrupamento hierárquico com a \emph{linkagem} average, que será melhor detalhada a seguir, e métrica de dissimilaridade euclidiana para o famoso conjunto de dados Iris ~\parencite{Iris}.

\begin{figure}[h!]
	\includegraphics[width=\linewidth]{figuras/dendrograma_exemplo.png}
	\caption{Exemplo de um dendrograma resultante do agrupamento hierárquico com a \emph{linkagem} average e métrica de dissimilaridade euclidiana para o conjunto de dados Iris.}
	\label{fig:dendrograma_exemplo}
\end{figure}

\begin{algorithm}
	\caption{Agrupamento aglomerativo.}
	\label{alg:aglomerative_clustering}
	\begin{algorithmic}[1]
		
		\STATE{\textbf{Entradas:  matriz X de dados com n instâncias}}
		\STATE{\textbf{Saída:  partição $\mathcal{C}$ dos dados }}
	%	\renewcommand{\REQUIRE}{\textbf{Entradas: }}
	%	\REQUIRE \text{Matriz quadrada de similaridade d de dimensão n}
%		\STATE
		\FOR{$i \leftarrow 1$ \TO $n$} 
			\STATE{$\mathcal{C}_i \leftarrow {i}$} //inicializa cada grupo;
		\ENDFOR
		\STATE $\mathcal{S} \leftarrow \{1,..,n\}$; // inicializa o conjunto de grupos a serem aglomerados;
		\WHILE{$|\mathcal{S}| > 1$}
			\STATE{$(j,k) \leftarrow arg \min_{j,k \in \mathcal{S}} d_{j,k} $}; //Escolhe os dois grupos mais similares
			\STATE{$\mathcal{C}_l \leftarrow \mathcal{C}_j \cup \mathcal{C}_k$}; // Cria o novo grupo
			\IF{$\mathcal{C}_l \neq \{1,...,n\}$}
				\STATE{$\mathcal{S} \leftarrow \mathcal{S} \cup \{l\}$}; // Marca o novo grupo como disponível para algomeração
			\ENDIF
			\renewcommand{\algorithmicforall}{\textbf{foreach}}
			\FORALL{$i \in \mathcal{S}$}
				\STATE{d(i,l)} ; // Atualiza a matriz de similaridade
			\ENDFOR			
		\ENDWHILE
		
	\end{algorithmic}
\end{algorithm}

Existem três variações principais na maneira de se definir a dissimilaridade entre os grupos, e cada uma delas pode acarretar em resultados bastante diferentes. A seguir seguem breves comentários de cada uma delas. 

Na abordagem \emph{Single link}, a distância entre dois grupos é definida como a menor distância entre duas instâncias de cada grupo. Por sua vez, na \emph{Complete link}, a distância entre dois grupos é a maior distância entre duas instâncias de cada grupo. Assim, se o diâmetro de um grupo é definida como a maior distância entre instâncias contidas no grupo, então a abordagem \emph{complete link} gera partições com grupos bastante compactos mas pouco dispersos, enquanto a estratégia \emph{single link} cria partições com agrupamentos pouco compactos mas muito dispersos entre si.

 Em uma tentativa de se alcançar as virtudes de ambos os métodos, foi proposta a \emph{average link}, no qual a distância entre dois grupos é definida como a distância média das distâncias entre todos os pares de instâncias que formam o grupo. Na prática, a abordagem \emph{average link} é a mais utilizada por obter partições que representam um compromisso entre compactação e dispersão dos dados.


\subsection{Algoritmos baseados em densidades}

Os algoritmos particionais possuem a grande desvantagem de ser necessário o conhecimento, \emph{a priori}, do número $k$ de grupos que os dados se encontram naturalmente particionados. Dessa maneira, faz-se necessário rodadas iterativas com diferentes valores de $k$ para se alcançar a partição desejada, o que em muitas aplicações, principalmente para aquelas com muitos dados, é proibitivo. Além disso, os algoritmos particionais têm performance degradada quando o formato dos grupos não é esférico. Por sua vez, os algoritmos hierárquicos possuem a deficiência da definição do ponto de poda do dendrograma para a obtenção dos grupos.

Em uma tentativa de sanar todos esses problemas foram criados os algoritmos baseados em densidades, nos quais regiões com alta densidade de instâncias se encontram os grupos e em regiões de baixa densidade, se encontram as regiões de separação destes. Essa visão é bastante intuitiva e pode ser generalizada para qualquer dimensão, onde a vizinhança de cada ponto do grupo, definida por certo raio, deve conter um número mínimo de pontos para este ser considerado pertencente ao grupo. O DBSCAN, \emph{Density-based spatial clustering of applications with noise}, é o precursor desta classe de algoritmos e é detalhado a seguir.

\subsubsection{DBSCAN} \label{sec:dbscan}

Proposto em ~\parencite{DBSCAN}, o DBSCAN é um algoritmo que consegue descobrir grupos de formas arbitrárias além de não ser necessário informar \emph{a priori} o número de grupos. Os autores do algoritmo advogam que ele é menos sensível à ruídos e mais eficiente que os algoritmos hierárquicos e particionais. Todos esses ganhos foram conquistados com a introdução de dois parâmetros que não são necessários nos demais algoritmos: \textbf{Eps}, que é o raio máximo entre duas instâncias para elas serem consideradas da mesma vizinhança e \textbf{MinPts}, o número mínimo de pontos que devem estar na vizinhança para um ponto ser considerado pertencente ao grupo. Ambos os parâmetros influenciam fortemente nos resultados do agrupamento e estimativa deles é baseada no conhecimento do problema e dos dados do agrupamento, ou então em análises gráficas como pode ser visto em ~\parencite{DBSCAN}.

A abordagem mais comum é fixar o valor de \textbf{MinPts} em aproximadamente $4$, e a partir daí traçar o gráfico \emph{k-dist} das instâncias do conjunto de dados com $k=MinPts$. O gráfico \emph{k-dist} é obtido pela ordenação das distância entre cada instância do conjunto de dados a ser agrupado e o seu k-ésimo vizinho mais próximo. A região do gráfico na qual ocorre um crescimento abrupto da curva, contém uma boa sugestão para \textbf{Eps}. O exemplo a seguir demonstra este procedimento.

Na figura ~\ref{fig:inital} se encontra um conjunto de dados gerado sinteticamente composto por 2500 instâncias, onde se nota a presença de dois grupos distintos. A distância euclidiana entre cada instância e as demais foi calculada, e a distância da terceira instância mais próxima, ou a distância do terceiro vizinho mais próximo, de cada instância foram armazenadas e em seguida ordenadas em ordem crescente. Ao gráfico obtido por estes valores ordenados dá-se o nome de \emph{k-dist}, que pode ser visualizado na figura ~\ref{fig:k_dist}. Nesta figura fica evidente que para uma distância de aproximadamente $0.086$, ocorre um ponto de inflexão do gráfico \emph{k-dist}, o que sugere que este é um bom valor para \emph{Eps}. Assim, com os valores de \emph{MinPts}=$3$ e \emph{Eps}=$0.086$, os dados foram agrupados pelo DBSCAN, e a partição resultante pode ser vista na figura ~\ref{fig:result}. É notável que algumas poucas instâncias, destacadas em verde, foram consideradas anômalas, ou \emph{outliers}, onde pôde-se verificar a robustez do DBSCAN frente a presença de dados ruidosos.

\begin{figure}[h!]
	\includegraphics[width=\linewidth]{figuras/initial.png}
	\caption{Conjunto de dados sintético formado por dois grupos.}
	\label{fig:inital}
\end{figure}

\begin{figure}[h!]
	\includegraphics[width=\linewidth]{figuras/k_dist.png}
	\caption{Gráfico \emph{k-dist} do conjunto de dados da figura ~\ref{fig:inital}, onde o crescimento abrupto ocorre para a distância 0.086 aproximadamente.}
	\label{fig:k_dist}
\end{figure}

\begin{figure}[h!]
	\includegraphics[width=\linewidth]{figuras/result.png}
	\caption{Partição dos dados da figura ~\ref{fig:inital} obtida pelo DBSCAN com \emph{MinPts}=$3$ e \emph{Eps}=$0.086$. O pontos verdes foram considerados \emph{outliers} pelo algoritmo.}
	\label{fig:result}
\end{figure}

Além do método demonstrado, existem outras tentativas para se estimar os melhores valores dos dois parâmetros, de forma automática, como em ~\parencite{DE_DBSCAN}, no qual os valores são obtidos por meio de algoritmos de evolução diferencial. Nesta abordagem, cada indivíduo da população recebe um valor de \textbf{Eps} e \textbf{MinPts}, e uma partição por meio do DBSCAN é obtida para cada indivíduo. A função \emph{fitness} do algoritmo de otimização é dada por índices de validação interna dos grupos, a serem detalhados na seção ~\ref{sec:indice_validacao_interna}, e o algoritmo continua até um critério de parada. Dessa maneira, espera-se alcançar valores de  \textbf{Eps} e \textbf{MinPts} que geram a partição com os melhores índices, que por sua vez refletem em uma melhor partição. Dessa maneira, pode-se afirmar que a abordagem não é muito diferente da estratégia iterativa de obtenção do valor de $k$ dos algoritmos particionais, nos quais os algoritmos têm que ser executados diversas vezes. Já em ~\parencite{OPTICS}, é proposto o algoritmo OPTICS, \emph{Ordering Points to Identify the Clustering Structure}, no qual se faz necessária somente a informação do valor de \textbf{MinPts}, sendo o \textbf{Eps} estimado pelo próprio algoritmo. 

Apesar de no exemplo apresentado a distância euclidiana ter sido utilizada, o DBSCAN pode receber como parâmetro uma matriz de dissimilaridade gerada por qualquer outra dissimilaridade, o que, no contexto de séries temporais, o torna bastante interessante. O custo computacional do DBSCAN é de $\mathcal{O}(n \log n)$.

\section{Índices de validação de agrupamentos}

Após a definição de uma métrica de dissimilaridade, da escolha de um algoritmo de agrupamento, e da realização do agrupamento em si, faz-se necessária a avaliação da partição obtida. Tal avaliação é realizada pelos índices de validação que podem ser divididos em dois grandes grupos, sendo o primeiro formado pelos índice de validação interna enquanto o segundo contém os índices de validação externa. 

Os índices de validação interna são utilizados quando não se dispões de informações externas dos dados, e dessa maneira, ele é obtido apenas pelo resultado do agrupamento e das dissimilaridades entre as instâncias. Por sua vez, os índices de validação externos, utilizam alguma outra informação, além do resultado do agrupamento, no cálculo do índice. Tal informação externa é o resultado esperado do agrupamento, ou seja, os rótulos de cada instância. Uma discussão mais profunda de cada um dos tipos de índices de validação de agrupamentos é realizada a seguir.

\subsection{Índices de validação interna} \label{sec:indice_validacao_interna}

Normalmente, nas tarefas de agrupamento de dados, os índices de validação interna são os mais utilizados, pelo fato de não se possuir \emph{a priori} a rotulação das instâncias. Os índices de validação internos podem ser úteis também para a definição de um valor $k$ que represente o número de grupos, no qual os dados se encontram naturalmente distribuídos. O número pode ser aproximado ao se obter diversas partições para valores iterativos de $k$, e a partição que acarretar no melhor valor do índice de validação interno será a que possui o valor de $k$ recomendado.

Basicamente duas características são mensuradas pelos índices de validação interna: a \textbf{compactação} dos grupos, ou seja, o quão próximas, ou similares, as instâncias contidas em um mesmo grupo são entre si, e a \textbf{separação} dos grupos, ou seja, o quão distantes, ou dissimilares as instâncias de diferentes grupos são entre si. Cada uma das estratégias descritas a seguir calcula e pondera de formas distintas essas duas características para a obtenção do índice de validação dos resultados do agrupamento.

\subsubsection{\emph{Silhouette}}

O silhouette é um índice de validação proposto por ~\parencite{Rousseeuw:1987:SGA:38768.38772} que tenta avaliar o quão compacta, ou quão similar, cada instância se encontra dentro de seu próprio grupo, e, simultaneamente, o quão distante, ou dissimilar, ela é da instância mais próxima, ou mais similar, à ela e que não pertence ao seu mesmo grupo.

Para cada instância $x_i$ do conjunto de dados $X$ contendo $n$ instâncias, agrupada no grupo $\mathcal{C}_k$, devem ser calculados dois valores, sendo o primeiro a média da métrica de dissimilaridade entre $x_i$ e as demais instâncias contidas em $\mathcal{C}_k$, denominado $a$, e o segundo como o valor mínimo dentre todas as dissimilaridades entre $x_i$ e as demais instâncias não contidas em $\mathcal{C}_k$, denominado $b$. Ou seja, o valor de silhouette $S_i$ do ponto $x_i$ é definido como:
\begin{equation}
S_i = \frac{(b-a)}{\max(a,b)}, \mathrm{onde}
\end{equation}

\begin{equation}
a = \frac{1}{n} \sum_{j}^{|\mathcal{C}_k|} dist(x_i,x_j); x_j, x_i \in \mathcal{C}_k,
\end{equation}

\begin{equation}
b =\min(dist(x_i,x_j)); x_i \in \mathcal{C}_k,  x_j \notin \mathcal{C}_k,
\end{equation}
sendo $dist(x_i,x_j)$ o valor da métrica de dissimilaridade entre as instâncias $x_i$ e $x_j$.

Os valores de $S_i$ variam de $-1$ a $1$, sendo o valor de $-1$ geralmente quando a atribuição da instância $x_i$ ao seu grupo é equivocada, e $1$ quando é acertada. Valores próximos de $0$ podem indicar sobreposição de grupos. O valor de silhouette da partição é definido como a média dos valores de cada instância do conjunto de dados:

\begin{equation}
S = \frac{1}{n} \sum_{i}^{n} S_i.
\end{equation}

Valores próximos de $1$ indicam uma boa partição, ao passo que valores próximos de $-1$ indicam uma partição ruim.

\subsubsection{Índice de Dunn}

Proposto por ~\parencite{Dunn}, o índice de Dunn é definido como a razão entre a mínima distância intragrupo e a máxima intergrupo. Sendo assim, dada uma partição com $K$ grupos, sendo $\mathcal{C}_k$ o k-ésimo grupo dessa partição, então $dist_{kk'}$ é definida como a menor distância entre cada uma das  $n_k$ e $n_k'$ instâncias dos grupos $\mathcal{C}_k$ e $\mathcal{C}_k'$ respectivamente:

\begin{equation}
dist_{kk'} = \min_{i=1,...,n_k\\j=1,...,n_{k'}} \text{			} \bigg(dist(x_{i}^{k},x_{j}^{k'})\bigg),
\end{equation}
e $d_{min}$ é a menor destas distâncias:

\begin{equation}
d_{min} = \min_{k \neq k'} dist_{kk'}.
\end{equation}

O diâmetro do grupo $\mathcal{C}_k$, é definido como a maior distância entre os distintos $n_k$ pontos contidos pelo grupo $\mathcal{C}_k$:

\begin{equation}
D_k = \max_{x_i,x_j \in \mathcal{C}_k, x_i \neq x_j} dist(x_i,x_j),
\end{equation}
e finalmente, o maior diâmetro da partição é definido como $d_{max}$:

\begin{equation}
d_{max} = \max_{1\leq k \leq K} D_k.
\end{equation}

O índice \emph{Dunn} é definido pela razão entre $d_{min}$ e $d_{max}$:

\begin{equation}
Dunn = \frac{d_{min}}{d_{max}}
\end{equation}

O índice \emph{Dunn} varia de $[0,\infty)$ e valores menores implicam em partições melhores.

\subsubsection{\emph{Davies-Bouldin Index}}

Proposto em ~\parencite{DBI}, o índice Davies-Bouldin é obtido da seguinte maneira: dada um partição com $K$ grupos e sendo $G_k$ o ponto definido como o ponto mais representativo do k-ésimo grupo, onde por mais representativo entende-se o centróide no espaço euclidiano e o medóide para outras métricas de dissimilaridade, então a  distância média $\delta_k$ entre cada instância e seu ponto mais representativo é definida como:

\begin{equation}
\delta_k = \frac{1}{n_k}\sum_{i}^{n_k} dist(x_i,G_k),
\end{equation}
onde $n_k$ é número de instâncias contidas no k-ésimo grupo e $x_i$ é a i-ésima instância do grupo $k$. Essa métrica de compactação de grupo deve ser calculada para cada uma dos $k$ grupos da partição, e a dispersão entre os grupos é calculada entre cada um dos pares de grupos que formam a partição:

\begin{equation}
\Delta_{kk'} = dist(G_k,G_k').
\end{equation}

Finalmente, o índice Davies-Bouldin é definido como:

\begin{equation}
DBI = \frac{1}{K}\sum_{k=1}^{K} \max_{k \neq k'} \bigg( \frac{\delta_k + \delta_k'}{\Delta_{kk'}}  \bigg).
\end{equation}

O DBI é um índice de validação interna definido em $[0,\infty)$, e quanto maior o seu valor melhor é considerada a partição.


\subsubsection{Índice de Calinski-Harabasz}

Em \parencite{CH} é proposto um índice de validação para a identificação de grupos em um espaço euclidiano. Para tal ele define uma métrica que pondera as variâncias intragrupo WGSS (\emph{within-group sum of sqaure distance }), entendida como a variância da distância entre o centróide de um grupo e as suas instâncias, e a dispersão entre os grupos BGSS (\emph{between group sum of square distances}), entendida como a dispersão entre os centróides de cada grupo e o centróide de todo o conjunto de dados. Os valores obtidos são ponderados em função do número de grupos $k$. Então, o índice de Calinski-Harabasz é definido como

\begin{equation} \label{eq:CH}
CH = \frac{\frac{\text{BGSS}}{k-1}}{\frac{\text{WGSS}}{n-k}},
\end{equation}
sendo
\begin{equation}
BGSS=\sum_i{n_i d^2(c_i,c)}
\end{equation}
e
\begin{equation}
WGSS=\sum_i\sum_{x \in C_i} d^2(x,c_i),
\end{equation}
onde $c_i$ e $c$ são os centróides do i-ésimo grupo e o centróide de todo o conjunto de dados respectivamente, e $d(x_1,x_2)$ é a distância euclidiana entre os pontos $x_1$ e $x_2$.

Apesar de ter sido definido para o espaço euclidiano, os próprios autores em \parencite{CH} fazem menção à possibilidade de extensão do índice a espaços não-euclidianos. Para tal, basta utilizar outros conceitos como o medóide, ao invés do centróide, e a própria distância não-euclidiana entre os medóides e as instâncias ao invés da distância euclidiana entre os centróides e as instâncias.

Uma partição ideal tem alta dispersão entre os grupos (BGSS), e baixa variância intragrupo (WGSS), de forma que pela equação ~\ref{eq:CH} fica evidente que valores grandes de CH são desejados. O índice em questão é definido no intervalo $]0,\infty)$.

%\subsubsection{Índice I}

%O índice I, é definido em \parencite{I} como

%\begin{equation} \label{eq:I}
%\bigg( \frac{1}{k} . \frac{\sum_{x \in X} d(x,c)}{\sum_i\sum_{x \in C_i} d(x,c_i)} . max_{i,j} d(c_i,c_j)\bigg)^p,
%\end{equation}
%onde $k$ é o número de grupos, $C_i$ é o i-ésimo grupo e $c_i$ é o ponto mais representativo do i-ésimo grupo, definido como o centróide para espaços euclidianos e medóide para espaços não-euclidianos. Por sua vez, $p$ é um valor de correção e sugerido igual a $2$.

%Como pode ser visto pela equação ~\ref{eq:I}, o índice é composto por três termos:

%\begin{itemize}
%	\item $\frac{1}{k}$ cuja função é reduzir o índice com o incremento de $k$,
%	\item $\frac{E_1}{E_k}$ onde o termo $E_1$ é constante para determinado conjunto de dados e o termo $E_k$ diminui com o crescimento de $k$ indicando 
%\end{itemize}

%\hl{TODO}

%\subsubsection{Índice de Xie-Beni}
%\hl{TODO}

\subsection{Índices de validação externa}

Em tarefas reais de agrupamento, os índices de validação externa não têm utilidade direta, pois como já dito, em tarefas de agrupamento não se dispõe \emph{a priori} do rótulo das instâncias. Eles são muito úteis para se validar parâmetros e estratégias de agrupamento, ao se obter um valor numérico que faça com que seja possível comparar o resultado obtido  de cada abordagem com o resultado esperado.

Assim, variando-se os algoritmos, métricas, estratégias de pré-processamento e outros parâmetros ou decisões da tarefa de agrupamento de determinado conjunto de dados rotulado, pode-se comparar cada partição obtida com a partição real, ou verdadeira. Ou seja, caso se deseje realizar o agrupamento de dados não rotulados, e não se sabe a melhor estratégia para tal, se, primeiramente,  o agrupamento for feito em uma base de dados que seja, de alguma forma, semelhante à base alvo e a primeira esteja rotulada, então os índices de validação externa podem ser úteis na obtenção de uma estratégia mais adequada para a base não rotulada. %A seguir seguem alguns comentários sobre os principais índices de validação externa.

Existem diversos índices de validação externa tais quais o \emph{Adjusted Rand Index} (ARI) ~\parencite{Hubert1985}, \emph{Adjuste Mutual Info } (AMI) ~\parencite{AMI}, \emph{V-measure} ~\parencite{V_measure}, \emph{Variation of Information} (VI) ~\parencite{Meila}, entre outros. Neste trabalho, o índice VI foi especialmente utilizado, pelo fato de, diferentemente dos demais, ser uma métrica no espaço de partições. Dada a sua importância, ele se encontra detalhado na seção seguinte.


\subsubsection{\emph{Variation of Information} - VI} \label{sec:VI}

Proposto em ~\parencite{Meila}, e baseado em técnicas da teoria da informação, o VI mede a quantidade de informação ganha ou perdida ao se mudar da partição $\mathcal{P}$ para a partição $\mathcal{P'}$. O VI possui a propriedade de ser uma métrica de distância no espaço de partições, ou seja, obedece aos quatro axiomas apresentados no início da seção ~\ref{sec:metricas}. Essa propriedade faz com que seja possível, por exemplo, realizar o agrupamento de partições. 

Assim, dada uma partição $\mathcal{P}$ formada por $n$ instâncias divididas em $k$ grupos, onde o k-ésimo grupo contém $n_k$ instâncias, podemos definir uma variável aleatória como a probabilidade de se escolher uma instância do k-ésimo grupo e a entropia da partição como:

\begin{equation}
P(k) = \frac{n_k}{n}\text{ e,}
\end{equation}

\begin{equation}
H(\mathcal{P}) = -\sum_{i}^{k}P(k) \log P(k).
\end{equation}

Sendo $P(k)$ e $P'(k')$ as variáveis aleatórias associadas às partições $\mathcal{P}$ e $\mathcal{P'}$ respectivamente. Então, pode-se definir a probabilidade de uma instância pertencer ao grupo $\mathcal{C}_k$ da partição $\mathcal{P}$ e ao grupo $\mathcal{C'}_{k'}$ da partição $\mathcal{P'}$ como:

\begin{equation}
P(k,k') = \frac{|\mathcal{C}_k\cap \mathcal{C'}_{k'}|}{n},
\end{equation}
e a informação mútua entre as partições $\mathcal{P}$ e $\mathcal{P'}$ como:

\begin{equation}
I(\mathcal{P},\mathcal{P'}) = \sum_{i}^{k} \sum_{i}^{k'} P(k,k') \log \frac{P(k,k')}{P(k)P'(k')}.
\end{equation}

Finalmente, a variação da informação entre as partições $\mathcal{P}$ e $\mathcal{P'}$ pode ser definida como:

\begin{equation}
VI(\mathcal{P},\mathcal{P'}) = H(\mathcal{P}) + H(\mathcal{P'}) - 2I(\mathcal{P},\mathcal{P'}).
\end{equation}
Pelo fato do índice em questão ser uma métrica no espaço de partições, se soubermos a partição esperada, ou verdadeira, de um conjunto de dados, denominada $\mathcal{P}_v$, e possuirmos duas partições de teste $\mathcal{P}_{t1}$ e $\mathcal{P}_{t2}$, a partição de teste que possuir menor VI em relação à partição verdadeira, é considerada uma melhor partição dos dados do que a outra partição de teste. Ou seja,:

\begin{equation}
\begin{cases}
\text{se } VI(\mathcal{P}_v,\mathcal{P}_{t1})  < VI(\mathcal{P}_v,\mathcal{P}_{t2})\text{, então $\mathcal{P}_{t1}$ é melhor, ou mais verossímel, que $\mathcal{P}_{t2}$} ,\\
\text{se } VI(\mathcal{P}_v,\mathcal{P}_{t1})  > VI(\mathcal{P}_v,\mathcal{P}_{t2})\text{, então $\mathcal{P}_{t2}$ é melhor, ou mais verossímel, que $\mathcal{P}_{t1}$},\\
\text{se } VI(\mathcal{P}_v,\mathcal{P}_{t1})  = VI(\mathcal{P}_v,\mathcal{P}_{t2})\text{, então $\mathcal{P}_{t2}$ e $\mathcal{P}_{t1}$ são aproximações equivalentes de $\mathcal{P}_v$}.
\end{cases}
\end{equation}

Tais propriedades fazem com que o índice de variação da informação seja extremamente interessante para se construir um ranking entre diversas variações de estratégias de agrupamento para uma mesma base de dados rotulada. Dessa maneira, o VI será utilizado nos testes teóricos a serem realizados no Capítulo ~\ref{cap:testes_teoricos}.

O valores de VI variam de 0, quando $2I(\mathcal{P},\mathcal{P'})=H(\mathcal{P}) + H(\mathcal{P'})$, o que ocorre somente quando $\mathcal{P}=\mathcal{P'}$, até $H(\mathcal{P}) + H(\mathcal{P'})$, o que ocorre quando $I(\mathcal{P},\mathcal{P'}) =0$. Dessa maneira, é possível que o índice VI seja normalizado, dividindo o valor obtido pelo seu valor máximo, que como mencionado é a soma da entropia de cada partição, confinando-o, assim, no intervalo $[0,1]$. No entanto, quando este procedimento é realizado o índice perde a sua característica de métrica no espaço de partições, como provado em ~\parencite{Meila}.

%\subsubsection{\emph{V-Measure Index}}
%\hl{TODO}

%\subsubsection{\emph{IGT - Index of ground truth}}
%\hl{TODO}

%\subsubsection{\emph{Adjusted Rand Index}}
%\hl{TODO}

\section{Conclusões do Capítulo}

Neste capítulo foram apresentados as diversas escolhas que podem ser feitas em tarefas de agrupamento de séries temporais. Inicialmente, foram apresentadas as possíveis operações a serem realizadas em uma fase de pré-processamento como a normalização dos dados, redução de dimensionalidade ou transformadas para se trabalhar no domínio da frequência. Em seguida, foram apresentadas diversas métricas de dissimilaridades e em seguida as possíveis invariâncias desejadas em tarefas de mineração de dados de séries temporais foram enumeradas e exemplificadas. Por fim, os algoritmos mais comuns utilizados no agrupamento de séries temporais, e as estratégias de avaliação de partições, que são realizados por meio de índices internos e externos, foram discutidas.

Após esta extensa revisão acerca do tema de mineração de dados de séries temporais, e, especificamente no contexto de agrupamento, tornou-se evidente a complexidade do assunto, dado, principalmente, às diversas possibilidades de se realizar a tarefa de agrupamento que surgem a partir das diversas combinações possíveis de pré-processamento, métrica de dissimilaridade, algoritmo de agrupamento e método de avaliação do resultado. A cada combinação dessas escolhas dá-se o o nome de estratégia de agrupamento e, uma vez que tal diversidade de estratégias torna difícil a proposição de uma estratégia de agrupamento genérica, faz-se necessário um estudo detalhado e restrito a um número reduzido de escolhas mais promissoras para avaliar a performance de tais escolhas e a consequente proposição de uma estratégia.

Tal estudo foi realizado e os resultados, comentários e conclusões se encontram no Capítulo seguinte (~\ref{cap:testes_teoricos}).

\chapter{Testes em bases de dados \emph{benchmark}} \label{cap:testes_teoricos}

Neste capítulo são apresentados resultados de experimentos com conjuntos de dados de séries temporais que se encontram rotulados e estão disponíveis em ~\parencite{UCRArchive}. Cada uma das $85$ bases do conjunto de dados pode ser visualizada no Capítulo ~\ref{cap:curvas_UCR}, onde estão detalhados os medóides e centróides de cada grupo que forma a partição real dos dados, obtidos por meio da distância euclidiana. Pela análise de tais figuras, pode-se perceber que existem determinados conjuntos de dados que a distinção entre as instâncias de grupos diferentes é extremamente árdua (vide Figura ~\ref{fig:Coffee})  ao passo que para outros conjuntos esta é uma tarefa relativamente simples (vide Figura ~\ref{fig:BirdChicken}). 

A heterogeneidade das bases que compõem o UCR \emph{Time Series} é notável no que diz respeito ao domínio do problema já que existem curvas que são: sinais biológicos, geradas em problemas de visão computacional, sinais de equipamentos elétricos, sinais multivariados, detre outros. Ademais, a diversidade no que diz respeito ao número de instâncias, grupos e tamanho das séries temporais também é bastante significativa. 

Se considerarmos que no Capítulo ~\ref{cap:agrupamento_series_temporais}:

\begin{itemize}
	\item apresentou-se 5 tipos possíveis de pré-processamento: normalização Z, normalização max, normalização min-max, transformada discreta de Fourier e \emph{Piecewise aggregate aproximation}, desconsiderando-se as possíveis combinações entre eles, como, por exemplo, \emph{Piecewise aggregate aproximation} e normalização Z,
	\item discutiu-se 12 métricas de dissimilaridades puras (euclidiana, \emph{manhattan}, \emph{chebyshev}, DTW, LB-Keogh, EDR, Lorentzian, \emph{dice}, \emph{jaccard}, \emph{pearson}, \emph{mahalanobis}, \emph{cosine}) além das versões corrigidas pela correlação temporal (CORT) e complexidade (CID) de cada uma delas, existem 36 métricas de dissimilaridades  possíveis,
	\item analisou-se 7 algoritmos diferentes, sendo 3 hierárquicos (cada um com uma das \emph{linkagens }\emph{average}, \emph{complete} e \emph{single}), 3 particionais (\emph{k-means},\emph{k-means++} e \emph{k-medoids}) e um baseado em densidades (DBSCAN),
\end{itemize}
 então, dentre estas opções, existe um limite de $5*36*7=1260$ maneiras distintas de-se realizar o agrupamento de dados de séries temporais. Este número é, na verdade, menor, devido ao fato de algoritmo \emph{k-means} fazer uso somente da distância euclidiana e de algumas dissimilaridades como a DTW, EDR, dentre outras serem inaplicáveis após a transformação dos dados pela transformada de Fourier. 
 
 No entanto, este número demonstra bem o grande número de caminhos possíveis a serem seguidos e a dificuldade enfrentada por um projetista quando ele deve escolher um destes caminhos para realizar a tarefa de agrupamento de séries temporais. Testes iniciais foram feitos com as diversas possibilidades enumeradas, no entanto a análise dos resultados obtidos se tornou árdua, pois, dado o grande número de estratégias testadas, não ficou claro se alguma delas, ou algum grupo delas, se sobressaiu em relação às demais de forma significativa considerando-se todas as bases de dados. Por isto, o número de possibilidades de agrupamento das bases de dados do UCR \emph{Time Series} neste capítulo é significativamente reduzido, e somente as escolhas mais comuns na literatura serão estudadas a fundo e comparadas entre si. 
 
 Pelo fato de todas as curvas do conjunto de dados já se encontrarem normalizadas pela normalização Z, e de estudos afirmarem que, em tarefas de classificação ela é superior às demais ~\parencite{trillions}, ela é a única escolha de pré-processamento a ser utilizada em todos os experimentos deste Capítulo. No que diz respeito à escolha de métrica de dissimilaridade, a distância euclidiana foi escolhida por ser uma das métricas mais utilizadas em agrupamentos de curva de carga ~\parencite{Chicco}, além da DTW por esta ser recorrente na literatura e ter apresentado performance significativamente superior à euclidiana nos testes realizados em ~\parencite{BatistaComparativo}. Além das distâncias puras, as respectivas versões que levam em conta a correlação temporal das séries (CORT) e complexidade relativa entre elas (CID) também serão utilizadas para se verificar eventuais ganhos com a aplicação destas.
 
Dentre os algoritmos estudados, serão testados os algoritmos hierárquicos com \emph{linkagens average} e \emph{complete}, além dos algoritmos particionais \emph{k-means} e \emph{k-medoids}. O \emph{k-means} foi utlizado por ser, assim como a distância euclidiana, onipresente em agrupamento de curvas de carga ~\parencite{Chicco}, e a sua variação, o \emph{k-medoids} para que seja possível a realização do agrupamento por meio da combinação de um algoritmo particional e uma métrica de dissimilaridade não-euclidiana como a DTW, por exemplo. Já os algoritmos hierárquicos foram utilizados por apresentarem performance superior aos algoritmos particionais nos experimentos apresentados em ~\parencite{k_shape}. Além da análise das métricas mais promissoras, já mencionadas, todas elas serão confrontadas com a estratégia derivada da combinação do algoritmo \emph{k-means} e distância euclidiana já que, como já dito, esta é a estratégia padrão no agrupamento de curvas de carga.

Os experimentos deste Capítulo foram conduzidos com a finalidade de responder as seguintes perguntas:

\begin{itemize}
	\item qual métrica de dissimilaridade é a mais apropriada para séries temporais no que diz respeito à qualidade da partição e custo computacional?
	\item qual é o algoritmo de agrupamento mais apropriado para a realização do agrupamento de séries temporais?
	\item qual índice de validação interno é capaz de nos indicar as melhores partições de dados no agrupamento de séries temporais?
\end{itemize}

São realizados dois experimentos, sendo o primeiro destinado a responder às duas primeiras perguntas, enquanto o segundo é destinado a responder à terceira pergunta. Em todos experimentos são realizados os testes estatísticos de significância não-paramétricos de Friedman, Nemenyi e Bonferroni-Dunn para a comparação de diferentes estratégias de agrupamento aplicadas a diversas bases de dados, como sugerido em ~\parencite{Demsar} e detalhado ne seção ~\ref{sec:testes_estatisticos}.

Todos experimentos apresentados nesta dissertação de mestrado foram realizados em uma máquina com CPU  Intel(R) Core(TM) i5-2400 3.10GHz e 8GB de RAM. Os códigos foram feitos em python onde foram utilizados, principalmente, os módulos numpy ~\parencite{numpy}, scipy ~\parencite{scipy}, Orange3 ~\parencite{JMLR:demsar13a} e scikit-learn ~\parencite{scikit-learn} e os códigos utilizados podem ser encontrados em ~\parencite{codigo_fonte}.

\section{Testes estatísticos de significância} \label{sec:testes_estatisticos}

No teste de Friedman ~\parencite{Friedman1940Comparison}, dado $M$ conjuntos de dados e $L$ algoritmos, ou estratégias, com seus rankings médios $R_t$, onde $t=1,...,L$, assume-se a hipótese nula $H_0$ de que os algoritmos, ou estratégias, são equivalentes. Neste caso, a estatística
\begin{equation}
F_F = \frac{(M-1)\chi_{F}^{2}}{M(L-1)-\chi_{F}^{2}}
\end{equation} \label{eq:Friedman_F}
é distribuída de acordo com a distribuição F com $L-1$ e $(L-1)(M-1)$ graus de liberdade, onde 

\begin{equation}
\chi_{F}^{2}=\frac{12M}{L(L+1)}\bigg(\sum_{t}^{} R_t^2 - \frac{L(L+1)^2}{4} \bigg).
\end{equation}

Caso a hipótese nula $H_0$ seja rejeitada pelo teste de Friedman, os testes \emph{post-hoc} de Nemenyi e Bonferroni-Dunn podem ser aplicados para apontar significância na comparação dos algoritmos, ou estratégias, um-contra-um e um-contra-todos respectivamente.

O teste de Nemenyi ~\parencite{Nem63} é utilizado quando se deseja realizar a comparação entre as estratégias para se verificar quais são equivalentes entre si. A performance entre duas estratégias é significativamente diferente se os rankings médios delas diferirem de, no mínimo,

\begin{equation} \label{eq:CD}
	CD = q_{\alpha}\sqrt{\frac{L(L+1)}{6M}}, 
\end{equation}
onde o valor crítico $q_{\alpha}$ é função do número de algoritmos e pode ser consultado em ~\parencite{Demsar}.

No entanto, caso se deseje comparar a performance de todas as estratégias em análise, com a performance de uma estratégia de controle, o método de Bonferrni-Dunn ~\parencite{bonferroni_dunn} deve ser utilizado, ao invés do teste de Nemenyi. O teste define quais as estratégias possuem performance significativamente equivalente à performance da estratégia de controle, e estas são aquelas estratégias cujo ranking médio se encontra dentro da faixa definida pelo valor de uma diferença crítica do teste de Bonferroni-Dunn centrada no ranking médio da estratégia de controle. Tal diferença crítica é calculada conforme equação ~\ref{eq:CD}, sendo seu valor diferenciado do valor da diferença crítica do teste de Nemenyi apenas pelo valor de $q_{\alpha}$, que também pode ser consultado em ~\parencite{Demsar}.

\section{Primeiro experimento - Métrica de dissimilaridade e algoritmo de agrupamento}

No contexto de agrupamento de séries temporais, particularmente no agrupamento de curvas de carga, a principal combinação de métrica de dissimilaridade e algoritmo de agrupamento utilizada é a distância euclidiana com o \emph{k-means} ~\parencite{Chicco}. O objetivo do experimento a ser realizado nesta seção é de verificar se existe alguma outra combinação que seja mais promissora que a combinação padrão. Para tal, utilizou-se $55$ bases de dados do UCR \emph{Time Series}, onde, pelo fato delas estarem rotuladas, o índice VI discutido na Seção ~\ref{sec:VI} foi utilizado para comparar a partição obtida por determinada escolha de algoritmo e métrica de dissimilaridade com a partição real dos dados.

Todas as séries do repositório UCR \emph{Time Series} já se encontram normalizadas pela normalização Z, e nenhuma estratégia de redução de dimensão dos dados, como a DFT ou o PAA, foi realizada. Além do índice VI, os tempos de processamento de cada estratégia também foram comparados para que fosse possível a comparação das estratégias em relação à este critério. Por tempo de processamento da combinação métrica de dissimilaridade e algoritmo de agrupamento, entende-se como o tempo resultante da soma do tempo de obtenção da matriz de dissimilaridade e do tempo da realização do agrupamento em si pelo algoritmo.

Dado a extensa possibilidade de combinações de algoritmos de agrupamento e métricas de dissimilaridade que podem ser obtidas a partir dos algoritmos e métricas discutidos no Capítulo ~\ref{cap:agrupamento_series_temporais}, apenas algumas delas foram analisadas nesta seção. Dessa forma, as métricas DTW e euclidiana, bem como suas variações que levam em conta a complexidade (CID) e correlação temporal (CORT), foram analisadas, já que segundo  estudos no contexto de classificação de séries temporais, elas se demonstraram superiores em relação às demais ~\parencite{BatistaComparativo}. Os experimentos são detalhados nas seções a seguir.

\subsection{Algoritmos de agrupamento} \label{sec:experimento_1}

Inicialmente, deseja-se verificar o efeito da escolha do algoritmo nos resultados do agrupamento. Assim, em um primeiro experimento, todas as bases serão agrupadas por diferentes algoritmos mas com a dissimilaridade fixa, e definida como a distância euclidiana, já que esta é uma das métricas de dissimilaridade mais comum. Foram escolhidos $4$ algoritmos particionais, com diferentes inicializações, e dois algoritmos hierárquicos, com diferentes \emph{linkagens}, e todas as combinações formadas podem ser vistas na Tabela ~\ref{tbl:1}.

%\begin{center}
	%\captionsetup{type=table}
	%\resizebox{\columnwidth}{!}{%
\begin{table}
	\centering	
	\begin{tabular}{llllll}
		\toprule
		{} &          Id &              Algoritmo & Dissimilaridade \\
		\midrule
		0 &  KDR\_\_eucl. &       k-medoids(random) &              eucl.\\
		1 &   HC\_\_eucl. &  hierarchical-complete &               eucl. \\
		2 &  KDP\_\_eucl. &              k-medoids(PAM) &               eucl. \\
		3 &  KMP\_\_eucl. &              k-means(++) &            eucl. \\
		4 &  KMR\_\_eucl. &                k-means(random) &               eucl. \\
		5 &   HA\_\_eucl. &   hierarchical-average &          eucl. \\
		\bottomrule
	\end{tabular}
	%}
%	\captionof{table}{Estratégias comparadas no experimento 1.}  \label{tbl:1}
	\caption{Estratégias comparadas no experimento 1.}  \label{tbl:1}
\end{table}
%\end{center}

Para as $55$ bases de dados e $6$ combinações (estratégias), o valor 
crítico da estatística de Friedman é de $2.24$. O valor médio dos 
\textit{rankings} forneceu uma estatística $F_F = 12.72$ rejeitando, portanto, $H_0$. 
Dessa forma, pode-se afirmar que a escolha do algoritmo de agrupamento 
tem efeito significativo no desempenho individual de uma dada medida de dissimilaridade. 

Finalmente, o teste \emph{post-hoc} de Bonferroni-Dunn foi aplicado, 
centrado na combinação padrão, \emph{k-means (random)} e distância Euclidiana. Os resultados 
deste teste podem ser vistos na Figura \ref{fig:VIP1}), e o traço presente nela,  é o valor da diferença crítica do teste de Bonferroni-Dunn, conforme detalhado em ~\ref{sec:testes_estatisticos}, centrado no ranking médio da estratégia padrão (\emph{k-means(random)} com distância Euclidiana ) . Os resultados mostram que o algoritmo hierárquico 
com linkagem \emph{average} possui desempenho significativamente superior 
a estratégia padrão, ao passo que as demais estratégias não possuem uma performance significativamente distinta dela.
\begin{figure}[h!]
	\includegraphics[width=\linewidth]{figuras/figure_combinacoes_1.png}
	\caption{Resultado do teste de Bonferroni-Dunn após o ranqueamento dos valores de VI. A estratégia de controle se ecnontra centrada no ranking médio da estratégia \emph{k-means(random)} com distância Euclidiana.}
	\label{fig:VIP1}
\end{figure}

\subsection{Métricas de dissimilaridade} \label{exp:2}

O objetivo deste segundo experimento é de comparar o desempenho
das diferentes métricas de dissimilaridade. Desde que o algoritmo exerce 
influência no desempenho individual da métrica, o procedimento usado 
neste experimento foi fixar o algoritmo de agrupamento 
e variar as métricas. O algoritmo fixado foi o agrupamento 
hierárquico com linkagem \emph{average}, por este ter apresentado o melhor 
resultado no experimento anterior. A combinação \emph{k-means} e 
distância Euclidiana foi também avaliada neste experimento, por se tratar 
da combinação padrão adotada na literatura. Todas as combinações 
deste experimento são mostradas na Tabela ~\ref{tbl:2}. 

%\begin{center}
	%\resizebox{\columnwidth}{!}{%
	\begin{table}		
	\centering
	\begin{tabular}{llllll}
		\toprule
		{} &              Id &             Algoritmo &      Similaridade \\
		\midrule
		0 &  HA\_\_cort-eucl. &  hierarchical-average &       cort-eucl. \\
		1 &         HA\_\_DTW &  hierarchical-average &        DTW \\
		2 &   HA\_\_cid-eucl. &  hierarchical-average &        cid-eucl. \\
		3 &      KMR\_\_eucl. &               k-means(random) &              eucl. \\
		4 &       HA\_\_eucl. &  hierarchical-average &         eucl. \\
		\bottomrule
	\end{tabular}
	%}
	%\captionof{table}{Estratégias comparadas no experimento 2.}  \label{tbl:2} 
	\caption{Estratégias comparadas no experimento 2.}  \label{tbl:2} 
		\end{table}
%\end{center}

Para $55$ bases de dados e $5$ combinações testadas, o valor crítico da estatística 
de Friedman é de $2.41$. A hipótese $H_0$ foi descartada, já que obteve-se $F_F=14.04$. Em seguida, 
deu-se prosseguimento ao teste \emph{post-hoc} de Bonferroni-Dunn, centrado na estratégia padrão (\emph{k-means(random)} com distância Euclidiana). O resultado, 
ilustrado na Figura ~\ref{fig:VIP2}, mostra que a combinação padrão é inferior à todas as demais, 
e além disso, que o uso da DTW e da CID acarretaram em melhoras significativas em relação à distância 
Euclidiana e CORT-Euclidiana. \\

\begin{figure}[h!]
	\includegraphics[width=\linewidth]{figuras/figure_combinacoes_2.png}
	\caption{Resultado teste Bonferroni-Dunn após o ranqueamento dos valores de VI, centrado no ranking médio da estratégia padrão (\emph{k-means(random)} com distância Euclidiana).}
	\label{fig:VIP2}
\end{figure}

\subsection{Comparação entre Métricas Promissoras} \label{exp:3}

Os experimentos anteriores mostraram que a adição de métricas 
robustas a distorções nos dados, tais como DTW e CID pode melhorar
significativamente a qualidade das partições obtidas. Com isso em mente, 
o objetivo deste experimento foi comparar as possíveis combinações 
envolvendo estas métricas promissoras nos quesitos de qualidade das 
partições (índice VI) e custo computacional (tempo de processamento). 
As combinações analisadas neste experimento se encontram na Tabela ~\ref{tbl:3}. 
Observe que, novamente, o algoritmo hierárquico com linkagem \emph{average} foi 
fixado.

%\begin{center}
	
	%	\resizebox{\columnwidth}{!}{%
	\begin{table}
		\centering		
	\begin{tabular}{llllll}
		\toprule
		{} &              Id &             Algoritmo & Similaridade \\
		\midrule
		0 &  HA\_\_cort-eucl. &  hierarchical-average & cort-eucl. \\
		1 &         HA\_\_DTW &  hierarchical-average &  DTW \\
		2 &   HA\_\_cid-eucl. &  hierarchical-average &  cid-ecul. \\
		3 &    HA\_\_cort-DTW &  hierarchical-average &  cort-DTW \\
		4 &     HA\_\_cid-DTW &  hierarchical-average &   cid-DTW \\
		\bottomrule
	\end{tabular}
	%	}
%	\captionof{table}{Estratégias do experimento 3.}  \label{tbl:3}
	\caption{Estratégias do experimento 3.}  \label{tbl:3}
		\end{table}
%\end{center}
%\setlength\parindent{20pt}

O teste de Friedman foi aplicado para testar a hipótese nula de que as estratégias 
são equivalentes. Para $55$ bases de dados e $5$ estratégias, o valor crítico 
é de $2.41$. Considerando o quesito tempo de processamento, 
obteve-se $F_F=543.03$ e, para o VI, $F_F=3.48$. Logo, 
para ambos os casos, a hipótese nula de que os rankings médios são iguais 
foi descartada. Deu-se então prosseguimento ao teste \emph{post-hoc} de Nemenyi, 
com o objetivo de verificar quais grupos de estratégias possuem desempenhos 
equivalentes. O diagrama resultante do teste de Nemenyi para o índice VI pode ser 
visto na Figura ~\ref{fig:VIP3}. O diagrama correlato, para o 
tempo de processamento, é mostrado na Figura ~\ref{fig:VIP4}. 

Pode-se concluir, com  base na Figura ~\ref{fig:VIP3}, que todas as combinações 
envolvendo as métricas DTW e CID são significamente melhores que a 
combinação envolvendo CORT e Euclidiana (HA\_cort-eucl.). 
A Figura ~\ref{fig:VIP4}, por sua vez, mostra que todas as combinações 
envolvendo DTW (HA-DTW, HA\_CID-DTW, HA\_cort-DTW) possuem maior custo computacional, 
sendo estatisticamente inferiores em relação as demais combinações. 
Uma análise conjunta, envolvendo os dois quesitos, apontam a combinação HA\_CID-eucl. 
como um bom equilíbrio entre a qualidade das partições obtidas e o menor tempo 
de processamento.

Com relação ao índice VI (Figura ~\ref{fig:VIP3}), observa-se ainda que, 
embora o teste não tenha apontado diferença estatística, os resultados 
em termos dos \textit{rankings} médios apontam que a correção CID (complexidade) 
produz efeito mais positivo do que a correção CORT (correlação temporal), 
independentemente da métrica convencional adotada: Euclidiana ou DTW. 
Além disso, todas as combinações envolvendo DTW se mostraram melhores 
que as outras em termos de \textit{rankings} médios, o que sugere a importância 
de um critério que seja invariante a deslocamentos locais no agrupamento 
de séries temporais.

\begin{figure}[h!]
	\includegraphics[width=\linewidth]{figuras/figure_combinacoes_4.png}	
	\caption{Resultado teste Nemenyi após o ranqueamento dos valores de VI.}
	\label{fig:VIP3}
\end{figure}

\begin{figure}[h!]	
	\includegraphics[width=\linewidth]{figuras/figure_combinacoes_3.png}
	\caption{Resultado teste Nemenyi após o ranqueamento dos valores de tempo de processamento.}
	\label{fig:VIP4}
\end{figure}


\section{Segundo experimento - Índice de validação interno} \label{sec:definicao_indice_interno}

Neste segundo experimento, tentaremos responder às seguintes perguntas:

\begin{itemize}
	\item Existe algum índice de validação interno que seja mais adequado para a definição do número de grupos um conjunto de dados de séries temporais?
	\item Uma vez definido o número de grupos, existe algum índice interno que seja mais indicado para avaliar as partições obtidas por diferentes estratégias de agrupamento?
\end{itemize}

Para responder à ambas perguntas os índices internos exibidos na Tabela ~\ref{tbl:indices_internos} foram comparados nos experimentos desta seção. Além desses índices da tabela, outros três índices farão parte da análise. Eles derivam da combinação de todos os índices da tabela, denominados índices puros, e são definidos como a moda (ou voto majoritário), média e mediana do número de grupos indicado por cada um dos $7$ índices.

\begin{center} 
	\begin{table}
		\caption{Principais índices internos de validação disponíveis na literatura}  \label{tbl:indices_internos}
		\resizebox{\columnwidth}{!}{%
			\begin{tabular}{lllll}
				\hline\noalign{\smallskip}
				Índice & Notação & Definição & Valor ótimo & Ref.\\
				\noalign{\smallskip}
				\hline
				\noalign{\smallskip}
				Calinski-Harabasz & CH & $\frac{\sum_i{n_i d^2(c_i,c)/(k-1)}}{\sum_i\sum_{x \in C_i} d^2(x,c_i)/(n-k) }$ & Max & \cite{CH} \\
				I index& I & $\bigg( \frac{1}{k} . \frac{\sum_{x \in X} d(x,c)}{\sum_i\sum_{x \in C_i} d(x,c_i)} . max_{i,j} d(c_i,c_j)\bigg)^p$ & Max & \cite{I}\\
				Dunn's index & Dunn & $min_i\bigg(min_j\bigg( \frac{min_{x \in C_i, y \in C_j} d(x,y))}{max_k(max_{x,y \in C_k} d(x,y))} \bigg)\bigg)$ & Max & \cite{Dunn} \\
				Silhouette & silhouette & $\frac{1}{k} . \sum_i^k \frac{1}{n_i} \sum_{x \in C_i} \frac{b(x)-a(x)}{max(b(x),a(x))}$ & Max & \cite{silhouette}\\
				Davies-Bouldin & DB & $\frac{1}{k} \sum_i max_{j,j \neq i} \bigg( \bigg( \frac{1}{n_i} \sum_{x \in C_i} d(x,c_i) + \frac{1}{n_j} \sum_{x \in C_j} d(x,c_j) \bigg) / d(c_i,c_j) \bigg)$ & Min & \cite{DBI}\\
				Xie-Beni & XB & $\frac{\sum_i \sum_{x \in C_i} d^2(x,c_i)}{n.min_{i,j \neq i} d^2(c_i,c_j)}$ & Min & \cite{XB}\\
				SD validity & SD & $Dis(k_{max}) Scat(k) + Dis(k)$ & Min & \cite{SD}\\
				{} & {} & $Scat(k) = \frac{1}{k} \frac{\sum_i^k ||\sigma(C_i)||}{||\sigma(X)||}$ & {} & {}\\
				{} & {} & $Dis(k) = \frac{max_{i,j}d(c_i,c_j)}{min_{i,j, i \neq j}d(c_i,c_j)} \sum_i^k \big( \sum_{j,j \neq i}^k d(c_i,c_j)  \big)^{-1}$ & {} & {}\\
				\hline
			\end{tabular}
		}
	\end{table}
\end{center}  

Na Tabela ~\ref{tbl:indices_internos} se encontram os índices internos a serem utilizados nesta seção,  onde $X$: é todo o conjunto de dados; $c$: é o centróide (quando a métrica de dissimilaridade escolhida é a distância euclidiana) ou medóide de $X$; $C_i$: $i$-ésimo grupo;  $n_i$: número de instâncias no $i$-ésimo grupo; $c_i$: centróide ou medóide do $i$-ésimo grupo; 
$d(x,y)$: medida de dissimilaridade entre $x$ e $y$; $n$: número de elementos em $X$; $k$: número de grupos em $X$, $||v|| = (v^T.v)^{\frac{1}{2}}$ e $\sigma(C_i)$: vetor de variância de todas as instâncias contidas em $C_i$; 
$p$: valor inteiro, recomendado $p=2$ ~\parencite{I}, $b(i)$: a menor dissimilaridade média da instância $i$ e cada instância contida em grupos diferentes do grupo de $i$ ;  $a(i)$ é a média de dissimilaridade da instância $i$ para todas as outras instâncias contidas no mesmo grupo de $i$.

\subsection{Índice para a definição do número de grupos} \label{sec:indice_iterativo}

Em tarefas de agrupamento, quando não se sabe \emph{a priori} o número de grupos nos quais os dados se encontram naturalmente divididos, são necessárias estratégias para se definir este número, já que ele é o principal parâmetro de entrada dos algoritmos particionais como o \emph{k-means} e o \emph{k-medoids}. Mesmo para algoritmos hierárquicos, esta é uma informação importante, pois a interpretação dos dendrogramas muitas vezes não é trivial, já que é difícil a definição do ponto de corte da árvore que define o particionamento dos dados. 

Uma estratégia comum é realizar diversos agrupamentos para diferentes valores de $k$ (número de grupos), dentro de determinado intervalo, no qual se espera que o valor ótimo $k_{opt}$, ou real, de grupos se encontra. Após a avaliação de cada partição resultante por um índice de validação interno, aquela que possuir um valor mais próximo do valor ótimo do índice, será a partição que apresenta o número ótimo ou mais indicado de grupos dentro do intervalo testado.

Dessa maneira, surge a pergunta: qual, ou quais, dos índices internos detalhados na literatura são mais interessante para a tarefa de se definir o número de agrupamentos de um conjunto de dados de séries temporais? Para responder à essa pergunta, iremos utilizar, novamente, o conjunto de dados UCR \emph{Time series} ~\parencite{UCRArchive}. Como cada conjunto de dados se encontra rotulado é sabido \emph{a priori} o número de grupos que os dados se encontram divididos em cada conjunto de dados. Assim, dado o valor esperado de determinado conjunto de dados, denominado $k_{opt}$, foi formado um intervalo $k = k_{opt}-6,k_{opt}-5,...,k_{opt},...,k_{opt}+6$, onde caso o primeiro valor seja menor que $2$, o intervalo é definido, necessariamente, a partir deste valor, já que em tarefas de agrupamento não faz sentido um valor de $k$ menor do que $2$.

Após a escolha de um algoritmo e uma métrica de dissimilaridade como, por exemplo, um algoritmo de agrupamento hierárquico com linkagem \emph{average} e a métrica CID-euclidiana, para cada valor de $k$ no intervalo criado, uma partição é obtida e avaliada com cada um dos índices da Tabela ~\ref{tbl:indices_internos}. A título de exemplo, tal procedimento, com o algoritmo e métrica de dissimilaridade acima mencionado, foi realizado no conjunto de dados \emph{Herring} (vide Figura ~\ref{fig:Herring}), que possui $k_{opt}=2$, e os resultados podem ser vistos na Tabela ~\ref{tbl:indices_iterativos}. Os índices DBI, XB e silhouette tiveram êxito na sugestão do valor de $k_{opt}$, ao passo que todos os demais falharam em tal tarefa. A partir da combinação da sugestão de cada um dos $7$ índices testados podem-se derivar os seguintes índices: moda dos valores sugeridos, cujo valor é 2; média dos valores sugeridos, que possui valor $3.57$ que é truncado para $4$ já que este é o valor inteiro mais próximo; e finalmente mediana dos valores sugeridos para $k_{opt}$, cujo valor é 3. Logo, para o teste em questão, dentre os índices originados pela combinação dos índices puros, somente a moda teve sucesso na indicação do valor de $k_{opt}$.

\begin{table} [!h]
	\centering
	\caption{Resultados dos experimentos iterativos para o conjunto de dados \emph{Herring}, cujo número de grupos é igual a 2. Os valores em negrito representam os melhores valores obtidos em cada índice.}  \label{tbl:indices_iterativos}
%	\resizebox{\columnwidth}{!}{%
		\begin{tabular}{lrrrrrrrr}
			\toprule
			{} &    k &      CH &    DBI &   Dunn &       I &     SD &     XB &  silhouette \\
			\midrule
   \rowcolor{lightgray} &  2.0 &   2.025 &  \bfseries{0.486} &  0.369 &   0.031 &  0.517 &  \bfseries{0.478} &       \bfseries{0.409} \\
   &  3.0 &  20.876 &  0.492 &  0.353 &  15.024 &  \bfseries{0.414} &  0.489 &       0.259 \\
   &  4.0 &  \bfseries{21.448} &  0.879 &  0.262 &  \bfseries{20.746} &  0.569 &  0.614 &       0.206 \\
   &  5.0 &  16.315 &  0.826 &  0.267 &  14.124 &  0.500 &  0.598 &       0.183 \\
   &  6.0 &  14.646 &  1.184 &  0.330 &  12.780 &  0.806 &  0.982 &       0.195 \\
   &  7.0 &  12.879 &  1.073 &  0.334 &  10.826 &  0.759 &  0.961 &       0.186 \\
   &  8.0 &  10.987 &  1.208 &  \bfseries{0.378} &   8.506 &  0.908 &  1.092 &       0.184 \\
			\bottomrule
		\end{tabular}
	%}
\end{table} 

Os valores em negrito na Tabela ~\ref{tbl:indices_iterativos} representam os melhores valores obtidos em cada índice, e quando este valor coincide com o valor de $k_{opt}$, então um acerto é computado para o índice naquela base. Em seguida, tal procedimento foi repetido para 70 bases de dados, com o objetivo de se construir uma tabela binária, representado os acertos e erros de cada índice em cada base. Ao final da tabela, o número de acertos de cada índice foi somado e dividido pelo número de bases, para se ter uma ideia da porcentagem de vezes que cada índice obteve sucesso. Devido às suas grandes dimensões (70x13), somente parte da tabela binária de tal experimento pode ser vista na Tabela ~\ref{tbl:indices_hits}, já que diversas linhas (datasets) e algumas colunas (índices) da tabela não se encontram presentes.

	\begin{table}[!h]
		\centering
		\caption{Resultados dos experimentos iterativos para cada conjunto de dados, onde o asterisco corresponde à um sucesso na sugestão do índice para a base de dados. Foi utilizado o algoritmo hierárquico com linkagem \emph{average} e métrica de dissimilaridade cid-euclidiana.}  \label{tbl:indices_hits} 
		%\resizebox{\columnwidth}{!}{
		%{\footnotesize
			\begin{tabular}{llrrrrrr}
				\toprule
				{} &                         Conjunto de dados &     CH &    DBI &   Dunn &      I &     SD &     XB \\
				\midrule
				0  &                         50words &    &    &    &    &    &   \\
				1  &                           Adiac &    &    &    &    &    &    \\
				2  &                       ArrowHead &  * &    &    &  * &    &  \\				
				26 &                         Herring &    &  * &    &    &    &  * \\			
				69 &                            yoga &  * &    &    &    &    &    \\
				 \rowcolor{lightgray} 70 &                           Total &  0.429 &  0.214 &  0.143 &  0.257 &  0.171 &  0.343\\
				 {} &                         {}  &     CH &    DBI &   Dunn &      I &     SD &     XB\\
			\bottomrule
		\end{tabular}
		%}
	\end{table}
	
	%	\begin{longtable}
	%	\caption{Resultados dos experimentos iterativos para cada conjunto de dados, onde o asterisco corresponde à um sucesso na sugestão do índice para a base de dados. Foi utilizado o algoritmo hierárquico com linkagem \emph{average} e métrica de dissimilaridade cid-euclidiana.}
	%\resizebox{\columnwidth}{!}{
	%{\footnotesize
	%	\begin{longtable}{llrrrrrrrrrr}
	%		\toprule
	%		{} &                         Conjunto de dados &     CH &    DBI &   Dunn &      I &     SD &     XB &  silhouette &  média &   moda &  mediana \\
	%		\midrule
	%		0  &                         50words &    &    &    &    &    &    &         &   * &    &      \\
	%		1  &                           Adiac &    &    &    &    &    &    &         &    &    &      \\
	%		2  &                       ArrowHead &  * &    &    &  * &    &    &         &    &    &      \\
	%	3  &                            Beef &    &    &    &    &    &    &         &   * &    &      \\
	%		4  &                       BeetleFly &  * &    &    &  * &    &    &         &    &    &      \\
	%		5  &                     BirdChicken &    &    &    &    &    &    &         &    &    &      \\
	%		6  &                             CBF &  * &    &    &  * &    &    &       * &    &  * &      \\
	%		7  &                             Car &    &    &    &    &    &    &         &    &    &      \\
	%		8  &           ChlorineConcentration &    &    &  * &    &    &    &         &   * &    &      \\
	%		9  &                          Coffee &  * &    &    &    &    &    &         &    &    &      \\
	%		10 &                       Computers &  * &    &    &    &  * &    &       * &    &  * &      \\
	%		11 &             DiatomSizeReduction &  * &    &    &    &    &    &         &    &    &      \\
	%		12 &    DistalPhalanxOutlineAgeGroup &  * &  * &    &  * &    &  * &       * &   * &  * &     * \\
	%		13 &     DistalPhalanxOutlineCorrect &    &  * &  * &    &  * &  * &       * &    &  * &     * \\
	%		14 &                 DistalPhalanxTW &    &    &    &    &    &    &         &    &    &      \\
	%		15 &                          ECG200 &  * &  * &    &  * &    &  * &       * &    &  * &     * \\
	%		16 &                         ECG5000 &  * &    &    &  * &    &    &         &    &    &      \\
	%		17 &                     ECGFiveDays &  * &    &    &    &    &  * &         &    &  * &      \\
	%		18 &                     Earthquakes &    &  * &  * &    &    &  * &       * &    &  * &     * \\
	%		19 &                            FISH &    &    &    &    &    &    &         &    &    &      \\
	%		20 &                         FaceAll &    &    &    &    &    &    &         &    &    &    \\
	%		21 &                        FaceFour &    &    &    &    &  * &    &         &    &    &      \\
	%		22 &                        FacesUCR &    &    &    &    &    &    &         &    &    &      \\
	%		23 &                           FordB &  * &    &    &  * &    &  * &       * &    &  * &     * \\
	%		24 &                       Gun\_Point &  * &  * &  * &  * &  * &  * &       * &   * &  * &     * \\
	%		25 &                             Ham &    &  * &  * &    &  * &  * &       * &   * &  * &     * \\
	%		26 &                         Herring &    &  * &    &    &    &  * &       * &    &  * &      \\
	%		27 &             InsectWingbeatSound &    &    &    &    &    &    &         &    &    &      \\
	%		28 &                ItalyPowerDemand &  * &    &    &    &  * &  * &       * &    &  * &     * \\
	%		29 &          LargeKitchenAppliances &    &    &    &    &    &    &         &    &    &      \\
	%		30 &                       Lighting2 &  * &  * &    &  * &  * &  * &       * &   * &  * &     * \\
	%		31 &                       Lighting7 &    &    &    &    &    &    &         &    &    &      \\
	%		32 &                            Meat &  * &    &  * &  * &  * &  * &       * &   * &  * &     * \\
	%		33 &                   MedicalImages &    &    &    &    &    &    &         &    &    &      \\
	%		34 &    MiddlePhalanxOutlineAgeGroup &  * &    &    &  * &    &  * &       * &   * &  * &     * \\
	%		35 &     MiddlePhalanxOutlineCorrect &  * &    &  * &    &    &  * &       * &    &  * &     * \\
	%		36 &                 MiddlePhalanxTW &    &    &    &    &    &    &         &    &    &      \\
	%		37 &                      MoteStrain &  * &  * &    &  * &  * &  * &       * &   * &  * &     * \\
	%		38 &                         OSULeaf &    &    &    &    &    &    &         &    &    &      \\
	%		39 &                        OliveOil &    &    &    &    &    &    &         &    &    &      \\
	%		40 &                           Plane &    &    &    &    &    &    &         &    &    &      \\
	%		41 &  ProximalPhalanxOutlineAgeGroup &    &    &    &    &    &    &         &   * &    &      \\
	%		42 &   ProximalPhalanxOutlineCorrect &    &    &    &    &    &    &       * &    &    &      \\
	%		43 &               ProximalPhalanxTW &    &    &    &    &    &    &         &    &    &      \\
	%		44 &            RefrigerationDevices &  * &    &    &    &    &    &         &   * &    &      \\
	%		45 &                      ScreenType &    &    &    &    &    &    &         &    &    &      \\
	%		46 &                     ShapeletSim &  * &  * &    &  * &    &    &         &    &    &      \\			
	%		47 &                       ShapesAll &    &    &    &    &    &    &         &    &    &      \\
	%		48 &          SmallKitchenAppliances &    &    &    &    &    &    &         &   * &    &      \\
	%		49 &            SonyAIBORobotSurface &  * &    &  * &  * &    &  * &       * &    &  * &     * \\
	%		50 &          SonyAIBORobotSurfaceII &  * &    &    &  * &    &    &         &    &    &      \\
	%		51 &                      Strawberry &  * &  * &    &    &    &  * &       * &    &  * &     * \\
	%	52 &                     SwedishLeaf &    &    &    &    &    &    &         &    &    &      \\
	%		53 &                         Symbols &    &    &    &    &    &    &         &    &    &      \\
	%	54 &                ToeSegmentation1 &  * &    &    &    &    &  * &       * &    &  * &      \\
	%		55 &                ToeSegmentation2 &  * &    &    &  * &    &    &       * &    &  * &      \\
	%		56 &                           Trace &    &    &    &    &    &    &         &    &    &      \\
	%		57 &                      TwoLeadECG &  * &    &    &    &  * &  * &       * &    &  * &     * \\
	%	58 &                    Two\_Patterns &    &    &    &    &    &    &         &    &    &      \\
	%		59 &          UWaveGestureLibraryAll &    &    &    &    &    &    &         &    &    &      \\
	%		60 &                            Wine &    &  * &  * &    &    &  * &       * &    &  * &     * \\
	%		61 &                   WordsSynonyms &    &    &    &    &    &    &         &    &    &      \\
	%		62 &                           Worms &    &    &    &    &    &    &         &    &    &      \\
	%		63 &                   WormsTwoClass &    &  * &    &    &  * &  * &       * &    &  * &     * \\
	%		64 &                 easy\_clustering &  * &  * &  * &  * &  * &  * &       * &   * &  * &     * \\
	%		65 &                             foo &  * &    &    &  * &    &  * &       * &    &  * &     * \\
	%		66 &                     synthetic\_6 &    &    &    &    &    &    &         &    &    &      \\
	%		67 &               synthetic\_control &    &    &    &    &    &    &         &    &    &      \\
	%		68 &                           wafer &  * &  * &    &    &    &  * &       * &    &  * &     * \\
	%		69 &                            yoga &  * &    &    &    &    &    &         &    &    &      \\
	%		\rowcolor{lightgray} 70 &                           Total &  0.429 &  0.214 &  0.143 &  0.257 &  0.171 &  0.343 &       0.386 &   0.2 &  0.386 &     0.3 \\
	%		{} &                         {}  &     CH &    DBI &   Dunn &      I &     SD &     XB &  silhouette &  média &   moda &  mediana \\
	%		\bottomrule
%		\end{longtable}\label{tbl:indices_hits} 
%	}
	%}
	%\end{longtable} \label{tbl:indices_iterativos} 

O total da tabela binária obtida pela combinação do algoritmo hierárquico com linkagem \emph{average} e métrica de dissimilaridade cid-euclidiana foi armazenado, assim como os totais obtidos por diversas outras combinações. Com todos os totais, obtidos por diferentes combinações de algoritmos de agrupamento e métricas de dissimilaridade, pôde-se formar uma terceira tabela, cuja as linhas são compostas pelas diferentes combinações de algoritmo e métrica de dissimilaridade, enquanto as colunas são os índices, puros e combinados, utilizados na tarefa iterativa de se apontar o número de agrupamentos. Os valores de tal tabela contêm cada um a porcentagem de vezes que um índice aponta corretamente o número de agrupamentos de um conjunto de dados de séries temporais para determinada combinação de métrica de dissimilaridade e algoritmo de agrupamento. 

A tabela resultante possui $78$ linhas, ou combinações de algoritmos e métricas de dissimilaridade, e $10$ colunas, ou índices testados, o que faz com que o valor crítico para o teste não-paramétrico de Friedman seja $F_{crit} = 1.79$, para uma significância de 5\%. Após a obtenção do ranking médio de cada índice e a realização do teste de Friedman, obteve-se o valor $F_F = 258.42$, o que descarta a hipótese nula $H_0$ de que os índices possuem performance equivalente frente às diferentes estratégias testadas. Dessa maneira, deu-se prosseguimento ao teste \emph{pos-hoc} de Nemenyi, para verificar se existe algum índice, ou grupo de índices, que possua performance superior aos demais. O resultado do teste de Nemenyi pode ser visto na Figura ~\ref{fig:resultado_iterativo}, onde se destaca o o grupo de índices formado pelo Calinski-Harabasz e a moda da sugestão de todos os índices, já que este grupo possui um ranking médio menor do que os demais gruposs.

Logo, conclui-se que o índice Calinski-Harabasz é o mais indicado para a tarefa de se definir o número de grupos em que um conjunto de dados de séries temporais se encontra naturalmente dividido, já que para a obtenção do índice combinado moda, que possui performance equivalente, faz-se necessária a obtenção de todos os demais índices puros além do citado índice, o que, do ponto de vista do custo computacional, não é interessante.

 \begin{figure}[h!]
 	%\vspace{2.5cm}
 	%\resizebox{\columnwidth}{!}{%
 		\includegraphics[width=\linewidth]{figuras/figura_alteracao_centroid.png}
 	%}
 	\caption{Resultado do teste de Nemenyi para o teste do melhor índice interno para apontar o número de grupos.}
 	 	\label{fig:resultado_iterativo}
 \end{figure} 


\subsection{Índice interno para avaliação de partições com o número de grupos definido} \label{sec:indice_interno_k_fixo}

Uma vez definido o número de grupos, surge outra pergunta: qual dos índices internos da Tabela ~\ref{tbl:indices_internos} é o mais indicado para avaliar a qualidade de diferentes partições? Ou, em outras palavras, qual índice interno é o melhor indicador que houve uma melhora, ou piora, na partição resultante após uma alteração nos parâmetros ou estratégias de agrupamento com exceção do número de grupos? Nos experimentos será utilizada a distância de kendall-tau que é uma medida de dissimilaridade entre dois rankings e definida como a contagem do número de discordâncias entre os dois rankings em análise ~\parencite{kendall_tau}. Vale ressaltar que a medida de Kendall-tau, que varia no intervalo $[-1,1]$, é uma métrica no espaço de rankings, de forma que valores maiores implicam em rankings mais próximos. Quando dois rankings são idênticos, o valor da métrica Kendall-Tau entre eles é 1.

Com o objetivo de elucidar a proposta deste experimento, foi realizado o seguinte experimento:
\begin{itemize}
	\item escolheu-se arbitrariamente uma base de dados do UCR \emph{Time Series} (\emph{synthetic_control} que possui 6 grupos, vide Figura ~\ref{fig:synthetic_control}),
	\item escolheu-se arbitrariamente uma métrica de dissimilaridade (distância Euclidiana),
	\item realizou-se seis agrupamentos do conjunto de dados com os seguintes algoritmos: \emph{hierarchical-average (HA), hierarchical-complete (HC), hierarchical-single (HS), k-means++ (KP), k-means (KR), k-medoids(PAM)(KP)},
	\item avaliou-se cada uma das seis partições obtidas com os seguintes índices: VI, silhouette, Calinski-Harabasz.
\end{itemize}

Os resultados numéricos podem ser vistos na Tabela ~\ref{tbl:rankings_3}, bem como os valores ranqueados na Tabela ~\ref{tbl:rankings}. Em seguida, calculou-se a distância Kendall-Tau entre o ranking do índice silhouette e o ranking do índice VI, onde obteve-se o valor -0.20, bem como a distância Kendall-Tau entre o ranking do índice Calinski-Harabasz e o ranking do índice VI, onde obteve-se o valor 0.33. Logo, como valores maiores de Kendall-Tau implicam em rankings mais coerentes entre si, pode-se afirmar que o índice Calinski-Harabasz, ao se variar o algoritmo de agrupamento, na base de dados e métrica de dissimilaridade escolhidas, tem comportamento, no que diz respeito a indicar uma piora ou melhora da partição, mais similar ao índice VI do que o índice silhouette.

A partir deste breve experimento, pode-se afirmar que, caso não se tivesse disponível a rotulação dos dados, mas somente o número de grupos, o índice Calinski-Harabasz é mais indicado que o índice silhouette para indicar uma melhora, ou piora, nas partições obtidas com a variação do algoritmo de agrupamento para a base de dados e métrica de dissimilaridade escolhida.

%\begin{center}
	\begin{table}[!h]
		\centering
		\caption{Rankings dos valores dos índices Para cada uma das partições obtidas variando-se o algoritmo de agrupamento.} 	\label{tbl:rankings_3}
		\begin{tabular}{lrrrrrr}
			\toprule
			{} &          HC &          HA &         HS &          KM &          KR &          KP \\
			\midrule
			VI &   1.41 &   1.39 &  1.84 &   0.84 &   0.73 &   2.57 \\
			Calinski-Harabasz &  42.04 &  49.11 &  1.02 &  55.73 &  55.87 &  26.27 \\
			silhouette &   0.23 &   0.25 & -0.01 &   0.21 &   0.21 &  -0.01 \\
			\bottomrule
		\end{tabular}
			\end{table}
%\end{center}

%\begin{center}
	\begin{table}[!h]
		\centering
		\caption{Rankings dos valores dos índices da tabela ~\ref{tbl:rankings_3}. Note que valores maiores para o Calinski-Harabasz e o silhouette implicam em rankings melhores, ao passo que para o VI em rankings piores.} 	\label{tbl:rankings}
		\begin{tabular}{lrrrrrr}
			\toprule
			{} &          HC &          HA &         HS &          KM &          KR &          KP \\			
			\midrule
			VI &  5 &  4 &  2 &  1 &  3 &  6 \\
			Calinski-Harabasz &  5 &  4 &  2 &  1 &  6 &  3 \\
			silhouette &  2 &  1 &  5 &  4 &  6 &  3 \\
			\bottomrule
		\end{tabular}
	\end{table}

%\end{center}

O experimento realizado restringiu-se a uma base  de dados, à dois índices internos e à avaliação da variação do algoritmo de agrupamento. Para realizar um estudo mais amplo, foram realizados agrupamentos com um pequeno grupo de $55$ combinações de algoritmo de agrupamento e métrica de dissimilaridade e de $71$ conjuntos de dados de séries temporais, novamente conjuntos de dados do UCR \emph{Time Series} ~\parencite{UCRArchive}. Cada partição obtida foi avaliada por cada um dos índices internos da Tabela ~\ref{tbl:indices_internos} e pelo índice externo VI exposto na seção ~\ref{sec:VI}. Como o VI é uma métrica no espaço de partições, uma melhora ou piora no seu valor indica uma respectiva melhora ou piora na aproximação da partição real em função de uma estratégia de agrupamento melhor ou pior para determinada base de dados. Dessa maneira, espera-se que os índices internos que acompanharem a melhora ou piora do VI em função das variações da estratégia de agrupamento sejam os mais indicados para avaliar partições com o número de grupos definido. 

Para se comparar a aderência, ou coerência, de cada índice interno ao VI, todos os índices para cada base foram ranqueados em função das estratégias escolhidas. Em seguida, para cada base e para cada índice foi calculada a distância de Kendall-Tau entre os rankings obtidos pelo VI e os rankings obtidos para cada índice. Os resultados podem ser vistos na Tabela ~\ref{tbl:coherence}, que contêm os valores $1-\tau$, onde $\tau$ é o valor Kendall-Tau, de forma que valores menores indicam uma maior coerência do índice interno com o índice VI.

\begin{longtable}{llllllll}
	\caption{Tabela com os valores de Kendall-Tau dos índices em relação ao VI.}
	 \label{tbl:coherence}\\
	% \endfirsthead
%	 \endhead
	\toprule
	{} &                         dataset &    CH &   DBI & silhouette &    SD &    XB &  Dunn \\
	\midrule
	0  &                         50words &  0.91 &  1.13 &       0.91 &  1.24 &  1.08 &  0.93 \\
	1  &                           Adiac &  0.98 &  0.68 &       1.03 &  0.92 &  0.74 &  1.08 \\
	2  &                       ArrowHead &  1.04 &  0.95 &       0.78 &  1.14 &  1.02 &  1.01 \\
	3  &                       BeetleFly &  0.87 &  1.11 &       1.11 &  1.11 &  1.29 &  0.97 \\
	4  &                             CBF &  1.11 &  1.15 &       0.74 &  0.84 &  1.05 &  0.92 \\
	5  &                             Car &  1.07 &  1.02 &       0.95 &  0.96 &  1.07 &  0.84 \\
	6  &           ChlorineConcentration &  1.08 &  1.08 &       0.66 &  1.09 &  0.98 &  0.81 \\
	7  &                  CinC\_ECG\_torso &  0.85 &  0.94 &       0.79 &  1.17 &  1.00 &  0.78 \\
	8  &                          Coffee &  0.82 &  1.08 &       0.81 &  1.00 &  0.99 &  1.09 \\
	9  &                       Computers &  1.02 &  1.06 &       0.97 &  1.15 &  1.08 &  0.79 \\
	10 &             DiatomSizeReduction &  0.68 &  1.11 &       0.58 &  1.21 &  1.19 &  0.60 \\
	11 &    DistalPhalanxOutlineAgeGroup &  1.08 &  0.87 &       0.66 &  0.91 &  0.93 &  0.82 \\
	12 &     DistalPhalanxOutlineCorrect &  1.34 &  1.44 &       0.53 &  1.01 &  1.40 &  0.69 \\
	13 &                 DistalPhalanxTW &  1.01 &  1.17 &       0.92 &  0.90 &  1.11 &  0.92 \\
	14 &                          ECG200 &  1.14 &  1.19 &       0.63 &  1.02 &  1.15 &  0.77 \\
	15 &                         ECG5000 &  1.17 &  1.07 &       0.69 &  1.02 &  1.16 &  0.96 \\
	16 &                     ECGFiveDays &  1.16 &  1.05 &       1.24 &  0.85 &  0.77 &  1.26 \\
	17 &                     Earthquakes &  1.03 &  0.79 &       0.86 &  1.19 &  0.84 &  1.03 \\
	18 &                            FISH &  1.07 &  1.07 &       0.76 &  0.95 &  1.23 &  0.90 \\
	19 &                        FaceFour &  1.01 &  1.10 &       0.93 &  1.12 &  1.07 &  1.00 \\
	20 &                        FacesUCR &  1.13 &  1.01 &       0.82 &  0.94 &  1.08 &  0.99 \\
	21 &                           FordA &  1.07 &  0.98 &       1.03 &  1.17 &  0.93 &  0.83 \\
	22 &                           FordB &  0.72 &  0.78 &       1.12 &  1.04 &  0.88 &  0.82 \\
	23 &                             Ham &  1.49 &  1.48 &       0.66 &  1.10 &  1.41 &  0.61 \\
	24 &                    HandOutlines &  1.32 &  1.27 &       0.61 &  1.03 &  1.27 &  0.61 \\
	25 &                         Haptics &  1.22 &  1.22 &       0.95 &  0.95 &  1.31 &  0.81 \\
	26 &                         Herring &  1.31 &  1.38 &       0.76 &  0.97 &  1.35 &  0.73 \\
	27 &                     InlineSkate &  1.13 &  1.30 &       0.65 &  0.98 &  1.23 &  1.01 \\
	28 &                ItalyPowerDemand &  0.95 &  1.20 &       0.84 &  1.15 &  1.12 &  0.82 \\
	29 &          LargeKitchenAppliances &  1.04 &  0.88 &       0.92 &  0.84 &  0.93 &  0.72 \\
	30 &                       Lighting2 &  0.89 &  1.11 &       0.88 &  0.81 &  1.16 &  0.93 \\
	31 &                       Lighting7 &  0.96 &  1.04 &       0.87 &  1.02 &  1.21 &  0.83 \\
	32 &                          MALLAT &  0.81 &  0.94 &       0.92 &  0.99 &  1.01 &  1.10 \\
	33 &                            Meat &  0.85 &  1.30 &       0.68 &  1.03 &  1.25 &  0.69 \\
	34 &                   MedicalImages &  1.06 &  1.11 &       0.80 &  1.01 &  1.03 &  0.72 \\
	35 &    MiddlePhalanxOutlineAgeGroup &  0.74 &  1.15 &       0.85 &  0.87 &  1.15 &  1.07 \\
	36 &     MiddlePhalanxOutlineCorrect &  1.12 &  1.10 &       0.96 &  0.90 &  0.93 &  0.41 \\
	37 &                 MiddlePhalanxTW &  1.17 &  1.38 &       0.83 &  0.84 &  1.19 &  0.64 \\
	38 &                      MoteStrain &  0.73 &  0.75 &       1.10 &  1.09 &  0.75 &  0.79 \\
	39 &     NonInvasiveFatalECG\_Thorax1 &  0.93 &  1.02 &       0.96 &  1.05 &  0.89 &  0.98 \\
	40 &     NonInvasiveFatalECG\_Thorax2 &  1.07 &  0.92 &       0.85 &  0.89 &  0.76 &  1.03 \\
	41 &                         OSULeaf &  1.03 &  1.15 &       0.82 &  0.77 &  1.04 &  0.99 \\
	42 &        PhalangesOutlinesCorrect &  1.37 &  1.40 &       0.86 &  1.30 &  1.30 &  0.74 \\
	43 &                         Phoneme &  0.94 &  1.04 &       0.84 &  1.20 &  1.03 &  1.06 \\
	44 &                           Plane &  0.84 &  1.02 &       0.74 &  1.14 &  1.06 &  0.85 \\
	45 &  ProximalPhalanxOutlineAgeGroup &  0.98 &  1.01 &       0.59 &  1.18 &  0.96 &  0.75 \\
	46 &   ProximalPhalanxOutlineCorrect &  1.05 &  0.72 &       1.01 &  1.12 &  0.95 &  0.87 \\
	47 &               ProximalPhalanxTW &  0.96 &  1.07 &       0.79 &  1.17 &  1.00 &  0.71 \\
	48 &            RefrigerationDevices &  1.11 &  1.18 &       0.91 &  1.05 &  1.18 &  0.87 \\
	49 &                      ScreenType &  1.14 &  0.96 &       0.88 &  1.15 &  0.97 &  1.02 \\
	50 &                       ShapesAll &  0.64 &  1.10 &       1.00 &  0.99 &  0.86 &  0.88 \\
	51 &          SmallKitchenAppliances &  1.15 &  1.03 &       0.72 &  0.93 &  1.06 &  0.92 \\
	52 &            SonyAIBORobotSurface &  1.16 &  1.01 &       0.50 &  0.93 &  0.88 &  0.74 \\
	53 &          SonyAIBORobotSurfaceII &  0.81 &  0.94 &       0.95 &  1.07 &  1.12 &  0.87 \\
	54 &                 StarLightCurves &  0.72 &  0.87 &       0.75 &  1.12 &  1.16 &  1.11 \\
	55 &                      Strawberry &  1.08 &  1.08 &       0.48 &  0.83 &  1.03 &  0.90 \\
	56 &                     SwedishLeaf &  0.72 &  1.02 &       0.96 &  0.96 &  0.99 &  1.01 \\
	57 &                         Symbols &  0.66 &  1.08 &       0.65 &  0.94 &  1.24 &  0.82 \\
	58 &                ToeSegmentation1 &  0.82 &  1.14 &       1.31 &  1.23 &  1.13 &  0.87 \\
	59 &                ToeSegmentation2 &  0.88 &  0.98 &       1.12 &  0.98 &  1.04 &  1.13 \\
	60 &                           Trace &  0.83 &  0.81 &       0.86 &  0.92 &  0.93 &  1.16 \\
	61 &                      TwoLeadECG &  1.06 &  1.13 &       0.78 &  0.95 &  1.15 &  0.61 \\
	62 &                    Two\_Patterns &  0.97 &  1.14 &       0.90 &  1.12 &  0.96 &  0.89 \\
	63 &          UWaveGestureLibraryAll &  0.82 &  1.03 &       0.78 &  1.09 &  0.97 &  0.88 \\
	64 &                            Wine &  1.39 &  1.36 &       0.67 &  1.15 &  1.30 &  0.56 \\
	65 &                   WordsSynonyms &  0.93 &  1.14 &       0.82 &  0.90 &  1.25 &  0.97 \\
	66 &                           Worms &  0.86 &  1.12 &       0.74 &  0.96 &  1.08 &  1.05 \\
	67 &                   WormsTwoClass &  0.73 &  0.84 &       0.87 &  1.14 &  0.89 &  1.09 \\
	68 &               synthetic\_control &  1.15 &  1.01 &       0.58 &  0.88 &  1.09 &  0.85 \\
	69 &                           wafer &  1.08 &  1.16 &       0.80 &  0.94 &  1.02 &  0.94 \\
	70 &                            yoga &  0.89 &  0.99 &       0.66 &  1.10 &  1.01 &  0.85 \\
	\midrule
	 \rowcolor{lightgray}    &			Rankings médios & 3.58 & 4.42 & 2.16 & 3.91 & 4.23 & 2.67 \\
	 	\midrule
	 {} &                         {} &    CH &   DBI & silhouette &    SD &    XB &  Dunn \\
	\bottomrule
\end{longtable}

Para $70$ bases e $6$ índices internos, o valor crítico da estatística de Friedman é de $F_{crit}=2.12$, e como  para os rankings médios destacados na Tabela ~\ref{tbl:coherence} foi obtido o valor $F_F = 18.11$, então a hipótese nula $H_0$ de que todos os índices possuem aderências equivalentes ao índice VI foi descartada. A partir desse resultado, deu-se prosseguimento ao teste de Nemenyi para verificar qual, ou quais, índice interno possui aderência significativamente maior que os demais índices. Os resultados podem ser vistos na Figura ~ \ref{fig:coherence}, que mostram que os índices silhouette e Dunn, possuem melhor aderência ao VI que os demais índices internos.

Logo, em uma tarefa de agrupamento real, no qual não se possui a rotulação dos dados e já se decidiu, ou já se sabe \emph{a priori}, o número de grupos que os dados se encontram divididos, os índices silhouette e Dunn são os mais indicados para avaliar as novas partições obtidas por diferentes estratégias de agrupamento.

 \begin{figure}[h!]
 	%\vspace{2.5cm}
 	%\resizebox{\columnwidth}{!}{%
 	\includegraphics[width=\linewidth]{figuras/best_coherence.png}
 	%}
 	\caption{Resultado do teste de Nemenyi para o teste de coerência dos índices internos com o VI.}
 	\label{fig:coherence}
 \end{figure}

\section{Conclusões do Capítulo}

Os experimentos deste Capítulo buscaram definir uma estratégia padrão de agrupamento de séries temporais, e as estratégias testadas foram confrontadas com a estratégia padrão, que consiste na realização do agrupamento pelo algoritmo \emph{k-means} e distância Euclidiana. Dessa maneira, foram realizados dois experimentos, sendo o primeiro para estudar as influências da escolha do algoritmo e métrica de dissimilaridade na qualidade da partição de dados de séries temporais, enquanto o segundo experimento teve como objetivo estudar os índices de validação interna na avaliação das partições resultantes do agrupamento de séries temporais.

Assim, no primeiro experimento, foram realizados dois testes:
\begin{itemize}
	\item no primeiro foram testados diversos algoritmos de agrupamento fixando-se a distância Euclidiana como métrica de dissimilaridade e verificou-se a superioridade do algoritmo hierárquico com \emph{linkagem average},
	\item já no segundo teste, fixou-se o algoritmo de agrupamento, como o algoritmo hierárquico com \emph{linkagem average}, e variou-se as dissimilaridades. onde  verificou-se a superioridade da performance das dissimilaridades DTW e CID-Euclidiana em relação às demais dissimilaridades testadas.
\end{itemize}

No segundo experimento, também foram realizados dois testes:
\begin{itemize}
	\item no primeiro, os índices de validação interna foram testados para indicar aquele que é o mais adequado para a tarefa de apontar o número de grupos que as instâncias do conjunto de dados se encontram naturalmente divididas, e, após a análise dos resultados, chegou-se a conclusão que o índice Calinski-Harabasz é o mais indicado para tal,
	\item no segundo teste, procurou-se encontrar um índice da validação interna que seja mais coerente com o VI e, dessa maneira, um índice que seja um melhor indicador da melhora ou piora da partição dos dados de séries temporais ao se alterar a estratégia de agrupamento, chegou-se a conclusão que os índices Dunn e silhouette são os índices mais coerentes com o VI.
\end{itemize}

Dessa maneira, partindo-se do pressuposto que não se possui, \emph{a priori}, nenhuma informação útil para a realização do agrupamento, como, por exemplo, o número de grupos que as instâncias se encontram naturalmente divididas, os experimentos realizados em bases de dados rotuladas sugerem a seguinte abordagem a ser realizada no agrupamento de séries temporais:

\begin{itemize}
	\item Realizar algum tipo de normalização das curvas. Sugere-se a normalização Z ou normalização max.
	\item Escolher alguma medida de dissimilaridade. Sugere-se a CID-DTW caso o conjunto de dados não seja muito grande, ou os recursos computacionais escassos. Em qualquer uma dessas duas hipóteses, sugere-se a utilização da CID-Euclidiana.
	\item Para se definir o número de grupos $k$, realizar o agrupamento do conjunto de dados já normalizados pelo algoritmo hierárquico com \emph{linkagem average} e medida de dissimilaridade escolhida. Em seguida, obter diferentes partições a partir da poda do dendrograma obtido por valores iterativos de $k$, em um intervalo no qual se espera que se encontre o valor de $k$.
	\item Em seguida, para cada uma das partições obtidas, realizar a avaliação de cada uma delas pelo índice Calinski-Harabasz, e aquela que possuir o melhor valor indicará o valor de $k$.
	\item Caso se deseje realizar um refinamento da partição obtida até o momento, uma vez definido o número de grupos $k$, pode-se realizar outras partições com diferentes escolhas de combinações de algoritmos e métricas de dissimilaridade e avaliá-las pelo índice silhouette. A partição que acarretar no melhor valor do índice é a partição final da tarefa de agrupamento.
\end{itemize}

A estratégia descrita foi empregada no estudo de caso do Capítulo ~\ref{cap:estudo_de_caso}. Vale destacar que esta estratégia não garante resultados ótimos, ela é apenas uma sugestão inicial quando não se dispõe de informações relevantes para o agrupamento acerca dos dados. Nos próprios dados sintéticos do UCR \emph{Time Series}, verificou-se que para alguns conjuntos de dados em específico, a metologia proposta teve performance inferior à performance da estratégia padrão (\emph{k-means} com distância Euclidiana). No entanto, na média, levando-se em conta todos os conjuntos, os testes estatísticos demonstraram que ela é significativamente superior às demais estratégias, como a estratégia padrão, por exemplo.



 %Mesmo para outros algoritmos de agrupamento, baseados em densidade por exemplo, como o DBSCAN, que fazem uso de outros parâmetros que não o número de agrupamentos, os índices de validação são importantes. Em ~\parencite{DE_DBSCAN}, por exemplo, é proposto um algoritmo de evolução diferencial para escolher os melhores parâmetros do DBSCAN, e a função objetivo deste problema de otimização é baseada em índices de validação interna.







%\chapter{Estudo de caso}

Neste capítulo será realizado, enfim, o agrupamento dos dados de curva de carga da Irlanda. O conjunto de dados da Irlanda é composto por medidas periódicas do consumo de energia elétrica de consumidores residenciais em diversas cidades irlandesas. O conjunto de dados contempla as medidas do dia $1^{\circ}$  de Janeiro de 2009 até o dia $1^{\circ}$ de Janeiro de 2011, possuindo, dessa maneira, $2$ anos de medições. Todas as medições foram realizadas em intervalos fixos de $30$ minutos, no início de cada hora e na metade delas, totalizando, assim, $48$ medidas diárias.

A carga consumida a cada $30$ minutos é a única informação disponível para a realização da tarefa de agrupamento, sendo todas as demais, como a localização, renda ou número de moradores em cada residência, não fornecidas por questões de privacidade. No estudo de caso, o objetivo é de se fazer uma análise de um ano completo do comportamento diário de cada carga, bem como avaliar as variações deste comportamento em função das estações do ano e do fato da curva diária ser ou não uma curva do final de semana. Assim, somente as medições de 21 de Dezembro de 2009 até o mesmo dia do ano de 2010 foram analisadas. Este dia foi escolhido por ele ser o dia do solstício de inverno no hemisfério norte, e dessa maneira, teremos o mesmo número de medições diárias para cada estação do ano.

Para se analisar as variações das cuvas de cargas diárias em função da estação do ano e dos dias da semana, o conjunto de dados foi divido em seis subconjuntos e as estratégias de agrupamento adotadas serão as mesmas das definidas no capítulo ~\ref{cap:testes_teoricos}. No entanto, antes delas serem aplicadas, faz-se necessária uma etapa de pré-processamento dos dados, que será melhor detalhada na seção seguinte.

\section{Pré-processamento}

No conjunto de dados da Irlanda, dentro do intervalo de um ano analisado, existem dados incompletos ou inconsistentes de alguns medidores, o que fez com que fosse necessária uma etapa de pré-processamento dos dados brutos fornecidos. Dessa maneira, medidores que apresentavam dados faltantes, ou seja, instantes de tempo para os quais não existe medição, tiveram o mesmo tratamento realizado em ~\parencite{Flath2012}: medidores com curvas diárias que tiverem mais de $4$ medições consecutivas ausentes são descartados, e para aqueles que tiverem $4$ ou menos medições consecutivas ausentes é feita uma regressão linear simples para se estimar o(s) valore(s) naquele(s) instante(s).

Após esta primeira etapa de pré-processamento, o conjunto de dados que continha $6435$ medidores foi reduzido para $5608$ medidores. Em seguida, o conjunto de dados foi dividido nas estações do ano: inverno, verão e transição, sendo esta última entendida como a união dos dados da primavera e do outono. Cada um dos três subconjuntos gerados foi dividido, novamente, em medições de dia de semana e medições de finais de semana e, assim, ao final, obteve-se um total de $6$ subconjuntos de dados para serem agrupados. Na figura ~\ref{fig:6938_raw} podem ser vistas as curvas do medidor 6938 obtidas nos dias de finais de semana do verão de 2010. Na figura ~\ref{fig:6938_Z} podem ser vistas as mesmas curvas após a normalização Z.

\begin{figure}
	\centering
	\begin{subfigure}{.5\textwidth}
		\centering
		\includegraphics[width=.9\linewidth]{figuras/irish/raw_data_6938.png}
		\caption{Versão não normalizada.}
		\label{fig:6938_raw}
	\end{subfigure}%
	\begin{subfigure}{.5\textwidth}
		\centering
		\includegraphics[width=.9\linewidth]{figuras/irish/Z_norm_6938.png}
		\caption{Versão z-normalizada.}
		\label{fig:6938_Z}
	\end{subfigure}
	\caption{Curvas de consumo diárias medidas pelo medidor 6938 no verão de 2010 contendo somente curvas de dias dos finais de semana (sábado e domingo).}
	\label{fig:Z_norm}
\end{figure}

%\begin{figure}[h!]
%	\includegraphics[width=\linewidth]{figuras/irish/raw_data_6938.png}
%	\caption{Curvas de consumo diárias medidas pelo medidor 6938 no verão de 2010 contendo somente curvas de dias dos finais de semana (sábado e domingo).}
%	\label{fig:6938_raw}
%\end{figure}

%\begin{figure}[h!]
%	\includegraphics[width=\linewidth]{figuras/irish/Z_norm_6938.png}
%	\caption{Versão z-normalizada das curvas da figura ~\ref{fig:6938_raw}}
%	\label{fig:6938_Z}
%\end{figure}

Após a criação desses sub-conjuntos de dados, foi escolhida uma curva diária de cada medidor para cada sub-conjunto, de forma que tal curva seja o mais representativa possível do medidor naquele subconjunto. Diferentemente de ~\parencite{Flath2012}, onde obteve-se o centróide para cada medidor em cada subconjunto e em seguida os dados foram normalizados pela normalização max \hl{TODO: citar normalizacao max}, neste trabalho,  primeiramente, todas as curvas diárias de cada medidor foram z-normalizadas, conforme descrito na seção ~\ref{sec:norm_Z}. Após esta etapa de normalização, as distâncias cid-euclidiana entre as curvas de cada medidor, em cada um dos $6$ sub-conjuntos foram calculadas, de forma que o medóide de cada medidor foi obtido e escolhido como o mais representativo de  cada medidor em cada subconjunto de dados. O medóide encontrado para as curvas do medidor 6938 representados na figura ~\ref{fig:6938_Z} pode ser visto na figura ~\ref{fig:6938_medoid},

\begin{figure}[h!]
	\includegraphics[width=\linewidth]{figuras/irish/medoid_6938.png}
	\caption{Medóide encontrado para as curvas da figura ~\ref{fig:6938_Z}}
	\label{fig:6938_medoid}
\end{figure}


Escolheu-se esta estratégia pelo fato todos os estudos terem sido feitos com a normalização Z e pela reconhecimento de sua superioridade em relação às demais normalizações no contexto de séries temporais. Já a preferência do medóide em detrimento do centróide, se deu pelo fato de esse ser menos sensível a \emph{outliers} do que este, de forma que se, dentro de determinada estação do ano, em um dia de semana por exemplo, houver algum consumo de energia atípico, como ocorre em feriados, tal efeito será desconsiderado, ou pelo menos atenuado, caso se escolha o medóide como curva mais representativa. No entanto, o mesmo não poderia ser dito caso fosse escolhido o centróide como a curva mais representativa.

Após esta etapa da escolha da curva do medidor mais representativo de cada sub-conjunto fez-se necessário o descarte de mais alguns medidores com dados anômalos. Primeiro todos os medidores que apresentavam valores constantes foram descartados, pois estes formam um grupo pequeno que possui um perfil muito bem definido e significativamente diferente dos perfis de consumo dos demais medidores. Caso os medidores com valores de medição constante não fossem descartados, isso faria com que a tarefa de agrupamento indicasse a presença de apenas $2$ grupos muito bem definidos, sendo um dos grupos formado por um pequeno número de medidores com valores constantes e o outro grupo com todos os outros medidores que possuem variações em sua medição. O número de medidores descartados em cada subconjunto nesta etapa se encontra detalhado na coluna 'Constantes descartados' da tabela ~\ref{tbl:sumario_descarte}.

Após o descarte dos medidores com medidas constantes foi calculada a matriz de dissimilaridade para cada um dos 6 subconjuntos, utilizando-se a métrica cid-euclidiana. Uma vez obtida a matriz de dissimilaridade para cada subconjunto, foi encontrado o medóide de cada subconjunto. Uma vez definido o medóide, que é a curva diária mais representativa de todos os medidores daquele subconjunto, calculou-se a média e o desvio padrão das dissimilaridades entre cada instância e o medóide. A partir daí, verificou-se que existem alguns medidores que se encontram significativamente afastados do medóide, de maneira que estes também foram descartados. Foi calculada a dissimilaridade média e o desvio padrão das instâncias até o medóide, e os medidores que apresentavam uma dissimlaridade cid-euclidiana em relação ao medóide maior que a média mais duas vezes o desvio padrão foram descartados. O número de medidores descartados nesta segundo etapa para cada subconjunto se encontra detalhado na coluna '\emph{outliers} descartados' da tabela ~\ref{tbl:sumario_descarte}. \hl{TODO: Ver com Cristiano se deve-se utilizar a distancia Mahalanobis aqui. Acho que nao}

Finalmente, na última etapa do pré-processamento, com a finalidade de manter-se a coerência entre os subconjuntos, ou seja, de que todos os subconjuntos contenham os dados de medição dos mesmos medidores, alguns medidores ainda foram descartados. O número de medidores descartados nesta etapa final em cada subconjunto e em cada uma das etapas de pré-processamento anteriores se encontra sumarizado na coluna 'Descartados na etapa final' da tabela ~\ref{tbl:sumario_descarte}. Ao final da etapa de pré-processamento todos os subconjuntos continham as medidas dos mesmos $4996$ medidores do conjunto inicial, ou seja, na etapa de pré-processamento foram descartados $1439$ medidores do conjunto original da Irlanda o que representa, aproximadamente, $22\%$ do número de medidores inicial.

\begin{center}
	\begin{table}			
		\caption{Tabela que sumariza os números de medidores descartados em cada uma das etapas de pré-processamento em cada um dos subconjuntos. Ao final, todos os subconjuntos contêm os mesmos $4996$ medidores. } 
	\resizebox{\columnwidth}{!}{%
	\begin{tabular}{llllll}
		\toprule
		Id &          Estação do Ano &              Dia de semana? & Constantes descartados & \emph{outliers} descartados & Descartados na etapa final\\
		\midrule
		S\_weekend  &  Verão       &  Não &             71 & 281 &  256\\
		S\_weekday  &   Verão      &  Sim &               63 & 233 & 316  \\
		W\_weekend &  Inverno    &  Não &               52 & 203 & 253\\
		W\_weekday &  Inverno    &  Sim &            41 & 65 & 306\\
		T\_weekend  &  Transição &   Não &               89 & 291 & 228\\
		T\_weekday  &   Transição &  Sim &          81 & 260 & 271\\
		\bottomrule
	\end{tabular}
	}
		\end{table} \label{tbl:sumario_descarte}
\end{center}

\section{Agrupamento das curvas de carga}

Após a etapa de pré-processamento, faz-se necessária a definição do número de grupos de cada um dos $6$ subconjuntos, já que é possível, e se espera, que existam subconjuntos com um número de grupos diferentes. Para tal, iremos utilizar as estratégias definidas na seção ~\ref{sec:definicao_indice_interno}, no qual cada subconjunto será agrupado pelo algoritmo hierárquico com linkagem \emph{average}, e o dendrograma resultante será podado para valores iterativos de $k$, que é o número de grupos. O valor de $k$ que acarretar no melhor valor do índice de Calinski-Harabasz será o número de grupos daquele subconjunto. Como dito, a escolha do índice se baseia nos experimentos da seção ~\ref{sec:indice_iterativo}, pois em ~\parencite{Flath2012} foi utilizado o índice Davies-Bouldin.

O intervalo de número de grupos testado foi de $2$ grupos até $9$ grupos, pois este é um intervalo no qual faz sentido a divisão de perfis de carga para a criação de políticas de tarifas diferentes para cada grupo. Ou seja, é possível que um valor de $k$ maior do que o limite superior do intervalo possua um valor melhor para o índice Calinski-Harabasz, no entanto, para a aplicação em questão, ele não faz sentido. O número de grupos escolhido para cada subconjunto se encontram na tabela ~\ref{tbl:valores_otimos_k} e os valores iterativos detalhados na tabela ~\ref{tbl:valores_iterativos_k}.

\begin{center}
			\begin{table}
				\caption{Tabela que mostra a variação do índice Calinski-Harabasz para valores iterativos do número de grupos. Os valores ótimos para cada subconjunto se encontram destacados. } 
%	\resizebox{\columnwidth}{!}{% 
		\begin{tabular}{llrrrrrrrr}
			\toprule
			{} &    subconjunto &    k=2 &     k=3 &      k=4 &     k=5 &     k=6 &     k=7 &     k=8 &     k=9 \\
			\midrule
			 &  S\_weekday &   7.31 &   47.14 &    28.17 &  \bfseries{799.96 } &  647.15 &  551.76 &  478.84 &  433.44 \\
			 &  S\_weekend &  83.77 &   60.10 &  \bfseries{1216.29} &  913.81 &  753.17 &  629.40 &  539.54 &  494.46 \\
			 &  T\_weekday &  51.96 &  172.64 &   128.53 &   96.47 &  101.20 &   84.51 &  \bfseries{469.90} &  411.93 \\
			 &  T\_weekend &  16.66 &   68.36 &    52.28 &  191.03 &  \bfseries{858.14} &  715.85 &  614.17 &  537.50 \\
			 &  W\_weekday &   7.74 &   44.12 &    30.20 &   24.00 &   19.21 &   15.86 &  \bfseries{549.97} &  481.31 \\
			 &  W\_weekend &   1.44 &   30.93 &    \bfseries{44.24} &   35.17 &   27.92 &   31.70 &   27.12 &   23.68 \\
			\bottomrule
		\end{tabular} 	
\end{table}\label{tbl:valores_iterativos_k}
%	}

\end{center}

\begin{center}
	\begin{table}
		\caption{Tabela que mostra os valores ótimos de cada subconjunto segundo o índice de Calinski-Harabasz.} 		
%	\resizebox{\columnwidth}{!}{%
\begin{tabular}{lll}
 	\toprule
 	{} &    subconjunto & $k\_\{opt\}$ \\
 	\midrule
 	 &  S\_weekday &       5 \\
 	 &  S\_weekend &       4 \\
 	 &  T\_weekday &       8 \\
 	 &  T\_weekend &       6 \\
 	 &  W\_weekday &       8 \\
 	 &  W\_weekend &       4 \\
 	\bottomrule
\end{tabular}	
	\end{table}\label{tbl:valores_otimos_k}
%	} 
\end{center} 	

Uma vez definido o número de grupos para cada subconjunto, agora para se encontrar a melhor partição realizaremos os procedimentos definidos na seção ~\ref{sec:indice_interno_k_fixo}. Nela chegou-se a conclusão que a melhor partição de um conjunto de dados de séries temporais é obtida ao se variar as estratégias de agrupamento para valores fixos de $k$, e a estratégia que apresentar os melhores valores do índice de silhouette será a partição escolhida. Por estratégia de agrupamento entende-se as variações de algoritmos de agrupamento e métrica de dissimilaridade. As estratégias testadas bem como os valores de silhouette para cada subconjunto podem ser vistos  na tabela ~\ref{tbl:valores_silhouette}.

\begin{center}
	\begin{table}
		\caption{Valores para o silhouette para diferentes estratégias em cada de cada subconjunto, onde HA representa a escolha do algoritmo hieráquico com linkagem \emph{average}.} 		
	\resizebox{\columnwidth}{!}{%
\begin{tabular}{lrrrrrrrrrrrrr}
	\toprule
	dataset & k-means,eucl & HA,DTW & HA,EDR & HA,chebyshev & HA,cid-DTW & HA,cid-EDR & HA,cid-chebyshev & HA,cid-eucl & HA,cort-DTW & HA,cort-EDR & HA,cort-chebyshev & HA,cort-eucl & HA,eucl \\
	\midrule
	S\_weekday &                0.06 &                     0.17 &                     0.17 &                           0.07 &                         0.19 &                         \bfseries{0.21} &                               0.09 &                                 0.07 &                          0.11 &                          0.17 &                                0.05 &                                  0.05 &                             0.05 \\
	S\_weekend &                0.06 &                     0.15 &                     0.25 &                           0.08 &                         0.20 &                         \bfseries{0.31} &                               0.08 &                                 0.13 &                          0.17 &                          0.24 &                                0.05 &                                  0.04 &                             0.06 \\
	T\_weekday &                0.06 &                     0.14 &                     \bfseries{0.15} &                           0.07 &                         0.11 &                         0.13 &                               0.04 &                                 0.04 &                          0.06 &                          0.12 &                                0.02 &                                  0.03 &                             0.05 \\
	T\_weekend &                0.06 &                     0.15 &                     0.21 &                           0.04 &                         0.20 &                         \bfseries{0.24} &                               0.05 &                                 0.10 &                          0.11 &                          0.17 &                                0.03 &                                  0.04 &                             0.06 \\
	W\_weekday &                0.05 &                     0.15 &                     0.15 &                           0.12 &                         \bfseries{0.22} &                         0.10 &                               0.13 &                                 0.05 &                          0.16 &                          0.14 &                                0.03 &                                  0.06 &                             0.10 \\
	W\_weekend &                0.05 &                     \bfseries{0.26} &                     0.23 &                           0.17 &                         0.16 &                         0.17 &                               0.20 &                                 0.15 &                          0.25 &                          0.21 &                                0.13 &                                  0.16 &                             0.15 \\
	\bottomrule
\end{tabular}
	}
	\end{table}\label{tbl:valores_silhouette}
\end{center} 

\begin{center}
	\begin{table}
		\caption{Tabela que mostra os resultados do agrupamento de cada subconjunto.} 		
		%	\resizebox{\columnwidth}{!}{%
		\begin{tabular}{lllll}
			\toprule
			subconjunto & $k\_\{opt\}$ & Algoritmo de agrupamento & Medida de dissimilaridade & Silhouette \\
			\midrule
			 S\_weekday &       5 & hierarchical-average & cid-EDR & 0.21 \\
			 S\_weekend &       4 & hierarchical-average & cid-EDR & 0.31\\
			 T\_weekday &       8 & hierarchical-average & EDR & 0.15 \\
			 T\_weekend &       6 & hierarchical-average & cid-EDR & 0.24 \\
			 W\_weekday &       8 & hierarchical-average & cid-DTW & 0.22 \\
			 W\_weekend &       4 & hierarchical-average & DTW & 0.25\\
			\bottomrule
		\end{tabular}	
	\end{table}\label{tbl:resultados_finais_agrupamento}
	%	} 
\end{center} 

Os medóides de cada grupo bem como o número de cargas em cada grupo podem ser vistos nas figuras a seguir.

\begin{figure}[h!]
	\includegraphics[width=\linewidth]{figuras/irish/2016-07-19_15:33:04.743952__irish_results/S_weekday_-_[cid-EDR,hierarchical-average]_-_5_grupos.png}
	\caption{Medóide encontrado para as curvas...}
	%\label{fig:6938medoid}
\end{figure}

\begin{figure}[h!]
	\includegraphics[width=\linewidth]{figuras/irish/2016-07-19_15:33:04.743952__irish_results/S_weekend_-_[cid-EDR,hierarchical-average]_-_4_grupos.png}
	\caption{Medóide encontrado para as curvas...}
	%\label{fig:6938medoid}
\end{figure}

\begin{figure}[h!]
	\includegraphics[width=\linewidth]{figuras/irish/2016-07-19_15:33:04.743952__irish_results/T_weekday_-_[EDR,hierarchical-average]_-_8_grupos.png}
	\caption{Medóide encontrado para as curvas...}
	%\label{fig:6938medoid}
\end{figure}

\begin{figure}[h!]
	\includegraphics[width=\linewidth]{figuras/irish/2016-07-19_15:33:04.743952__irish_results/T_weekend_-_[cid-EDR,hierarchical-average]_-_6_grupos.png}
	\caption{Medóide encontrado para as curvas...}
	%\label{fig:6938medoid}
\end{figure}

\begin{figure}[h!]
	\includegraphics[width=\linewidth]{figuras/irish/2016-07-19_15:33:04.743952__irish_results/W_weekday_-_[cid-DTW,hierarchical-average]_-_8_grupos.png}
	\caption{Medóide encontrado para as curvas...}
	%\label{fig:6938medoid}
\end{figure}

\begin{figure}[h!]
	\includegraphics[width=\linewidth]{figuras/irish/2016-07-19_15:33:04.743952__irish_results/W_weekend_-_[DTW,hierarchical-average]_-_4_grupos.png}
	\caption{Medóide encontrado para as curvas...}
	%\label{fig:6938medoid}
\end{figure}

\section{Análise e discussão dos resultados}



\chapter{Estudo de caso} \label{cap:estudo_de_caso}

Neste capítulo, é descrita a metodologia para o agrupamento das curvas de carga obtidas pela medição periódica do consumo de energia elétrica das residências australianas. O objetivo deste agrupamento é verificar as influências que existem na potência consumida em cada carga em função dos dias de semana e da estação do ano. 

Em geral, tarefas de mineração de dados, independente do domínio do problema, buscam  extrair conhecimento inovador, para auxiliar as tomadas de decisões a nível gerencial ~\parencite{carvalhointeligencia}. Assim, o agrupamento dos perfis de consumo de carga buscam gerar perfis característicos de consumo em diferentes estações do ano, e em diferentes dias da semana. Ademais, espera-se realizar a segmentação dos consumidores para a proposição de uma eventual taxa diferenciada de cobrança para cada grupo de consumidores encontrado. Segundo \parencite{Flath2012}, tal faixa diferenciada criada pelas concessionárias de energia a partir da segmentação de seus consumidores residenciais tem o poder de influenciar o consumo destes, podendo, assim, diminuir os picos de consumo.

\section{Metodologia}

O conjunto de dados australiano consiste das medições de $13735$ cargas, realizadas de 6/10/2011 a 4/3/2014, no entanto, na avaliação foram utilizadas somente as medições das primeiras 5000 cargas durante um ano compreendido de 21/12/2012 a 21/12/2013. A redução do número de cargas a serem analisadas foi necessária por restrições de memória RAM, já que esta se encontrava limitada em 8GB. Já a restrição e escolha das datas de início e fim do período em análise, se deu pela necessidade de se avaliar as curvas diárias dos medidores em função da variação das estações do ano, assim, o intervalo de análise das medições foi reduzido para um ano, e teve como dia de início e fim o dia do solstício de verão do hemisfério sul, de forma que existe um número igual de medições para cada uma das quatro estações do ano.

A potência consumida foi medida pelos medidores australianos a cada $30$ minutos, totalizando $48$ medições diárias. Existem dados incompletos (medições faltantes) ou inconsistentes (medições repetidas em um mesmo instante), o que fez com que fosse necessária uma etapa de filtragem dos dados. Nesta etapa, se um medidor apresentasse dados de medição incompletos ou inconsistentes, ele seria descartado. Isso fez com que o número de medidores utilizados no agrupamento fosse reduzido drasticamente para $427$ cargas.

A tarefa de agrupamento foi realizada em seis conjuntos de dados distintos, mas cada conjunto contendo os mesmo medidores. A formação dos conjuntos de dados se deu a partir da combinação das estações do ano; inverno, verão e transição, sendo esta última compreendida pelas estações outono e inverno, com os dois tipos de dia de semana; final de semana e dia de semana. Cada um desses subconjuntos é composto pela curva diária mais representativa de cada medidor nos dias correspondentes ao grupo, portanto, cada subconjunto contém, ao final, $427$ curvas. Tal divisão é a mesma realizada em ~\parencite{Flath2012}, onde curvas de carga de consumidores residenciais, comerciais, industriais e rurais da Alemanha foram agrupados.

A formação destes $6$ subconjuntos pode ser visualizada na Figura ~\ref{fig:diagrama_dados}. A primeira caixa (Perfis diários anuais) contém $155.855$ curvas ($427*365$), e em seguida, cada uma das curvas é direcionada para o segundo nível. As curvas obtidas por medições em dias do final de semana (sábado e domingo) são direcionadas para a caixa à esquerda (Final de semana), enquanto as curvas obtidas por medições em dias de semana (segunda-feira, terça-feira, quarta-feira, quinta-feira e sexta-feira) são direcionadas para a caixa à direita (Dia de semana). Finalmente, em cada uma das caixas, cada curva é novamente direcionada para um novo subconjunto em função da estação do ano que ela foi obtida. Cada um dos subconjuntos do terceiro nível da figura contém todas as medições, no período definido pelo subconjunto, para cada um dos $427$ medidores, e, dessa maneira, para a realização do agrupamento, faz-se necessária a obtenção de uma curva representativa para cada um dos $427$ medidores em cada um dos $6$ subconjuntos.

 \begin{figure}[!h]
	%\vspace{2.5cm}
	%\resizebox{\columnwidth}{!}{%
	\includegraphics[width=\linewidth]{figuras/diagrama_dados.png}
	%}
	\caption{Formação dos $6$ subconjuntos de dados a partir dos perfis diários.}
	\label{fig:diagrama_dados}
\end{figure}

A curva diária mais representativa de cada medidor foi escolhida como o centróide de todas as medições naquele subconjunto, pois esta escolha, em detrimento do medóide, fez com que as curvas características de cada medidor se tornassem mais suaves, o que facilitou a tarefa de agrupamento e de interpretação dos resultados. Após a obtenção do centróide de cada medidor, em cada um dos seis subconjuntos, obteve-se seis matrizes de $427x48$, já que em um dia foram realizadas 48 medições por cada um dos $427$ medidores em cada um dos seis subconjunto de dados. Uma vez que todos os centróides possuíam valores positivos e a normalização max é a mais utilizada no agrupamento de curvas de carga ~\parencite{Chicco}, então neste estudo de caso, cada centróide foi normalizado pela normalização max descrita na Seção ~\ref{sec:norm_Z}.

\section{Agrupamento}

Cada um dos seis conjuntos de dados foi agrupado segundo a estratégia definida no Capítulo ~\ref{cap:testes_teoricos}, ou seja, algoritmo hierárquico com \emph{linkagem average} e métrica de dissimilaridade, e avaliação do índice Calinski-Harabasz para a indicação do número de grupos, que foi variado de $2$ a $20$. Notou-se que os melhores valores desses índices eram obtidos em partições altamente desbalanceadas, ou seja, partições que possuíam um dos grupos contendo cerca de $95\%$ dos medidores. Uma partição com tamanho desbalanceamento não é, na maior parte das vezes, adequada, pois não se pode extrair nenhuma informação útil dela, já que não se pode realizar a segmentação dos consumidores de energia.

A partir dos resultados insatisfatórios obtidos a partir da estratégia indicada no final do ~\ref{cap:testes_teoricos}, foram realizados agrupamentos com os seguintes algoritmos: k-means, k-medoids, hierarchical-average . Para cada um dos algoritmos, com exceção do k-means que foi realizado somente para a distância euclidiana, foi realizado um agrupamento com as seguintes dissimilaridades: euclidiana, cid-euclidiana, cort-euclidiana, chebyshev, cid-chebyshev, cort-chebyshev, DTW, cid-DTW, cort-DTW, EDR, cid-EDR, cort-EDR. Para cada combinação possível de algoritmo e dissimilaridade, o valor $k$ do número de grupos foi, novamente, variado de $2$ a $20$. A partir das combinações de algoritmo, métrica de dissimilaridade e $k$, formou-se $684$ partições diferentes para cada um dos  seis conjuntos de dados.

Cada partição foi avaliada pelos índices internos silhouette e Calinski-Harabasz, onde, novamente, todas as soluções apontadas pelos índices se encontravam altamente desbalanceadas.  Em uma tentativa de se obter partições mais balanceadas, descartou-se as instâncias que compunham os grupos minoritários, por pensar que estas poderiam ser \emph{outliers}, e o agrupamento foi realizado novamente somente no grupo majoritário, que continha quase a totalidade das instâncias. Os resultados acarretaram, novamente, em partições altamente desbalanceadas. Este procedimento foi repetido reiteradamente, até que se encontra-se um número muito pequeno de cargas, cerca de 10\% do número inicial.

Em uma terceira tentativa, para realizar o agrupamento dos subconjuntos, utilizou-se o DBSCAN, devido, principalmente, à sua robustez frente a conjuntos de dados ruidosos. Foi utilizada a métrica CID-Euclidiana, e os critérios de \emph{MinPts} e \emph{Eps}, foram escolhidos pelo gráfico \emph{k-dist}, conforme exemplo da seção ~\ref{sec:dbscan}. Novamente, o DBSCAN apontou uma partição dos dados composta por um único grupo majoritário e um número reduzido de instâncias que ele classificou como \emph{outliers}. Assim, como na abordagem com o algoritmo hierárquico, as instâncias que foram consideradas \emph{outliers} foram descartadas e o DBSCAN foi novamente empregado nas instâncias que compunham o grupo majoritário. Uma única partição majoritária foi novamente obtida, com um número reduzido de \emph{outliers}. Assim, chegou-se a conclusão que o grupo de dados australiano é composto, efetivamente, por instâncias compactas, de forma que o agrupamento das mesmas deverá ser realizado por outra técnica não abordada até o momento.

No trabalho realizado em ~\parencite{Flath2012}, que foi inspiração para este estudo de caso, não foram obtidas partições tão desbalanceadas  com um procedimento similar ao empregado na base de dados australiana, salvo à utilização do índice DBI ao invés do índice Calnski-Harabasz para a indicação do número de grupos. Dentre as 215 cargas alemãs utilizadas no referido estudo, existiam cargas industriais, residenciais, rurais e comerciais, diferentemente da base de dados australiana que contém somente dados de medidores inteligentes instalados em cargas residenciais. Dessa maneira, dada a própria heterogeneidade dos tipos de cargas a serem agrupadas, no estudo alemão, foi possível a obtenção de grupos mais balanceados com as técnicas empregadas até o momento.

\subsection{Estratégia alternativa de agrupamento}

Após os resultados insatisfatórios obtidos nas partições, deu-se prosseguimento à uma nova estratégia para a escolha da partição de cada um dos seis conjuntos de dados construídos. Calculou-se as porcentagens de cargas  contidas em cada grupo, de cada partição obtida. A partir dos valores de porcentagem de cada grupo, obteve-se os valores mínimos, máximos, médios, e o desvio padrão das porcentagens de cada partição. Partições cujo valor máximo excedeu $50\%$, ou cujo valor mínimo foi inferior à $5\%$ foram descartadas. O outro critério a ser utilizado na escolha da melhor partição de cada grupo foi o índice silhouette. Este índice foi escolhido por ter um limite superior, ou seja, valor máximo igual a $1$, o que faz com que seja possível a sua utilização na comparação de duas partições obtidas por diferentes dissimilaridades, por exemplo. Ao se utilizar o índice Calinski-Harabasz, valores muito discrepantes são obtidos para as diferentes partições, e por não possuir um valor máximo, os valores de cada partição não podem ser normalizados para que a sua comparação seja realizada efetivamente.

Observou-se que as partições que possuíam valores maiores do índice silhouette eram as mais desbalanceadas, ou seja, com um alto valor de desvio padrão, ao passo que as partições com um baixo desvio padrão, ou balanceadas, possuíam um valor de silhouette aproximadamente igual a $0$, o que indica partições cujo agrupamento foi realizado de forma próxima à aleatória. Assim, dada à característica conflitante observada entre esses dois objetivos; qualidade da partição, associada ao valor de silhouette, e balanceamento da mesma, associada ao desvio padrão das porcentagens de cargas em cada grupo, ficou caracterizado um problema de otimização multiobjetivo. Logo, fez-se necessário encontrar as soluções, ou partições, não dominadas e da escolha de uma solução, ou partição, dentre aquelas soluções que formam a fronteira Pareto-ótima.

Não houve a necessidade de se utilizar um método de otimização multiobjetivo, pois todas as  partições já estavam calculadas. O que se fez foi encontrar as soluções não dominadas que formam a fronteira Pareto-ótima. Uma vez encontrada as soluções que formam a fronteira de Pareto-ótima, calculou-se a distância da solução ótima $[1,0]$, ou seja valor de silhouette igual a $1$ e desvio padrão igual à $0$, a cada uma das soluções da fronteira Pareto-ótima. Aquela que possuía a menor distância foi escolhida como partição adequada para aquele conjunto de dados. Os resultados, ou seja, as partições escolhidas, podem ser encontrados na tabela ~\ref{tbl:resultado_australia}.

O espaço de soluções bem como as curvas características de cada curva em cada um dos seis conjuntos de dados podem ser visualizados nas figuras ~\ref{fig:pareto_FDS_verao}, ~\ref{fig:pareto_FDS_inverno}, ~\ref{fig:pareto_FDS_transicao}, ~\ref{fig:pareto_DDS_verao}, ~\ref{fig:pareto_DDS_inverno} e ~\ref{fig:pareto_DDS_transicao}. Nestas figuras, cada ponto representa uma partição, e pode-se notar que existem cerca de $20$ pontos em cada figura, ao invés das $684$ partições inicialmente obtidas. Isso ocorreu devido à restrição imposta ao se descartar partições cujo grupo majoritário contêm mais de \%50 das instâncias ou o grupo minoritário contém menos de \%5 do total de instâncias. Esta redução drástica no número de partições, reforça o caráter compacto do grupo, já que para diferentes escolhas de algorimo e métrica de dissimilaridade, a maioria das partições obtidas não atendeu à restrição imposta. Os pontos em vermelho nas figuras compõem a fronteira Pareto-ótima, e a solução, ou partição, indicada é representada pelo ponto amarelo.

%\begin{center}
	%\begin{table}			
		%\caption{Tabela que apresenta as informações das partições escolhidas para cada um dos seis sub-conjuntos gerados.} 
		%\resizebox{\columnwidth}{!}{%
		%\begin{tabular}{lllllllrrrrrr}
			%\toprule
			%{} & Norm. &    method & Dia de semana & Estação &        Alg. &              diss. &   k &       max &      mean &       min &  silhouette &       std \\
			%\midrule
			%292  &  max &  centroid &      False &      T &    k-means &      minkowski\_2 &   9 &  0.175644 &  0.111111 &  0.067916 &    0.124528 &  0.032075 \\
			%977  &  max &  centroid &      False &      W &    k-means &      minkowski\_2 &  10 &  0.142857 &  0.100000 &  0.060890 &    0.113796 &  0.026455 \\
			%1659 &  max &  centroid &      False &      S &    k-means &      minkowski\_2 &   8 &  0.189696 &  0.125000 &  0.065574 &    0.115326 &  0.034893 \\
			%2262 &  max &  centroid &       True &      T &  k-medoids &  cid-minkowski\_2 &   3 &  0.360656 &  0.333333 &  0.306792 &    0.066607 &  0.021997 %\\
		%	3028 &  max &  centroid &       True &      W &    k-means &      minkowski\_2 &   9 &  0.152225 &  0.111111 &  0.060890 &    0.126416 &  0.026659 \\
		%	3706 &  max &  centroid &       True &      S &    k-means &      minkowski\_2 &   3 &  0.355972 &  0.333333 &  0.295082 &    0.165268 &  0.027199 \\
		%	\bottomrule
	%	\end{tabular}

%	}
%	\end{table} \label{tbl:resultado_australia}
%\end{center}

\begin{center}
	\begin{table}			
		\caption{Tabela que apresenta as informações das partições escolhidas para cada um dos seis sub-conjuntos gerados.} 
		\resizebox{\columnwidth}{!}{%
			\begin{tabular}{lllllllrrrrrr}
				\toprule
				Norm. &    Dia de semana & Estação &        Alg. &              diss. &   k &       max \%&      méd. \% &       min \% &  silhouette &       std \% \\
				\midrule
				  max &        {} &      Transição &    k-means &      eucl. &   9 &  0.17 &  0.11 &  0.07 &    0.12 &  0.03 \\
				  max &  		{} &      Inverno &    k-means &      eucl. &  10 &  0.14 &  0.10 &  0.060 &    0.11 &  0.02 \\
				  max &  {} &      Verão &    k-means &      eucl. &   8 &  0.18 &  0.12 &  0.06 &    0.11 &  0.03 \\
				  max &  \checkmark &      Transição &  k-medoids &  cid-eucl. &   3 &  0.36 &  0.33 &  0.31 &    0.07 &  0.02 \\
				  max &  \checkmark &      Inverno &    k-means &      eucl. &   9 &  0.15 &  0.11 &  0.06 &    0.13 &  0.02 \\
				  max &  \checkmark &      Verão &    k-means &      eucl. &   3 &  0.36 &  0.33 &  0.30 &    0.17&  0.03 \\
				\bottomrule
			\end{tabular}
			
		}
	\label{tbl:resultado_australia}
	\end{table} 
\end{center}

\begin{figure}[!h]
	\centering
	\begin{subfigure}{.5\textwidth}
		\centering
		\includegraphics[width=.9\linewidth]{figuras/australia_5000/pareto_Final_de_Semana_Verao.png}
		\caption{Soluções no espaço de soluções.}
		\label{fig:pareto_FDS_verao}
	\end{subfigure}%
	\begin{subfigure}{.5\textwidth}
		\centering
		\includegraphics[width=.9\linewidth]{figuras/australia_5000/Final_de_Semana_Verao.png}
		\caption{Medóides de cada um dos grupos da partição escolhida.}
		\label{fig:FDS_verao}
	\end{subfigure}
	\caption{Resultado final para os curvas diárias das cargas medidas nos dias do final de semana e no verão. Os pontos em azul, na Figura ~\ref{fig:pareto_FDS_verao}, representam partições dominadas pelas demais, enquanto os pontos vermelhos formam a fronteira de Pareto-ótima. O ponto em amarelo é a partição escolhida por ser a mais próxima do ponto ótimo.}
	\label{fig:FDS_verao_}
\end{figure}
\begin{figure}[!h]
	\centering
	\begin{subfigure}{.5\textwidth}
		\centering
		\includegraphics[width=.9\linewidth]{figuras/australia_5000/pareto_Final_de_Semana_Inverno.png}
		\caption{Soluções no espaço de soluções.}
		\label{fig:pareto_FDS_inverno}
	\end{subfigure}%
	\begin{subfigure}{.5\textwidth}
		\centering
		\includegraphics[width=.9\linewidth]{figuras/australia_5000/Final_de_Semana_Inverno.png}
		\caption{Medóides de cada um dos grupos da partição escolhida.}
		\label{fig:FDS_inverno}
	\end{subfigure}
	\caption{Resultado final para os curvas diárias das cargas medidas nos dias do final de semana e no inverno. Os pontos em azul, na Figura ~\ref{fig:pareto_FDS_inverno}, representam partições dominadas pelas demais, enquanto os pontos vermelhos formam a fronteira de Pareto-ótima. O ponto em amarelo é a partição escolhida por ser a mais próxima do ponto ótimo.}
	\label{fig:FDS_inverno_}
\end{figure}

\begin{figure}[!h]
	\centering
	\begin{subfigure}{.5\textwidth}
		\centering
		\includegraphics[width=.9\linewidth]{figuras/australia_5000/pareto_Final_de_Semana_Transicao.png}
		\caption{Soluções no espaço de soluções.}
		\label{fig:pareto_FDS_transicao}
	\end{subfigure}%
	\begin{subfigure}{.5\textwidth}
		\centering
		\includegraphics[width=.9\linewidth]{figuras/australia_5000/Final_de_Semana_Transicao.png}
		\caption{Medóides de cada um dos grupos da partição escolhida.}
		\label{fig:FDS_transicao}
	\end{subfigure}
	\caption{Resultado final para os curvas diárias das cargas medidas nos dias do final de semana e nas estações de transição. Os pontos em azul, na Figura ~\ref{fig:pareto_FDS_transicao}, representam partições dominadas pelas demais, enquanto os pontos vermelhos formam a fronteira de Pareto-ótima. O ponto em amarelo é a partição escolhida por ser a mais próxima do ponto ótimo}
	\label{fig:FDS_transicao_}
\end{figure}

\begin{figure}[!h]
	\centering
	\begin{subfigure}{.5\textwidth}
		\centering
		\includegraphics[width=.9\linewidth]{figuras/australia_5000/pareto_Dia_de_semana_Verao.png}
		\caption{Soluções no espaço de soluções.}
		\label{fig:pareto_DDS_verao}
	\end{subfigure}%
	\begin{subfigure}{.5\textwidth}
		\centering
		\includegraphics[width=.9\linewidth]{figuras/australia_5000/Dia_de_semana_Verao.png}
		\caption{Medóides de cada um dos grupos da partição escolhida.}
		\label{fig:DDS_verao}
	\end{subfigure}
	\caption{Resultado final para os curvas diárias das cargas medidas nos dias de semana e no verão. Os pontos em azul, na Figura ~\ref{fig:pareto_DDS_verao}, representam partições dominadas pelas demais, enquanto os pontos vermelhos formam a fronteira de Pareto-ótima. O ponto em amarelo é a partição escolhida por ser a mais próxima do ponto ótimo}
	\label{fig:DDS_verao_}
\end{figure}

\begin{figure}[!h]
	\centering
	\begin{subfigure}{.5\textwidth}
		\centering
		\includegraphics[width=.9\linewidth]{figuras/australia_5000/pareto_Dia_de_semana_Inverno.png}
		\caption{Soluções no espaço de soluções.}
		\label{fig:pareto_DDS_inverno}
	\end{subfigure}%
	\begin{subfigure}{.5\textwidth}
		\centering
		\includegraphics[width=.9\linewidth]{figuras/australia_5000/Dia_de_semana_Inverno.png}
		\caption{Medóides de cada um dos grupos da partição escolhida.}
		\label{fig:DDS_inverno}
	\end{subfigure}
	\caption{Resultado final para os curvas diárias das cargas medidas nos dias de semana e no inverno. Os pontos em azul, na Figura ~\ref{fig:pareto_DDS_inverno}, representam partições dominadas pelas demais, enquanto os pontos vermelhos formam a fronteira de Pareto-ótima. O ponto em amarelo é a partição escolhida por ser a mais próxima do ponto ótimo}
	\label{fig:DDS_inverno_}
\end{figure}

\begin{figure}[!h]
	\centering
	\begin{subfigure}{.5\textwidth}
		\centering
		\includegraphics[width=.9\linewidth]{figuras/australia_5000/pareto_Dia_de_semana_Transicao.png}
		\caption{Soluções no espaço de soluções.}
		\label{fig:pareto_DDS_transicao}
	\end{subfigure}%
	\begin{subfigure}{.5\textwidth}
		\centering
		\includegraphics[width=.9\linewidth]{figuras/australia_5000/Dia_de_semana_Transicao.png}
		\caption{Medóides de cada um dos grupos da partição escolhida.}
		\label{fig:DDS_transicao}
	\end{subfigure}
	\caption{Resultado final para os curvas diárias das cargas medidas nos dias de semana e na transição. Os pontos em azul, na Figura ~\ref{fig:pareto_DDS_transicao}, representam partições dominadas pelas demais, enquanto os pontos vermelhos formam a fronteira de Pareto-ótima. O ponto em amarelo é a partição escolhida por ser a mais próxima do ponto ótimo}
	\label{fig:DDS_transicao_}
\end{figure}

\section{Análise dos resultados}

A abordagem tradicional do agrupamento, com a combinação mais promissora apontada no Capítulo ~\ref{cap:testes_teoricos}, resultou em partições altamente desbalanceadas, o que não é adequado do ponto de vista prático para a descoberta de perfis e segmentação de consumidores residenciais. Foi então proposta uma metodologia que consiga encontrar partições dos dados que consiga estabelecer um compromisso entre o balanceamento da partição, e a qualidade da mesma, ou seja, partições com alta densidade intergrupo e sem a presença de um grupo que contém quase a totalidade das instâncias.

Diferentemente dos resultados obtidos no capítulo \ref{cap:testes_teoricos}, onde se tinha disponível a rotulação dos dados, e os melhores resultados foram obtidos com os algoritmos hierárquicos e métrica de dissimilaridade CID-DTW, neste estudo de caso, o conjunto de dados não possui rotulação, as partições obtidas por meio do algoritmo k-means aliado à métrica euclidiana foram os que alcançaram os resultados mais satisfatórios. As demais combinações testadas de algoritmo e métrica de dissimilaridade acarretaram em partições altamente desbalanceadas. Das $684$ partições obtidas ao se variar o algoritmo de agrupamento, a métrica de dissimilaridade e o número de grupos, após a seleção de partições cujo maior grupo representava menos de 50\% dos medidores e o menor mais do que 5\%, o número de soluções disponíveis foi reduzido drasticamente. Isso pode ser percebido nos gráficos do espaço de soluções ~\ref{fig:pareto_FDS_verao}, ~\ref{fig:pareto_FDS_inverno}, ~\ref{fig:pareto_FDS_transicao}, ~\ref{fig:pareto_DDS_verao}, ~\ref{fig:pareto_DDS_inverno} e ~\ref{fig:pareto_DDS_transicao},  já que, em cada um deles, existem aproximadamente $20$ soluções possíveis das $684$ iniciais.

Assim, pode-se inferir que o algoritmo \emph{k-means} é o mais adequado, no contexto de agrupamento de curvas de carga, quando se deseja obter partições balanceadas. Outro resultado que parece interessante é o fato do número de grupos encontrado ser maior nos conjuntos de dados obtidos a partir do final de semana. Nestes, o valor de $k$ foi em torno de $9$, ao passo que nos conjuntos de dados contendo medições realizadas nos dias de semana, o número $k$ foi igual a 3. Isso leva a crer que durante o final de semana a diversidade de perfis de consumo de carga aumenta em relação aos dias de semana. Além disso, este estudo indica que este fator, o tipo de dia de semana, tem maior impacto no número de grupos $k$ do que a variação das estações do ano. Isto pode ser dito no contexto australiano, pois, caso este estudo fizesse uso de dados medidos em outros países, como no Brasil, por exemplo, os resultados poderiam ter sido bem diferentes. 

Em relação aos perfis de consumo em cada subconjunto, pode-se perceber, pelas Figuras  ~\ref{fig:pareto_FDS_verao}, ~\ref{fig:pareto_FDS_inverno}, ~\ref{fig:pareto_FDS_transicao}, ~\ref{fig:pareto_DDS_verao}, ~\ref{fig:pareto_DDS_inverno} e ~\ref{fig:pareto_DDS_transicao}, que:
\begin{itemize}
	\item na maioria dos perfis característicos, as curvas têm um crescimento abrupto aproximadamente na décima segunda medição (às 6 horas da manhã),
	\item na maioria dos perfis característicos, as curvas têm um decrescimento abrupto aproximadamente na quadragésima medição (às 20 horas da noite),
	\item na maioria dos perfis característicos dos finais de semana, as curvas apresentam um vale centrado na vigésima quarta medição (às 12 horas),
	\item os perfis de dia de semana no verão (vide Figura ~\ref{fig:DDS_verao}) e nos meses de transição (vide Figura ~\ref{fig:DDS_transicao}) são bastante similares, o que reforça a suposição de que nas residências australianas a variação da estação do ano não influencia no comportamento de consumo de energia elétrica dos australianos, mas o tipo de dia de semana sim, além de ambos apresentarem uma tendência de crescimento durante o dia,
	\item no inverno e dia de semana (vide Figura ~\ref{fig:DDS_inverno}), existem alguns perfis que divergem significativamente dos demais, apresentando picos significativos de consumo,
\end{itemize}
todas essas observações são, de certa forma esperadas para curvas de carga de residências, o que mostra que os resultados são coerentes. No entanto, a interpretação das curvas aliado à outros dados dos consumidores que compõem cada grupo dos subconjuntos em análise, tais como a renda, localização, número de moradores, dentre outros, poderiam gerar informações mais ricas e úteis para as concessionárias de energia.

%Como estas outras informações não se encontram disponíveis, pode-se realizar algumas especulações, acercas dos resultados:
%\begin{itemize}
%	\item 
%\end{itemize}

Há de se destacar que a estratégia implementada neste estudo de caso foi realizada em uma base de dados reduzida, com 427 cargas, e dessa maneira foi possível realizar as $684$ partições diferentes de cada um dos seis conjuntos de dados com um esforço computacional moderado. O agrupamento dos dados gerados por medidores inteligentes instalados por uma concessionária de energia elétrica como a CEMIG, que possui milhões de consumidores residenciais ~\parencite{DadosCemig}, demandará um custo computacional relevante. Assim, em um cenário real, após a massificação da instalação de medidores inteligentes no sistema elétrico de potência brasileiro, é mister que seja definida uma única abordagem, em relação à escolha do algoritmo e métrica de dissimilaridade.

Assim, pelos resultados obtidos, sugere-se que o agrupamento de curvas de cargas residenciais seja realizado, após a normalização max, por meio do algoritmo \emph{k-means} e métrica de dissimilaridade euclidiana. Esse resultado vai de encontro com a literatura ~\parencite{Chicco}. Sugere-se ainda que caso seja feita a mesma divisão dos dados realizados neste estudo de caso, ou seja, em tipos de dia da semana e em estações do ano, que o parâmetro $k$ do algoritmo \emph{k-means} seja aproximadamente igual à 3 para os conjuntos de dados contendo medições nos dias de semana, e aproximadamente $10$ para os dias dos finais de semana. 

A indicação algoritmo k-means para a realização do agrupamento de curvas de carga, por gerar partições que possuem um bom compromisso entre balanceamento dos grupos e qualidade da partição, possui ainda a vantagem deste ser um algoritmo altamente escalável, ou em outras palavras,  altamente paralelizável ~\parencite{zhao2009parallel, stoffel1999parallel, bahmani2012scalable}, tanto localmente quanto distribuidamente, o que, por sua vez, também é adequado para uma massa de dados de milhões de consumidores.

\chapter{Conclusões e propostas de trabalhos futuros} \label{cap:conclusao}

Este trabalho teve como objetivo a realização de um estudo das estratégias de agrupamento de séries temporais para a posterior aplicação das técnicas estudadas no agrupamento de dados de curvas de consumo de cargas residenciais. Dessa maneira, inicialmente, as possíveis escolhas ao se realizar o agrupamento de séries temporais foram analisadas.

As técnicas de pré-processamento foram discutidas na qual foram expostas algumas técnicas de redução de dimensionalidade como a PAA ou TDF, e se verificou que algum tipo de normalização das curvas deve ser realizada para que a comparação entre as curvas faça sentido. Em seguida, foram expostas as medidas de dissimilaridade entre as curvas, no qual obtiveram destaque algumas medidas como a DTW e EDR, além da correção das dissimilaridades pela CID \parencite{CID}. Logo após, foram apresentados os principais algoritmos encontrados na literatura como os algoritmos \emph{k-means} e hierárquicos. Finalmente, os índices de validação de agrupamentos foram discutidos, dando-se certo destaque para o índice de validação externo de Variação da Informação (VI), pelo fato deste ser uma métrica no espaço de partições.

A partir da revisão bibliográfica realizada no Capítulo ~\ref{cap:agrupamento_series_temporais}, no Capítulo ~\ref{cap:testes_teoricos} realizou-se um estudo empírico de agrupamento de bases de dados rotuladas, onde os resultados foram comparados com o resultado esperado pelo índice VI. Nos testes foram realizados agrupamentos, em cada base de dados, variando-se os principais algoritmos e medidas de dissimilaridade mencionados na literatura. Testes estatísticos de significância apontaram que os algoritmos hierárquicos, principalmente ao se utilizar a \emph{linkagem} \emph{average}, possuem performance melhor que os algoritmos particionais. 

Realizando-se os agrupamentos com o algoritmo hierárquico e \emph{linkagem average} e diferentes métricas de dissimilaridade, verificou-se que as dissimilaridades DTW \parencite{DTW} possuem desempenho superior. No entanto, ao se aplicar a CID à dissimilaridades como a euclidiana, estas apresentaram desempenho estatisticamente equivalente ao desempenho da DTW quando também corrigida pela CID. Este resultado é interessante já que existe grande diferença de custos computacionais para se obter a DTW e distância euclidiana entre duas séries temporais, e que este custo para a última é significativamente menor.

Em seguida, foram realizados testes para analisar a performance dos índices de validação interno. Em um primeiro experimento, os índices internos foram testados para indicar o número de grupos  existentes em uma partição e, neste experimento, após as análises de significância, chegou-se a conclusão que o índice Calinski-Harabazs \parencite{CH} é o mais adequado para a realização de tal tarefa no contexto de séries temporais. Uma vez definido o número de grupos, surge a pergunta de qual índice de validação interno é o mais adequado para avaliar as partições obtidas por diferentes abordagens no que diz respeito às escolhas da técnica de pré-processamento, algoritmo de agrupamento e métrica de dissimilaridade. Para responder à esta pergunta, novos testes foram realizados nas bases de dados rotuladas, e os resultados mostraram que os índices silhouette \parencite{silhouette} e Dunn ~\parencite{Dunn} possuem performance significativamente superior aos outros quatro índices testados.

Dessa maneira, partindo-se do pressuposto que não se possui, \emph{a priori}, nenhuma informação útil para a realização do agrupamento, como, por exemplo, o número de grupos que as instâncias se encontram naturalmente divididas, os experimentos realizados em bases de dados rotuladas sugerem a seguinte abordagem a ser realizada no agrupamento de séries temporais:

\begin{itemize}
	\item Realizar algum tipo de normalização das curvas. Sugere-se a normalização Z ou normalização max.
	\item Escolher alguma medida de dissimilaridade. Sugere-se a CID-DTW caso o conjunto de dados não seja muito grande, ou os recursos computacionais escassos. Em qualquer uma dessas duas hipóteses, sugere-se a utilização da CID-Euclidiana.
	\item Para se definir o número de grupos $k$, realizar o agrupamento do conjunto de dados já normalizados pelo algoritmo hierárquico com \emph{linkagem average} e medida de dissimilaridade escolhida. Em seguida, obter diferentes partições a partir da poda do dendrograma obtido por valores iterativos de $k$, em um intervalo no qual se espera que se encontre o valor de $k$.
	\item Em seguida, para cada uma das partições obtidas, realizar a avaliação de cada uma delas pelo índice Calinski-Harabasz, e aquela que possuir o melhor valor indicará o valor de $k$.
	\item Caso se deseje realizar um refinamento da partição obtida até o momento, uma vez definido o número de grupos $k$, pode-se realizar outras partições com diferentes escolhas de combinações de algoritmos e métricas de dissimilaridade e avaliá-las pelo índice silhouette. A partição que acarretar no melhor valor do índice é a partição final da tarefa de agrupamento.
\end{itemize}

Após a definição dos passos descritos anteriormente, obtidos por meio de testes em bases de dados rotuladas, estes foram aplicados em uma base de dados não rotulada no Capitulo ~\ref{cap:estudo_de_caso}. O estudo de caso consistiu no agrupamento dos dados de consumo de cargas residenciais, gerados por medidores inteligentes instalados em residências australianas. As curvas diárias de consumo de cada residência foram divididas em seis subconjuntos, formados cada um pela média das curvas diárias em cada uma das seis possíveis combinações entre as estações do ano; inverno, verão e transição (outono e primavera), e os tipos de dias da semana; dias de final de semana e dias de semana.

Os resultados obtidos, em todos os subconjuntos, por meio dos passos descritos anteriormente, não foram satisfatórios, pois partições altamente desbalanceadas foram obtidas. Por partições altamente desbalanceadas entende-se partições dos dados nas quais um único grupo contém quase a totalidade dos dados, ao passo que os demais grupos, denominados grupos minoritários, contêm poucas instâncias. Partições com essas características são indesejadas, pois, na maioria das vezes, não é possível identificar padrões relevantes no agrupamento dos dados. Mesmo após a eliminação das cargas contidas nos grupos minoritários, por estas terem sido, inicialmente, consideradas \emph{outliers}, a realização de novos agrupamentos acarretou, novamente, em partições altamente desbalanceadas.

Assim, decidiu-se utilizar uma técnica de avaliação de partição, que pelo menos no estudo realizado nesta dissertação, não foi visto em nenhuma das fontes bibliográficas utilizadas e mencionadas neste trabalho. Na abordagem proposta, cada uma das partições obtidas foram avaliadas por dois critérios: pela qualidade da partição, indicada pelo índice interno silhouette, e pelo balanceamento da mesma, indicado pelo desvio padrão da porcentagem de instâncias nos grupos da partição em análise. Nas partições obtidas por meio de diversas combinações de algoritmo de agrupamento e métrica de dissimilaridade, verificou-se que estes dois critérios são conflitantes, ou seja, partições que apresentavam os melhores valores de silhouette possuíam os piores valores de desvio padrão, ao passo que as partições com os melhores valores de desvio padrão possuíam os piores valores de silhouette.

Dado o caráter conflitante entre as duas características almejadas para a partição final, ficou caracterizado um problema de otimização multi-objetivo, e assim fez-se necessário encontrar as soluções, ou partições não dominadas. A partição final escolhida, foi aquela que pertence à fronteira Pareto-ótima que se encontra mais perto do ponto ótimo, ou seja do ponto $[1,0]$, no qual a partição possui silhouette igual a $1$ e desvio padrão igual a $0$. As partições encontradas para cada um dos seis subgrupos formados apresentaram alguns padrões interessantes:

\begin{itemize}
	\item Todos os subgrupos encontrados utilizaram como algoritmo de agrupamento o \emph{k-means}, apesar de outros também terem sido testados, do que se  pode inferir que o \emph{k-means} gera partições que apresentam um bom compromisso entre a qualidade da partição e o balanceamento do número de instâncias nos grupos.
	\item Os subgrupos que continham dados de curva de carga obtidas nos finais de semana foram particionados em um número de grupos significativamente maior do que os subgrupos contendo curvas obtidas nos dias de semana, o que indica que existe uma maior heterogeneidade de perfis de consumo residenciais nos finais de semana em comparação com os dias de semana.
\end{itemize}

O fato da estratégia definida a partir dos experimentos realizados em bases de dados rotuladas no Capítulo ~\ref{cap:testes_teoricos} não ter gerado um resultado satisfatório ao ser aplicada em uma base de dado não rotulada no Capítulo ~\ref{cap:estudo_de_caso} não a desmerece, já que ela deve ser vista como uma sugestão inicial para problemas de agrupamento de séries temporais, e não como uma solução definitiva. Vale destacar que algumas decisões da tarefa de agrupamento de séries temporais, como a escolha da métrica de dissimilaridade, dependem, muitas vezes, das características intrínsecas do problema em questão. No estudo realizado com as bases rotuladas, foram utilizadas séries temporais de diversas aplicações e áreas do conhecimento como sinais biomédicos, problemas de visão computacional, detecção de movimento, sinais de equipamentos elétricos, entre outros. Assim, reforça-se a sua sugestão como uma primeira alternativa para a realização do agrupamento de séries temporais, e caso resultados insatisfatórios sejam obtidos com ela, outras abordagens podem ser adotadas.

No contexto de curvas de cargas obtidas por medidores inteligentes, o estudo de caso realizado no Capítulo ~\ref{cap:estudo_de_caso} indica que o algoritmo \emph{k-means}, aliado à métrica de dissimilaridade euclidiana, fornece resultados satisfatórios. Tais escolhas estão em consonância com a literatura ~\parencite{Chicco} e são interessantes para a realização de agrupamento de grandes massas de dados, como as que serão geradas após a massificação da instalação de medidores inteligentes no sistema elétrico de potência brasileiro, pois a distância euclidiana tem custo computacional significativamente menor que suas principais concorrentes (DTW e EDR), além do algoritmo \emph{k-means} ser altamente escalável \parencite{zhao2009parallel, stoffel1999parallel, bahmani2012scalable}. 

Assim, pode-se considerar que este contribui na indicação de metodologias (escolhas), no que diz respeito a pré-processamento, algoritmo de agrupamento e métrica de dissimilaridade, dentre as várias existentes na literatura, a serem utilizadas no agrupamento de curvas de carga obtidas por medidores inteligentes.

\section{Propostas de trabalhos futuros}

A dissertação de mestrado apresentada é somente uma fração do estudo de agrupamento de séries temporais e do seu emprego em agrupamentos de curvas de cargas, assim, sugere-se o aprofundamento nos seguintes temas:

\begin{itemize}
	\item Avaliar, no contexto de agrupamento, a performance da medida de dissimilaridade TWED (\emph{Time Warp Edit Distance }) proposta em \parencite{10.1109/TPAMI.2008.76}, que em experimentos realizados em problemas de classificação, expostos em \parencite{Serra},  apresentou performance superior à distância euclidiana, DTW, EDR e suas respectivas correções pela CID.
	\item Estudar a variação no cálculo da DTW descrita em \parencite{Petitjean:2011:GAM:1890011.1890193} que faz com que seja possível a utilização dela no algoritmo \emph{k-means}.
	\item Verificar os métodos descritos em \parencite{Morse:2007:EAM:1247480.1247544} que reduzem os custos computacionais envolvidos no cálculo da DTW e EDR.
	\item Estudar as técnicas de agrupamento consensuais, que combinam diversas partições para a obtenção de uma partição final, no contexto de séries temporais ~\parencite{consenso}.
	\item Investigar a robustez da proposta de agrupamento multi-objetiva realizada nesta dissertação para o agrupamento de conjuntos de dados desbalanceados.
\end{itemize}
% Referência bibliográfica.
\newpage
\phantomsection
\addcontentsline{toc}{chapter}{Referências}


\printbibliography[title=Referencias]

%
% FIXME Se não for utilizar apêndices (o glossario é um apêndice), comentar a
% linha abaixo.
%\appendix
% FIXME Remover as 4 linha abaixo.
%\chapter*{Glossário\markboth{Glossário}{}}  % \markboth{}{} é utilizado para corrigir o cabeçalho.
\addcontentsline{toc}{chapter}{Glossário}

\begin{description}
  % FIXME Remover as duas explicações abaixo e incluir as que serão
  % utilizadas.
    \item[REI] Redes Elétricas Inteligentes
    \item[MI] Medidores Inteligentes
    \item[VI] \emph{Variaotion of information}, índice interno de validação de agrupamentos
    \item[RAM] \emph{Random Acess Memory}, ou memória de acesso aleatório
    \item[XB] Xie-Beni \emph{index}, índice de validação externa de Xie-Beni
    
\end{description}

\input{ex_ape1}
%\input{ex_ape2}
%\input{ex_ape3}
%\input{ex_ape4}
% TODO Inserir os arquivos referentes aos apêndices.
%
% TODO Se não for utilizar anexos nem uma licença Creative Commons, comentar a linha abaixo.
%\annex
% FIXME Remover a linha abaixo.
%\input{ex_ane1}
% TODO Inserir os arquivos referentes aos anexos.
%
% FIXME Escolher, opcionalmente, uma das licença Creative Commons abaixo.
% Para isso remova a linha correspondente a licença não desejada.
%\input{cc-by}  % Creative Commons Atribuição 3.0 Não Adaptada
%q\input{cc-by-sa}  % Creative Commons AtribuiçãoCompartilhaIgual 3.0 Não Adaptada
%
\clearpage
\backmatter
\phantomsection
\addcontentsline{toc}{chapter}{Índice Remissivo}
\printindex
\end{document}
